% -------------------------------------------------------------
% NOTE ON THE DETAILED AND SHORT VERSIONS:
% -------------------------------------------------------------
% This paper comes in two versions, a detailed and a short one.
% The short version should be more than sufficient for any
% reasonable use; the detailed one was written purely to
% convince the author of its correctness.
% To switch between the two versions, find the line containing
% "\newenvironment{noncompile}{}{}" in this LaTeX file.
% Look at the two lines right beneath this line.
% To compile the detailed version, they should be as follows:
%   \includecomment{verlong}
%   \excludecomment{vershort}
% To compile the short version, they should be as follows:
%   \excludecomment{verlong}
%   \includecomment{vershort}
% As a rule, the line
%   \excludecomment{noncompile}
% should stay as it is.
% -------------------------------------------------------------
% NOTES ON SOME HACKS USED IN THIS FILE:
% -------------------------------------------------------------
% One of my pet peeves with amsthm is its use of italics in the theorem and
% proposition environments; this makes math and text indistinguishable in said
% environments. To avoid this, I redefine the environments to use the standard
% font and to use a hanging indent, along with a bold vertical bar to its
% left, to distinguish these environments from surrounding text. (Along with
% the advantage of distinguishing math from text, this also allows nesting
% several such environments inside each other, like a definition inside a
% remark. I'm not sure how good of an idea this is, though. There are also
% downsides related to the hanging indentation, such as footnotes out of it
% being painful to do right.) This is done starting from the line
%   \theoremstyle{definition}
% and until the line
%   {\end{leftbar}\end{exmp}}

\documentclass[numbers=enddot,12pt,final,onecolumn,notitlepage]{scrartcl}%
\usepackage[headsepline,footsepline,manualmark]{scrlayer-scrpage}
\usepackage{amsfonts}
\usepackage{amssymb}
\usepackage{amsmath}
\usepackage{amsthm}
\usepackage{framed}
\usepackage{comment}
\usepackage[breaklinks=true]{hyperref}
\usepackage[sc]{mathpazo}
\usepackage[T1]{fontenc}
\usepackage{needspace}
%TCIDATA{OutputFilter=latex2.dll}
%TCIDATA{Version=5.50.0.2960}
%TCIDATA{LastRevised=Wednesday, August 30, 2023 17:03:21}
%TCIDATA{SuppressPackageManagement}
%TCIDATA{<META NAME="GraphicsSave" CONTENT="32">}
%TCIDATA{<META NAME="SaveForMode" CONTENT="1">}
%TCIDATA{BibliographyScheme=Manual}
%TCIDATA{Language=American English}
%BeginMSIPreambleData
\providecommand{\U}[1]{\protect\rule{.1in}{.1in}}
%EndMSIPreambleData
\theoremstyle{definition}
\newtheorem{theo}{Theorem}[section]
\newenvironment{theorem}[1][]
{\begin{theo}[#1]\begin{leftbar}}
{\end{leftbar}\end{theo}}
\newtheorem{lem}[theo]{Lemma}
\newenvironment{lemma}[1][]
{\begin{lem}[#1]\begin{leftbar}}
{\end{leftbar}\end{lem}}
\newtheorem{prop}[theo]{Proposition}
\newenvironment{proposition}[1][]
{\begin{prop}[#1]\begin{leftbar}}
{\end{leftbar}\end{prop}}
\newtheorem{defi}[theo]{Definition}
\newenvironment{definition}[1][]
{\begin{defi}[#1]\begin{leftbar}}
{\end{leftbar}\end{defi}}
\newtheorem{remk}[theo]{Remark}
\newenvironment{remark}[1][]
{\begin{remk}[#1]\begin{leftbar}}
{\end{leftbar}\end{remk}}
\newtheorem{coro}[theo]{Corollary}
\newenvironment{corollary}[1][]
{\begin{coro}[#1]\begin{leftbar}}
{\end{leftbar}\end{coro}}
\newtheorem{conv}[theo]{Convention}
\newenvironment{convention}[1][]
{\begin{conv}[#1]\begin{leftbar}}
{\end{leftbar}\end{conv}}
\newtheorem{exam}[theo]{Example}
\newenvironment{example}[1][]
{\begin{exam}[#1]\begin{leftbar}}
{\end{leftbar}\end{exam}}
\newenvironment{statement}{\begin{quote}}{\end{quote}}
\newenvironment{fineprint}{\begin{small}}{\end{small}}
\iffalse
\newenvironment{proof}[1][Proof]{\noindent\textbf{#1.} }{\ \rule{0.5em}{0.5em}}
\newenvironment{question}[1][Question]{\noindent\textbf{#1.} }{\ \rule{0.5em}{0.5em}}
\fi
\newcommand\arxiv[1]{\href{http://www.arxiv.org/abs/#1}{\texttt{arXiv:#1}}}
\let\sumnonlimits\sum
\renewcommand{\sum}{\sumnonlimits\limits}
\voffset=0cm
\hoffset=-0.7cm
\setlength\textheight{22.5cm}
\setlength\textwidth{15.5cm}
\newenvironment{verlong}{}{}
\newenvironment{vershort}{}{}
\newenvironment{noncompile}{}{}
\excludecomment{verlong}
\includecomment{vershort}
\excludecomment{noncompile}
\ihead{Some basic properties of compositions}
\ohead{page \thepage}
\cfoot{}
\begin{document}

\title{Some basic properties of compositions}
\author{Darij Grinberg and Ekaterina A. Vassilieva}
\date{30 July 2024}
\maketitle
\tableofcontents

\section*{***}

This is a companion note to \cite{main}. The purpose of this note is to prove
some elementary properties of integer compositions that are used in
\cite{main}. All of these proofs are elementary and generally quite easy, but
they are hard to find written down and often left to the reader to prove.

\section{Notations}

We let $\mathbb{N}=\left\{  0,1,2,\ldots\right\}  $.

A \emph{composition} means a finite list $\left(  \alpha_{1},\alpha_{2}%
,\ldots,\alpha_{k}\right)  $ of positive integers. The set of all compositions
will be denoted by $\operatorname*{Comp}$.

The \emph{empty composition }is defined to be the composition $\left(
{}\right)  $, which is a $0$-tuple. It is denoted by $\varnothing$.

The \emph{length }$\ell\left(  \alpha\right)  $ of a composition
$\alpha=\left(  \alpha_{1},\alpha_{2},\ldots,\alpha_{k}\right)  $ is defined
to be the number $k$.

If $\alpha=\left(  \alpha_{1},\alpha_{2},\ldots,\alpha_{k}\right)  $ is a
composition, then the nonnegative integer $\alpha_{1}+\alpha_{2}+\cdots
+\alpha_{k}$ is called the \emph{size} of $\alpha$ and is denoted by
$\left\vert \alpha\right\vert $. For any $n\in\mathbb{N}$, we define a
\emph{composition of }$n$ to be a composition that has size $n$. We let
$\operatorname*{Comp}\nolimits_{n}$ be the set of all compositions of $n$ (for
given $n\in\mathbb{N}$). The notation \textquotedblleft$\alpha\models
n$\textquotedblright\ is short for \textquotedblleft$\alpha\in
\operatorname*{Comp}\nolimits_{n}$\textquotedblright. For example, $\left(
1,5,2,1\right)  $ is a composition with size $9$ (since $\left\vert \left(
1,5,2,1\right)  \right\vert =1+5+2+1=9$), so that $\left(  1,5,2,1\right)
\in\operatorname*{Comp}\nolimits_{9}$, or, in other words, $\left(
1,5,2,1\right)  \models9$. Note that the empty composition $\varnothing$ is a
composition of $0$. In other words, $\varnothing\in\operatorname*{Comp}%
\nolimits_{0}$.

For any $n\in\mathbb{Z}$, we let $\left[  n\right]  $ denote the set $\left\{
1,2,\ldots,n\right\}  $. This set is empty whenever $n\leq0$, and otherwise
has size $n$.

If $X$ is any set, then $\mathcal{P}\left(  X\right)  $ shall denote the
\emph{powerset} of $X$. This is the set of all subsets of $X$.

\section{The maps $D$ and $\operatorname*{comp}$}

It is well-known that any positive integer $n$ has exactly $2^{n-1}$
compositions. This has a standard bijective proof (\textquotedblleft stars and
bars\textquotedblright) which relies on the following bijections:

\begin{definition}
\label{def.comps.D-comp}Let $n\in\mathbb{N}$.

\begin{enumerate}
\item[\textbf{(a)}] We define a map $D:\operatorname*{Comp}\nolimits_{n}%
\rightarrow\mathcal{P}\left(  \left[  n-1\right]  \right)  $ by
setting\footnotemark%
\begin{align*}
D\left(  \alpha_{1},\alpha_{2},\ldots,\alpha_{k}\right)   &  =\left\{
\alpha_{1}+\alpha_{2}+\cdots+\alpha_{i}\ \mid\ i\in\left[  k-1\right]
\right\} \\
&  =\left\{  \alpha_{1},\ \ \alpha_{1}+\alpha_{2},\ \ \alpha_{1}+\alpha
_{2}+\alpha_{3},\ \ \ldots,\ \ \alpha_{1}+\alpha_{2}+\cdots+\alpha
_{k-1}\right\}
\end{align*}
for each $\left(  \alpha_{1},\alpha_{2},\ldots,\alpha_{k}\right)
\in\operatorname*{Comp}\nolimits_{n}$. (It is easy to see that this map $D$ is
well-defined; see \cite[detailed version, Lemma 10.4]{Grinbe15} for a detailed proof.)

\item[\textbf{(b)}] We define a map $\operatorname*{comp}:\mathcal{P}\left(
\left[  n-1\right]  \right)  \rightarrow\operatorname*{Comp}\nolimits_{n}$ as
follows: For any $I\in\mathcal{P}\left(  \left[  n-1\right]  \right)  $, we
set%
\[
\operatorname*{comp}\left(  I\right)  =\left(  i_{1}-i_{0},\ \ i_{2}%
-i_{1},\ \ \ldots,\ \ i_{m}-i_{m-1}\right)  ,
\]
where $i_{0},i_{1},\ldots,i_{m}$ are the elements of the set $I\cup\left\{
0,n\right\}  $ listed in increasing order (so that $i_{0}<i_{1}<\cdots<i_{m}$,
therefore $i_{0}=0$ and $i_{m}=n$ and $\left\{  i_{1},i_{2},\ldots
,i_{m-1}\right\}  =I$). (It is easy to see that this map $\operatorname*{comp}%
$ is well-defined; see \cite[detailed version, Lemma 10.15 (d)]{Grinbe15} for
a detailed proof.)
\end{enumerate}

The maps $D$ and $\operatorname*{comp}$ are mutually inverse bijections. (See
\cite[detailed version, Proposition 10.17]{Grinbe15} for a detailed proof of this.)
\end{definition}

\footnotetext{The notation \textquotedblleft$D\left(  \alpha_{1},\alpha
_{2},\ldots,\alpha_{k}\right)  $\textquotedblright\ means $D\left(  \left(
\alpha_{1},\alpha_{2},\ldots,\alpha_{k}\right)  \right)  $ (that is, the image
of the composition $\left(  \alpha_{1},\alpha_{2},\ldots,\alpha_{k}\right)  $
under the map $D$).}We note that both of these maps $D$ and
$\operatorname*{comp}$ depend on $n$. Thus, they should be denoted by $D_{n}$
or $\operatorname*{comp}\nolimits_{n}$ to avoid ambiguity. Otherwise, for
example, the expression \textquotedblleft$\operatorname*{comp}\left(  \left\{
2,3\right\}  \right)  $\textquotedblright\ would have different meanings
depending on whether $n$ is $4$ or $5$. However, we shall not use the map
$\operatorname*{comp}$ in what follows. As for the map $D$, we need not be
afraid of any ambiguity, since the value of $D\left(  \alpha\right)  $ for a
given composition $\alpha$ is uniquely determined (indeed, the expression
\textquotedblleft$D\left(  \alpha\right)  $\textquotedblright\ makes sense
only for one value of $n$, namely for $n=\left\vert \alpha\right\vert $; no
other value of $n$ would satisfy $\alpha\in\operatorname*{Comp}\nolimits_{n}%
$). Thus, we shall freely use the notation \textquotedblleft$D\left(
\alpha\right)  $\textquotedblright\ without explicitly specifying $n$.

The notation $D$ we just introduced presumably originates in the word
\textquotedblleft\textbf{d}escent\textquotedblright, but the connection
between $D$ and actual descents is indirect and rather misleading. I prefer to
call $D$ the \textquotedblleft partial sum map\textquotedblright\ (as
$D\left(  \alpha\right)  $ consists of the partial sums of the composition
$\alpha$) and its inverse $\operatorname*{comp}$ the \textquotedblleft
interstitial map\textquotedblright\ (as $\operatorname*{comp}\left(  I\right)
$ consists of the lengths of the intervals into which the elements of $I$
split the interval $\left[  n\right]  $).

\begin{example}
Let $n=10$.

\begin{enumerate}
\item[\textbf{(a)}] The map $D$ defined in Definition \ref{def.comps.D-comp}
\textbf{(a)} satisfies%
\begin{align*}
D\left(  1,4,2,3\right)   &  =\left\{  1,\ \ 1+4,\ \ 1+4+2\right\}  =\left\{
1,5,7\right\}  ;\\
D\left(  3,5,2\right)   &  =\left\{  3,\ \ 3+5\right\}  =\left\{  3,8\right\}
;\\
D\left(  1,1,1,1,1,1,1,1,1,1\right)   &  =\left\{  1,2,3,4,5,6,7,8,9\right\}
=\left[  9\right]  =\left[  n-1\right]  ;\\
D\left(  10\right)   &  =\left\{  {}\right\}  =\varnothing.
\end{align*}


\item[\textbf{(b)}] The map $\operatorname*{comp}$ defined in Definition
\ref{def.comps.D-comp} \textbf{(b)} satisfies%
\[
\operatorname*{comp}\left(  \left\{  2,3,7\right\}  \right)  =\left(
2-0,\ \ 3-2,\ \ 7-3,\ \ 10-7\right)  =\left(  2,1,4,3\right)
\]
(since $0,2,3,7,10$ are the elements of the set $\left\{  2,3,7\right\}
\cup\left\{  0,10\right\}  $ listed in increasing order).
\end{enumerate}
\end{example}

Our first observation about the bijections $D$ and $\operatorname*{comp}$
concerns the relation between the size of $D\left(  \alpha\right)  $ and the
length $\ell\left(  \alpha\right)  $ of $\alpha$. Namely, we shall show that
every composition $\alpha$ of size $\left\vert \alpha\right\vert >0$ satisfies
$\left\vert D\left(  \alpha\right)  \right\vert =\ell\left(  \alpha\right)
-1$:

\begin{proposition}
\label{prop.D-size.pos}Let $\alpha$ be a composition such that $\left\vert
\alpha\right\vert >0$. Then, $\left\vert D\left(  \alpha\right)  \right\vert
=\ell\left(  \alpha\right)  -1$.
\end{proposition}

Note that the \textquotedblleft$\left\vert \alpha\right\vert >0$%
\textquotedblright\ assumption in Proposition \ref{prop.D-size.pos} is
necessary, since Proposition \ref{prop.D-size.pos} would fail if $\alpha$ was
the empty composition $\varnothing=\left(  {}\right)  $ (because $D\left(
\varnothing\right)  =\varnothing$ and thus $\left\vert D\left(  \varnothing
\right)  \right\vert =0\neq\ell\left(  \varnothing\right)  -1$).

\begin{proof}
[Proof of Proposition \ref{prop.D-size.pos}.]Write the composition $\alpha$ in
the form $\alpha=\left(  \alpha_{1},\alpha_{2},\ldots,\alpha_{k}\right)  $.
Then, $\ell\left(  \alpha\right)  =k$ (by the definition of $\ell\left(
\alpha\right)  $) and $\left\vert \alpha\right\vert =\alpha_{1}+\alpha
_{2}+\cdots+\alpha_{k}$ (by the definition of $\left\vert \alpha\right\vert
$). If we had $k=0$, then we would have%
\begin{align*}
\left\vert \alpha\right\vert  &  =\alpha_{1}+\alpha_{2}+\cdots+\alpha
_{k}=\alpha_{1}+\alpha_{2}+\cdots+\alpha_{0}\ \ \ \ \ \ \ \ \ \ \left(
\text{since }k=0\right) \\
&  =\left(  \text{empty sum}\right)  =0,
\end{align*}
which would contradict $\left\vert \alpha\right\vert >0$. Thus, we cannot have
$k=0$. Hence, $k\neq0$, so that $k\geq1$ (since $k\in\mathbb{N}$).

From $\alpha=\left(  \alpha_{1},\alpha_{2},\ldots,\alpha_{k}\right)  $, we
obtain%
\begin{align}
D\left(  \alpha\right)   &  =D\left(  \alpha_{1},\alpha_{2},\ldots,\alpha
_{k}\right) \nonumber\\
&  =\left\{  \alpha_{1},\ \ \alpha_{1}+\alpha_{2},\ \ \alpha_{1}+\alpha
_{2}+\alpha_{3},\ \ \ldots,\ \ \alpha_{1}+\alpha_{2}+\cdots+\alpha
_{k-1}\right\}  \label{pf.prop.D-size.pos.4}%
\end{align}
(by the definition of the map $D$). However, it is easy to see that the chain
of inequalities
\[
\alpha_{1}<\alpha_{1}+\alpha_{2}<\alpha_{1}+\alpha_{2}+\alpha_{3}%
<\cdots<\alpha_{1}+\alpha_{2}+\cdots+\alpha_{k-1}%
\]
holds\footnote{\textit{Proof.} Let $i\in\left[  k\right]  $. Then, $\alpha
_{i}$ is an entry of $\alpha$ (since $\alpha=\left(  \alpha_{1},\alpha
_{2},\ldots,\alpha_{k}\right)  $).
\par
Recall that $\alpha$ is a composition, i.e., a finite list of positive
integers. Hence, $\alpha_{i}$ is a positive integer (since $\alpha_{i}$ is an
entry of $\alpha$). Therefore, $\alpha_{i}>0$. Hence,%
\[
\alpha_{1}+\alpha_{2}+\cdots+\alpha_{i}=\alpha_{1}+\alpha_{2}+\cdots
+\alpha_{i-1}+\underbrace{\alpha_{i}}_{>0}>\alpha_{1}+\alpha_{2}+\cdots
+\alpha_{i-1}.
\]
In other words, $\alpha_{1}+\alpha_{2}+\cdots+\alpha_{i-1}<\alpha_{1}%
+\alpha_{2}+\cdots+\alpha_{i}$.
\par
Forget that we fixed $i$. We thus have proved the inequality $\alpha
_{1}+\alpha_{2}+\cdots+\alpha_{i-1}<\alpha_{1}+\alpha_{2}+\cdots+\alpha_{i}$
for each $i\in\left[  k\right]  $. Hence, in particular, this inequality holds
for each $i\in\left\{  2,3,\ldots,k-1\right\}  $. In other words, we have the
chain of inequalities
\[
\alpha_{1}<\alpha_{1}+\alpha_{2}<\alpha_{1}+\alpha_{2}+\alpha_{3}%
<\cdots<\alpha_{1}+\alpha_{2}+\cdots+\alpha_{k-1}.
\]
}. Thus, the $k-1$ numbers $\alpha_{1},\ \ \alpha_{1}+\alpha_{2}%
,\ \ \alpha_{1}+\alpha_{2}+\alpha_{3},\ \ \ldots,\ \ \alpha_{1}+\alpha
_{2}+\cdots+\alpha_{k-1}$ are distinct. Therefore, the set of these $k-1$
numbers has size $k-1$. In other words, we have
\[
\left\vert \left\{  \alpha_{1},\ \ \alpha_{1}+\alpha_{2},\ \ \alpha_{1}%
+\alpha_{2}+\alpha_{3},\ \ \ldots,\ \ \alpha_{1}+\alpha_{2}+\cdots
+\alpha_{k-1}\right\}  \right\vert =k-1.
\]
In view of (\ref{pf.prop.D-size.pos.4}), we can rewrite this as $\left\vert
D\left(  \alpha\right)  \right\vert =k-1$. In other words, $\left\vert
D\left(  \alpha\right)  \right\vert =\ell\left(  \alpha\right)  -1$ (since
$\ell\left(  \alpha\right)  =k$). This proves Proposition
\ref{prop.D-size.pos}.
\end{proof}

The analogue of Proposition \ref{prop.D-size.pos} for $\left\vert
\alpha\right\vert =0$ is almost trivial:

\begin{proposition}
\label{prop.D-size.0}Let $\alpha$ be a composition such that $\left\vert
\alpha\right\vert =0$. Then, $\alpha=\varnothing$ and $\ell\left(
\alpha\right)  =0$ and $D\left(  \alpha\right)  =\varnothing$.
\end{proposition}

\begin{proof}
[Proof of Proposition \ref{prop.D-size.0}.]Write the composition $\alpha$ in
the form $\alpha=\left(  \alpha_{1},\alpha_{2},\ldots,\alpha_{k}\right)  $.
Then, $\left\vert \alpha\right\vert =\alpha_{1}+\alpha_{2}+\cdots+\alpha_{k}$
(by the definition of $\left\vert \alpha\right\vert $) and $\ell\left(
\alpha\right)  =k$ (by the definition of $\ell\left(  \alpha\right)  $).

Assume (for the sake of contradiction) that $k\neq0$. Thus, $k\geq1$ (since
$k\in\mathbb{N}$).

However, $\alpha$ is a composition, i.e., a finite list of positive integers.
In other words, $\left(  \alpha_{1},\alpha_{2},\ldots,\alpha_{k}\right)  $ is
a finite list of positive integers (since $\alpha=\left(  \alpha_{1}%
,\alpha_{2},\ldots,\alpha_{k}\right)  $). Thus, $\alpha_{1},\alpha_{2}%
,\ldots,\alpha_{k}$ are positive integers. Therefore, in particular,
$\alpha_{2},\alpha_{3},\ldots,\alpha_{k}$ are positive integers. Hence,
$\alpha_{2}+\alpha_{3}+\cdots+\alpha_{k}\geq0$ (since a sum of positive
integers is always $\geq0$). However, from $\left\vert \alpha\right\vert =0$,
we obtain%
\begin{align*}
0  &  =\left\vert \alpha\right\vert =\alpha_{1}+\alpha_{2}+\cdots+\alpha
_{k}=\alpha_{1}+\underbrace{\left(  \alpha_{2}+\alpha_{3}+\cdots+\alpha
_{k}\right)  }_{\geq0}\ \ \ \ \ \ \ \ \ \ \left(  \text{since }k\geq1\right)
\\
&  \geq\alpha_{1}>0\ \ \ \ \ \ \ \ \ \ \left(  \text{since }\alpha_{1}\text{
is a positive integer}\right)  ,
\end{align*}
which is absurd. This contradiction shows that our assumption (that $k\neq0$)
was false. Hence, $k=0$.

Now,
\begin{align*}
\alpha &  =\left(  \alpha_{1},\alpha_{2},\ldots,\alpha_{k}\right)  =\left(
\alpha_{1},\alpha_{2},\ldots,\alpha_{0}\right)  \ \ \ \ \ \ \ \ \ \ \left(
\text{since }k=0\right) \\
&  =\left(  {}\right)  =\varnothing\ \ \ \ \ \ \ \ \ \ \left(  \text{recall
that }\varnothing\text{ denotes the empty composition}\right)  .
\end{align*}
Moreover, $\ell\left(  \alpha\right)  =k=0$. Finally, from $\alpha=\left(
{}\right)  $, we obtain $D\left(  \alpha\right)  =D\left(  {}\right)
=\varnothing$ (by the definition of the map $D:\operatorname*{Comp}%
\nolimits_{0}\rightarrow\mathcal{P}\left(  \left[  0-1\right]  \right)  $).
Thus, Proposition \ref{prop.D-size.0} is proved.
\end{proof}

We can unite Proposition \ref{prop.D-size.pos} with Proposition
\ref{prop.D-size.0} by using the \emph{Iverson bracket notation}:

\begin{convention}
\label{conv.iverson}If $\mathcal{A}$ is a logical statement, then $\left[
\mathcal{A}\right]  $ shall denote the truth value of $\mathcal{A}$; this is
the integer defined by%
\[
\left[  \mathcal{A}\right]  =%
\begin{cases}
1, & \text{if }\mathcal{A}\text{ is true;}\\
0, & \text{if }\mathcal{A}\text{ is false.}%
\end{cases}
\]


For example, $\left[  2+2=4\right]  =1$ (since the statement $2+2=4$ is true)
and $\left[  2+2=5\right]  =0$ (since the statement $2+2=5$ is false).
\end{convention}

Now, Proposition \ref{prop.D-size.pos} with Proposition \ref{prop.D-size.0}
can be combined into the following:

\begin{corollary}
\label{cor.D-size.1}Let $n\in\mathbb{N}$. Let $\alpha\in\operatorname*{Comp}%
\nolimits_{n}$. Then, $\ell\left(  \alpha\right)  =\left\vert D\left(
\alpha\right)  \right\vert +\left[  n\neq0\right]  $.
\end{corollary}

\begin{proof}
[Proof of Corollary \ref{cor.D-size.1}.]From $\alpha\in\operatorname*{Comp}%
\nolimits_{n}$, we see that $\alpha$ is a composition of $n$ (since
$\operatorname*{Comp}\nolimits_{n}$ is the set of all compositions of $n$). In
other words, $\alpha$ is a composition having size $n$. Therefore, $\left\vert
\alpha\right\vert =n$ (since $\left\vert \alpha\right\vert $ is the size of
$\alpha$, but we know that $\alpha$ has size $n$).

We are in one of the following two cases:

\textit{Case 1:} We have $n=0$.

\textit{Case 2:} We have $n\neq0$.

Let us first consider Case 1. In this case, we have $n=0$. Hence, we don't
have $n\neq0$. Thus, $\left[  n\neq0\right]  =0$.

However, $\left\vert \alpha\right\vert =n=0$. Thus, Proposition
\ref{prop.D-size.0} yields $\alpha=\varnothing$ and $\ell\left(
\alpha\right)  =0$ and $D\left(  \alpha\right)  =\varnothing$. From $D\left(
\alpha\right)  =\varnothing$, we obtain $\left\vert D\left(  \alpha\right)
\right\vert =\left\vert \varnothing\right\vert =0$. Thus,
$\underbrace{\left\vert D\left(  \alpha\right)  \right\vert }_{=0}%
+\underbrace{\left[  n\neq0\right]  }_{=0}=0$. Comparing this with
$\ell\left(  \alpha\right)  =0$, we obtain $\ell\left(  \alpha\right)
=\left\vert D\left(  \alpha\right)  \right\vert +\left[  n\neq0\right]  $.
Hence, Corollary \ref{cor.D-size.1} is proved in Case 1.

Let us now consider Case 2. In this case, we have $n\neq0$. Hence, $\left[
n\neq0\right]  =1$. Also, from $n\neq0$, we obtain $n>0$ (since $n\in
\mathbb{N}$). Thus, $\left\vert \alpha\right\vert =n>0$. Hence, Proposition
\ref{prop.D-size.pos} yields $\left\vert D\left(  \alpha\right)  \right\vert
=\ell\left(  \alpha\right)  -1$. Hence, $\ell\left(  \alpha\right)
=\left\vert D\left(  \alpha\right)  \right\vert +1$. Comparing this with
$\left\vert D\left(  \alpha\right)  \right\vert +\underbrace{\left[
n\neq0\right]  }_{=1}=\left\vert D\left(  \alpha\right)  \right\vert +1$, we
obtain $\ell\left(  \alpha\right)  =\left\vert D\left(  \alpha\right)
\right\vert +\left[  n\neq0\right]  $. Thus, Corollary \ref{cor.D-size.1} is
proved in Case 2.

We have now proved Corollary \ref{cor.D-size.1} in both Cases 1 and 2. Hence,
Corollary \ref{cor.D-size.1} always holds.
\end{proof}

\begin{corollary}
\label{cor.D-size.2}Let $n\in\mathbb{N}$. Let $\alpha\in\operatorname*{Comp}%
\nolimits_{n}$ and $\beta\in\operatorname*{Comp}\nolimits_{n}$. Then,
$\ell\left(  \beta\right)  -\ell\left(  \alpha\right)  =\left\vert D\left(
\beta\right)  \right\vert -\left\vert D\left(  \alpha\right)  \right\vert $.
\end{corollary}

\begin{proof}
[Proof of Corollary \ref{cor.D-size.2}.]Corollary \ref{cor.D-size.1} yields
$\ell\left(  \alpha\right)  =\left\vert D\left(  \alpha\right)  \right\vert
+\left[  n\neq0\right]  $. Corollary \ref{cor.D-size.1} (applied to $\beta$
instead of $\alpha$) yields $\ell\left(  \beta\right)  =\left\vert D\left(
\beta\right)  \right\vert +\left[  n\neq0\right]  $. Hence,%
\begin{align*}
\underbrace{\ell\left(  \beta\right)  }_{=\left\vert D\left(  \beta\right)
\right\vert +\left[  n\neq0\right]  }-\underbrace{\ell\left(  \alpha\right)
}_{=\left\vert D\left(  \alpha\right)  \right\vert +\left[  n\neq0\right]  }
&  =\left(  \left\vert D\left(  \beta\right)  \right\vert +\left[
n\neq0\right]  \right)  -\left(  \left\vert D\left(  \alpha\right)
\right\vert +\left[  n\neq0\right]  \right) \\
&  =\left\vert D\left(  \beta\right)  \right\vert -\left\vert D\left(
\alpha\right)  \right\vert .
\end{align*}
This proves Corollary \ref{cor.D-size.2}.
\end{proof}

\section{Reversals}

We shall now discuss a certain operation on compositions:

\begin{definition}
\label{def.rev-comp}If $\alpha=\left(  \alpha_{1},\alpha_{2},\ldots,\alpha
_{k}\right)  $ is a composition, then the \emph{reversal} of $\alpha$ is
defined to be the composition $\left(  \alpha_{k},\alpha_{k-1},\ldots
,\alpha_{1}\right)  $. It is denoted by $\operatorname*{rev}\alpha$.

Thus, we have defined a map $\operatorname*{rev}:\operatorname*{Comp}%
\rightarrow\operatorname*{Comp}$ that sends each composition $\alpha$ to the
composition $\operatorname*{rev}\alpha$.
\end{definition}

\begin{example}
We have%
\begin{align*}
\operatorname*{rev}\left(  2,3,6\right)   &  =\left(  6,3,2\right)  ;\\
\operatorname*{rev}\left(  4,1,1,2\right)   &  =\left(  2,1,1,4\right)  ;\\
\operatorname*{rev}\varnothing &  =\varnothing.
\end{align*}

\end{example}

\begin{proposition}
\label{prop.rev-comp-size}Let $\alpha\in\operatorname*{Comp}$. Then,
$\left\vert \operatorname*{rev}\alpha\right\vert =\left\vert \alpha\right\vert
$.
\end{proposition}

\begin{proof}
[Proof of Proposition \ref{prop.rev-comp-size}.]Write the composition $\alpha$
in the form $\alpha=\left(  \alpha_{1},\alpha_{2},\ldots,\alpha_{k}\right)  $.
Then, $\operatorname*{rev}\alpha=\left(  \alpha_{k},\alpha_{k-1},\ldots
,\alpha_{1}\right)  $ (by Definition \ref{def.rev-comp}) and $\left\vert
\alpha\right\vert =\alpha_{1}+\alpha_{2}+\cdots+\alpha_{k}$ (by the definition
of $\left\vert \alpha\right\vert $). Now,%
\begin{align*}
\left\vert \operatorname*{rev}\alpha\right\vert  &  =\left\vert \left(
\alpha_{k},\alpha_{k-1},\ldots,\alpha_{1}\right)  \right\vert
\ \ \ \ \ \ \ \ \ \ \left(  \text{since }\operatorname*{rev}\alpha=\left(
\alpha_{k},\alpha_{k-1},\ldots,\alpha_{1}\right)  \right) \\
&  =\alpha_{k}+\alpha_{k-1}+\cdots+\alpha_{1}\ \ \ \ \ \ \ \ \ \ \left(
\text{by the definition of }\left\vert \left(  \alpha_{k},\alpha_{k-1}%
,\ldots,\alpha_{1}\right)  \right\vert \right) \\
&  =\alpha_{1}+\alpha_{2}+\cdots+\alpha_{k}\\
&  =\left\vert \alpha\right\vert \ \ \ \ \ \ \ \ \ \ \left(  \text{since
}\left\vert \alpha\right\vert =\alpha_{1}+\alpha_{2}+\cdots+\alpha_{k}\right)
.
\end{align*}
This proves Proposition \ref{prop.rev-comp-size}.
\end{proof}

\begin{proposition}
\label{prop.rev-rev}Let $\alpha\in\operatorname*{Comp}$. Then,
$\operatorname*{rev}\left(  \operatorname*{rev}\alpha\right)  =\alpha$.
\end{proposition}

\begin{proof}
[Proof of Proposition \ref{prop.rev-rev}.]Write the composition $\alpha$ in
the form $\alpha=\left(  \alpha_{1},\alpha_{2},\ldots,\alpha_{k}\right)  $.
Then, Definition \ref{def.rev-comp} yields $\operatorname*{rev}\alpha=\left(
\alpha_{k},\alpha_{k-1},\ldots,\alpha_{1}\right)  $. However, Definition
\ref{def.rev-comp} also yields $\operatorname*{rev}\left(  \alpha_{k}%
,\alpha_{k-1},\ldots,\alpha_{1}\right)  =\left(  \alpha_{1},\alpha_{2}%
,\ldots,\alpha_{k}\right)  $. Now,%
\[
\operatorname*{rev}\underbrace{\left(  \operatorname*{rev}\alpha\right)
}_{=\left(  \alpha_{k},\alpha_{k-1},\ldots,\alpha_{1}\right)  }%
=\operatorname*{rev}\left(  \alpha_{k},\alpha_{k-1},\ldots,\alpha_{1}\right)
=\left(  \alpha_{1},\alpha_{2},\ldots,\alpha_{k}\right)  =\alpha.
\]
This proves Proposition \ref{prop.rev-rev}.
\end{proof}

\begin{corollary}
\label{cor.rev-bij}The map%
\begin{align*}
\operatorname*{Comp}  &  \rightarrow\operatorname*{Comp},\\
\delta &  \mapsto\operatorname*{rev}\delta
\end{align*}
is a bijection.
\end{corollary}

\begin{proof}
[Proof of Corollary \ref{cor.rev-bij}.]Let us denote this map by
$\operatorname*{rev}$ (since the image of any $\delta\in\operatorname*{Comp}$
under this map is already being called $\operatorname*{rev}\delta$). Thus, we
must prove that this map $\operatorname*{rev}$ is a bijection.

But this is easy: Every $\alpha\in\operatorname*{Comp}$ satisfies
\begin{align*}
\left(  \operatorname*{rev}\circ\operatorname*{rev}\right)  \left(
\alpha\right)   &  =\operatorname*{rev}\left(  \operatorname*{rev}%
\alpha\right)  =\alpha\ \ \ \ \ \ \ \ \ \ \left(  \text{by Proposition
\ref{prop.rev-rev}}\right) \\
&  =\operatorname*{id}\left(  \alpha\right)  .
\end{align*}
Thus, $\operatorname*{rev}\circ\operatorname*{rev}=\operatorname*{id}$. Hence,
the map $\operatorname*{rev}$ is inverse to itself. Thus, the map
$\operatorname*{rev}$ is invertible, i.e., bijective. In other words, it is a
bijection. This proves Corollary \ref{cor.rev-bij}.
\end{proof}

We also define a related operation on subsets of $\left[  n-1\right]  $:

\begin{definition}
\label{def.rev-subset}Let $n\in\mathbb{N}$. For any subset $X$ of $\left[
n-1\right]  $, we let $\operatorname*{rev}\nolimits_{n}X$ denote the set
$\left\{  n-x\ \mid\ x\in X\right\}  $.
\end{definition}

\begin{example}
If $n=7$, then%
\begin{align*}
\operatorname*{rev}\nolimits_{n}\left(  \left\{  2,4\right\}  \right)   &
=\left\{  7-2,\ 7-4\right\}  =\left\{  5,3\right\}  =\left\{  3,5\right\}  ;\\
\operatorname*{rev}\nolimits_{n}\left(  \left\{  1,2,5,6\right\}  \right)   &
=\left\{  7-1,\ 7-2,\ 7-5,\ 7-6\right\}  =\left\{  6,5,2,1\right\}  =\left\{
1,2,5,6\right\}  ;\\
\operatorname*{rev}\nolimits_{n}\left(  \varnothing\right)   &  =\varnothing
;\\
\operatorname*{rev}\nolimits_{n}\left(  \left[  6\right]  \right)   &
=\left[  6\right]  .
\end{align*}

\end{example}

Informally speaking, the set $\operatorname*{rev}\nolimits_{n}X$ defined in
Definition \ref{def.rev-subset} is the reflection of the set $X$ across the
midpoint of the interval $\left[  n-1\right]  $ (where we regard numbers as
points on the number line). From this point of view, all claims of the
following theorem are visually obvious:

\begin{theorem}
\label{thm.rev-rev}Let $n\in\mathbb{N}$. Then:

\begin{enumerate}
\item[\textbf{(a)}] We have $\operatorname*{rev}\nolimits_{n}X\subseteq\left[
n-1\right]  $ for each subset $X$ of $\left[  n-1\right]  $.

\item[\textbf{(b)}] We have $\operatorname*{rev}\nolimits_{n}\left(
\operatorname*{rev}\nolimits_{n}X\right)  =X$ for any subset $X$ of $\left[
n-1\right]  $.

\item[\textbf{(c)}] If two subsets $X$ and $Y$ of $\left[  n-1\right]  $
satisfy $X\subseteq Y$, then $\operatorname*{rev}\nolimits_{n}X\subseteq
\operatorname*{rev}\nolimits_{n}Y$.

\item[\textbf{(d)}] We have $\left\vert \operatorname*{rev}\nolimits_{n}%
X\right\vert =\left\vert X\right\vert $ for any subset $X$ of $\left[
n-1\right]  $.

\item[\textbf{(e)}] We have $\operatorname*{rev}\nolimits_{n}X=\left\{
i\in\left[  n-1\right]  \ \mid\ n-i\in X\right\}  $ for any subset $X$ of
$\left[  n-1\right]  $.

\item[\textbf{(f)}] We have $\operatorname*{rev}\nolimits_{n}\left(
X\setminus Y\right)  =\left(  \operatorname*{rev}\nolimits_{n}X\right)
\setminus\left(  \operatorname*{rev}\nolimits_{n}Y\right)  $ for any subsets
$X$ and $Y$ of $\left[  n-1\right]  $.

\item[\textbf{(g)}] We have $\operatorname*{rev}\nolimits_{n}\left(  X\cap
Y\right)  =\left(  \operatorname*{rev}\nolimits_{n}X\right)  \cap\left(
\operatorname*{rev}\nolimits_{n}Y\right)  $ for any subsets $X$ and $Y$ of
$\left[  n-1\right]  $.

\item[\textbf{(h)}] We have $\operatorname*{rev}\nolimits_{n}\left(  \left[
n-1\right]  \right)  =\left[  n-1\right]  $.

\item[\textbf{(i)}] We have $D\left(  \operatorname*{rev}\alpha\right)
=\operatorname*{rev}\nolimits_{n}\left(  D\left(  \alpha\right)  \right)  $
for any composition $\alpha\in\operatorname*{Comp}\nolimits_{n}$.
\end{enumerate}
\end{theorem}

\begin{proof}
[Proof of Theorem \ref{thm.rev-rev}.]\textbf{(a)} Let $X$ be a subset of
$\left[  n-1\right]  $. Then, $n-x\in\left[  n-1\right]  $ for each $x\in
X$\ \ \ \ \footnote{\textit{Proof.} Let $x\in X$. Then, $x\in X\subseteq
\left[  n-1\right]  =\left\{  1,2,\ldots,n-1\right\}  $, so that
$n-x\in\left\{  1,2,\ldots,n-1\right\}  =\left[  n-1\right]  $.
\par
Forget that we fixed $x$. We thus have shown that $n-x\in\left[  n-1\right]  $
for each $x\in X$.}. In other words,%
\[
\left\{  n-x\ \mid\ x\in X\right\}  \subseteq\left[  n-1\right]  .
\]
This rewrites as $\operatorname*{rev}\nolimits_{n}X\subseteq\left[
n-1\right]  $ (since $\operatorname*{rev}\nolimits_{n}X$ is defined to be
$\left\{  n-x\ \mid\ x\in X\right\}  $). This proves Theorem \ref{thm.rev-rev}
\textbf{(a)}. \medskip

\textbf{(b)} Let $X$ be a subset of $\left[  n-1\right]  $. Let
$Y=\operatorname*{rev}\nolimits_{n}X$.

Let $p\in\operatorname*{rev}\nolimits_{n}Y$. We shall show that $p\in X$.

We have%
\begin{align*}
p  &  \in\operatorname*{rev}\nolimits_{n}Y=\left\{  n-x\ \mid\ x\in Y\right\}
\ \ \ \ \ \ \ \ \ \ \left(  \text{by the definition of }\operatorname*{rev}%
\nolimits_{n}Y\right) \\
&  =\left\{  n-y\ \mid\ y\in Y\right\}  \ \ \ \ \ \ \ \ \ \ \left(
\text{here, we have renamed the index }x\text{ as }y\right)  .
\end{align*}
In other words, $p=n-y$ for some $y\in Y$. Consider this $y$. Now,%
\[
y\in Y=\operatorname*{rev}\nolimits_{n}X=\left\{  n-x\ \mid\ x\in X\right\}
\ \ \ \ \ \ \ \ \ \ \left(  \text{by the definition of }\operatorname*{rev}%
\nolimits_{n}X\right)  .
\]
In other words, $y=n-x$ for some $x\in X$. Consider this $x$. Now,
$p=n-\underbrace{y}_{=n-x}=n-\left(  n-x\right)  =x\in X$.

Forget that we fixed $p$. We thus have shown that $p\in X$ for each
$p\in\operatorname*{rev}\nolimits_{n}Y$. In other words, $\operatorname*{rev}%
\nolimits_{n}Y\subseteq X$.

On the other hand, let $q\in X$. Then, $n-q$ has the form $n-x$ for some $x\in
X$ (namely, for $x=q$). In other words, $n-q\in\left\{  n-x\ \mid\ x\in
X\right\}  $. Since $Y=\operatorname*{rev}\nolimits_{n}X=\left\{
n-x\ \mid\ x\in X\right\}  $ (by the definition of $\operatorname*{rev}%
\nolimits_{n}X$), we can rewrite this as $n-q\in Y$.

Furthermore, $q=n-\left(  n-q\right)  $. Hence, $q$ has the form $n-x$ for
some $x\in Y$ (namely, for $x=n-q$). In other words, $q\in\left\{
n-x\ \mid\ x\in Y\right\}  $. Since $\operatorname*{rev}\nolimits_{n}%
Y=\left\{  n-x\ \mid\ x\in Y\right\}  $ (by the definition of
$\operatorname*{rev}\nolimits_{n}Y$), we can rewrite this as $q\in
\operatorname*{rev}\nolimits_{n}Y$.

Forget that we fixed $q$. We thus have shown that $q\in\operatorname*{rev}%
\nolimits_{n}Y$ for each $q\in X$. In other words, $X\subseteq
\operatorname*{rev}\nolimits_{n}Y$.

Combining this with $\operatorname*{rev}\nolimits_{n}Y\subseteq X$, we obtain
$\operatorname*{rev}\nolimits_{n}Y=X$. In other words, $\operatorname*{rev}%
\nolimits_{n}\left(  \operatorname*{rev}\nolimits_{n}X\right)  =X$ (since
$Y=\operatorname*{rev}\nolimits_{n}X$). This proves Theorem \ref{thm.rev-rev}
\textbf{(b)}. \medskip

\textbf{(c)} Let $X$ and $Y$ be two subsets of $\left[  n-1\right]  $ that
satisfy $X\subseteq Y$. The definition of $\operatorname*{rev}\nolimits_{n}Y$
yields $\operatorname*{rev}\nolimits_{n}Y=\left\{  n-x\ \mid\ x\in Y\right\}
$.

Let $p\in\operatorname*{rev}\nolimits_{n}X$. Then, $p\in\operatorname*{rev}%
\nolimits_{n}X=\left\{  n-x\ \mid\ x\in X\right\}  $ (by the definition of
$\operatorname*{rev}\nolimits_{n}X$). In other words, $p=n-x$ for some $x\in
X$. Consider this $x$, and denote it by $z$. Thus, $z\in X$ and $p=n-z$.

Now, $z\in X\subseteq Y$ and $p=n-z$. Therefore, $p=n-x$ for some $x\in Y$
(namely, for $x=z$). In other words, $p\in\left\{  n-x\ \mid\ x\in Y\right\}
$. This rewrites as $p\in\operatorname*{rev}\nolimits_{n}Y$ (since
$\operatorname*{rev}\nolimits_{n}Y=\left\{  n-x\ \mid\ x\in Y\right\}  $).

Forget that we fixed $p$. We thus have shown that $p\in\operatorname*{rev}%
\nolimits_{n}Y$ for each $p\in\operatorname*{rev}\nolimits_{n}X$. In other
words, $\operatorname*{rev}\nolimits_{n}X\subseteq\operatorname*{rev}%
\nolimits_{n}Y$. This proves Theorem \ref{thm.rev-rev} \textbf{(c)}. \medskip

\textbf{(d)} Let $X$ be a subset of $\left[  n-1\right]  $. Let
$Y=\operatorname*{rev}\nolimits_{n}X$.

The definition of $\operatorname*{rev}\nolimits_{n}X$ yields
$\operatorname*{rev}\nolimits_{n}X=\left\{  n-x\ \mid\ x\in X\right\}  $.
Thus, the elements of $\operatorname*{rev}\nolimits_{n}X$ are precisely the
numbers $n-x$ for $x\in X$. Clearly, there are at most $\left\vert
X\right\vert $ many such numbers (since there are $\left\vert X\right\vert $
many elements $x\in X$). Hence, the set $\operatorname*{rev}\nolimits_{n}X$
has at most $\left\vert X\right\vert $ many elements. In other words,
$\left\vert \operatorname*{rev}\nolimits_{n}X\right\vert \leq\left\vert
X\right\vert $.

The same argument (applied to $Y$ instead of $X$) yields $\left\vert
\operatorname*{rev}\nolimits_{n}Y\right\vert \leq\left\vert Y\right\vert $.
However, from $Y=\operatorname*{rev}\nolimits_{n}X$, we obtain
$\operatorname*{rev}\nolimits_{n}Y=\operatorname*{rev}\nolimits_{n}\left(
\operatorname*{rev}\nolimits_{n}X\right)  =X$ (by Theorem \ref{thm.rev-rev}
\textbf{(b)}). In view of this, we can rewrite $\left\vert \operatorname*{rev}%
\nolimits_{n}Y\right\vert \leq\left\vert Y\right\vert $ as $\left\vert
X\right\vert \leq\left\vert Y\right\vert $.

But from $Y=\operatorname*{rev}\nolimits_{n}X$, we also obtain $\left\vert
Y\right\vert =\left\vert \operatorname*{rev}\nolimits_{n}X\right\vert
\leq\left\vert X\right\vert $. Combining this inequality with $\left\vert
X\right\vert \leq\left\vert Y\right\vert $, we find $\left\vert X\right\vert
=\left\vert Y\right\vert =\left\vert \operatorname*{rev}\nolimits_{n}%
X\right\vert $. In other words, $\left\vert \operatorname*{rev}\nolimits_{n}%
X\right\vert =\left\vert X\right\vert $. This proves Theorem \ref{thm.rev-rev}
\textbf{(d)}. \medskip

\textbf{(e)} Let $X$ be a subset of $\left[  n-1\right]  $. Let $Y=\left\{
i\in\left[  n-1\right]  \ \mid\ n-i\in X\right\}  $. We shall show that
$\operatorname*{rev}\nolimits_{n}X=Y$.

Note that $\operatorname*{rev}\nolimits_{n}X=\left\{  n-x\ \mid\ x\in
X\right\}  $ (by the definition of $\operatorname*{rev}\nolimits_{n}X$).

Let $p\in\operatorname*{rev}\nolimits_{n}X$. Then, $p\in\operatorname*{rev}%
\nolimits_{n}X=\left\{  n-x\ \mid\ x\in X\right\}  $. In other words, $p=n-x$
for some $x\in X$. Consider this $x$. Thus, $p=n-x$, so that $n=p+x$.
Therefore, $n-p=x\in X$. Also, $p\in\operatorname*{rev}\nolimits_{n}%
X\subseteq\left[  n-1\right]  $ (by Theorem \ref{thm.rev-rev} \textbf{(a)}).
Hence, $p$ is an element $i$ of $\left[  n-1\right]  $ satisfying $n-i\in X$
(since $n-p\in X$). In other words, $p\in\left\{  i\in\left[  n-1\right]
\ \mid\ n-i\in X\right\}  $. In other words, $p\in Y$ (since $Y=\left\{
i\in\left[  n-1\right]  \ \mid\ n-i\in X\right\}  $).

Forget that we fixed $p$. We thus have shown that $p\in Y$ for each
$p\in\operatorname*{rev}\nolimits_{n}X$. In other words, $\operatorname*{rev}%
\nolimits_{n}X\subseteq Y$.

Now, let $q\in Y$. Thus, $q\in Y=\left\{  i\in\left[  n-1\right]
\ \mid\ n-i\in X\right\}  $. In other words, $q$ is an $i\in\left[
n-1\right]  $ satisfying $n-i\in X$. In other words, $q\in\left[  n-1\right]
$ and $n-q\in X$. Furthermore, $q=n-\left(  n-q\right)  $. Hence, $q$ has the
form $n-x$ for some $x\in X$ (namely, for $x=n-q$). In other words,
$q\in\left\{  n-x\ \mid\ x\in X\right\}  $. This rewrites as $q\in
\operatorname*{rev}\nolimits_{n}X$ (since $\operatorname*{rev}\nolimits_{n}%
X=\left\{  n-x\ \mid\ x\in X\right\}  $).

Forget that we fixed $q$. We thus have shown that $q\in\operatorname*{rev}%
\nolimits_{n}X$ for each $q\in Y$. In other words, $Y\subseteq
\operatorname*{rev}\nolimits_{n}X$.

Combining this with $\operatorname*{rev}\nolimits_{n}X\subseteq Y$, we obtain
$\operatorname*{rev}\nolimits_{n}X=Y=\left\{  i\in\left[  n-1\right]
\ \mid\ n-i\in X\right\}  $. This proves Theorem \ref{thm.rev-rev}
\textbf{(e)}. \medskip

\textbf{(f)} Let $X$ and $Y$ be two subsets of $\left[  n-1\right]  $. Then,
$X\setminus Y$ is a subset of $\left[  n-1\right]  $ as well (since
$X\setminus Y\subseteq X\subseteq\left[  n-1\right]  $). Thus,
$\operatorname*{rev}\nolimits_{n}\left(  X\setminus Y\right)  \subseteq\left[
n-1\right]  $ (by Theorem \ref{thm.rev-rev} \textbf{(a)}, applied to
$X\setminus Y$ instead of $X$). Also, $\left(  \operatorname*{rev}%
\nolimits_{n}X\right)  \setminus\left(  \operatorname*{rev}\nolimits_{n}%
Y\right)  \subseteq\operatorname*{rev}\nolimits_{n}X\subseteq\left[
n-1\right]  $ (by Theorem \ref{thm.rev-rev} \textbf{(a)}).

Theorem \ref{thm.rev-rev} \textbf{(e)} yields
\[
\operatorname*{rev}\nolimits_{n}X=\left\{  i\in\left[  n-1\right]
\ \mid\ n-i\in X\right\}  .
\]
Hence, for any $i\in\left[  n-1\right]  $, we have the logical equivalence%
\begin{equation}
\left(  i\in\operatorname*{rev}\nolimits_{n}X\right)  \Longleftrightarrow
\left(  n-i\in X\right)  . \label{pf.thm.rev-rev.f.1}%
\end{equation}
The same argument (applied to $Y$ instead of $X$) shows that for any
$i\in\left[  n-1\right]  $, we have the logical equivalence%
\begin{equation}
\left(  i\in\operatorname*{rev}\nolimits_{n}Y\right)  \Longleftrightarrow
\left(  n-i\in Y\right)  . \label{pf.thm.rev-rev.f.2}%
\end{equation}
The same argument (applied to $X\setminus Y$ instead of $Y$) shows that for
any $i\in\left[  n-1\right]  $, we have the logical equivalence
\begin{equation}
\left(  i\in\operatorname*{rev}\nolimits_{n}\left(  X\setminus Y\right)
\right)  \Longleftrightarrow\left(  n-i\in X\setminus Y\right)  .
\label{pf.thm.rev-rev.f.3}%
\end{equation}


Now, for each $i\in\left[  n-1\right]  $, we have the following chain of
logical equivalences:%
\begin{align}
\left(  i\in\operatorname*{rev}\nolimits_{n}\left(  X\setminus Y\right)
\right)  \  &  \Longleftrightarrow\ \left(  n-i\in X\setminus Y\right)
\ \ \ \ \ \ \ \ \ \ \left(  \text{by (\ref{pf.thm.rev-rev.f.3})}\right)
\nonumber\\
&  \Longleftrightarrow\ \left(  n-i\in X\text{ and }n-i\notin Y\right)
\nonumber\\
&  \Longleftrightarrow\ \left(  \underbrace{n-i\in X}%
_{\substack{\Longleftrightarrow\ \left(  i\in\operatorname*{rev}%
\nolimits_{n}X\right)  \\\text{(by (\ref{pf.thm.rev-rev.f.1}))}}}\text{ but
not }\underbrace{n-i\in Y}_{\substack{\Longleftrightarrow\ \left(
i\in\operatorname*{rev}\nolimits_{n}Y\right)  \\\text{(by
(\ref{pf.thm.rev-rev.f.2}))}}}\right) \nonumber\\
&  \Longleftrightarrow\ \left(  i\in\operatorname*{rev}\nolimits_{n}X\text{
but not }i\in\operatorname*{rev}\nolimits_{n}Y\right) \nonumber\\
&  \Longleftrightarrow\ \left(  i\in\operatorname*{rev}\nolimits_{n}X\text{
and }i\notin\operatorname*{rev}\nolimits_{n}Y\right) \nonumber\\
&  \Longleftrightarrow\ \left(  i\in\left(  \operatorname*{rev}\nolimits_{n}%
X\right)  \setminus\left(  \operatorname*{rev}\nolimits_{n}Y\right)  \right)
. \label{pf.thm.rev-rev.f.4}%
\end{align}


Now, from $\operatorname*{rev}\nolimits_{n}\left(  X\setminus Y\right)
\subseteq\left[  n-1\right]  $, we obtain%
\begin{align}
\operatorname*{rev}\nolimits_{n}\left(  X\setminus Y\right)   &  =\left[
n-1\right]  \cap\left(  \operatorname*{rev}\nolimits_{n}\left(  X\setminus
Y\right)  \right) \nonumber\\
&  =\left\{  i\in\left[  n-1\right]  \ \mid\ \underbrace{i\in
\operatorname*{rev}\nolimits_{n}\left(  X\setminus Y\right)  }%
_{\substack{\Longleftrightarrow\ \left(  i\in\left(  \operatorname*{rev}%
\nolimits_{n}X\right)  \setminus\left(  \operatorname*{rev}\nolimits_{n}%
Y\right)  \right)  \\\text{(by (\ref{pf.thm.rev-rev.f.4}))}}}\right\}
\nonumber\\
&  =\left\{  i\in\left[  n-1\right]  \ \mid\ i\in\left(  \operatorname*{rev}%
\nolimits_{n}X\right)  \setminus\left(  \operatorname*{rev}\nolimits_{n}%
Y\right)  \right\}  . \label{pf.thm.rev-rev.f.5}%
\end{align}


However, from $\left(  \operatorname*{rev}\nolimits_{n}X\right)
\setminus\left(  \operatorname*{rev}\nolimits_{n}Y\right)  \subseteq\left[
n-1\right]  $, we obtain%
\begin{align*}
\left(  \operatorname*{rev}\nolimits_{n}X\right)  \setminus\left(
\operatorname*{rev}\nolimits_{n}Y\right)   &  =\left[  n-1\right]  \cap\left(
\left(  \operatorname*{rev}\nolimits_{n}X\right)  \setminus\left(
\operatorname*{rev}\nolimits_{n}Y\right)  \right) \\
&  =\left\{  i\in\left[  n-1\right]  \ \mid\ i\in\left(  \operatorname*{rev}%
\nolimits_{n}X\right)  \setminus\left(  \operatorname*{rev}\nolimits_{n}%
Y\right)  \right\}  .
\end{align*}
Comparing this with (\ref{pf.thm.rev-rev.f.5}), we find $\operatorname*{rev}%
\nolimits_{n}\left(  X\setminus Y\right)  =\left(  \operatorname*{rev}%
\nolimits_{n}X\right)  \setminus\left(  \operatorname*{rev}\nolimits_{n}%
Y\right)  $. This proves Theorem \ref{thm.rev-rev} \textbf{(f)}. \medskip

\textbf{(g)} Recall that
\begin{equation}
A\setminus\left(  A\setminus B\right)  =A\cap B \label{pf.thm.rev-rev.g.AB}%
\end{equation}
for any two sets $A$ and $B$.

Let $X$ and $Y$ be two subsets of $\left[  n-1\right]  $. Then, $X\setminus Y$
is a subset of $\left[  n-1\right]  $ as well (since $X\setminus Y\subseteq
X\subseteq\left[  n-1\right]  $). Hence, Theorem \ref{thm.rev-rev}
\textbf{(f)} (applied to $X\setminus Y$ instead of $Y$) yields%
\begin{align}
\operatorname*{rev}\nolimits_{n}\left(  X\setminus\left(  X\setminus Y\right)
\right)   &  =\left(  \operatorname*{rev}\nolimits_{n}X\right)  \setminus
\underbrace{\left(  \operatorname*{rev}\nolimits_{n}\left(  X\setminus
Y\right)  \right)  }_{\substack{=\left(  \operatorname*{rev}\nolimits_{n}%
X\right)  \setminus\left(  \operatorname*{rev}\nolimits_{n}Y\right)
\\\text{(by Theorem \ref{thm.rev-rev} \textbf{(f)})}}}\nonumber\\
&  =\left(  \operatorname*{rev}\nolimits_{n}X\right)  \setminus\left(  \left(
\operatorname*{rev}\nolimits_{n}X\right)  \setminus\left(  \operatorname*{rev}%
\nolimits_{n}Y\right)  \right) \nonumber\\
&  =\left(  \operatorname*{rev}\nolimits_{n}X\right)  \cap\left(
\operatorname*{rev}\nolimits_{n}Y\right)  \label{pf.thm.rev-rev.g.1}%
\end{align}
(by (\ref{pf.thm.rev-rev.g.AB}), applied to $A=\operatorname*{rev}%
\nolimits_{n}X$ and $B=\operatorname*{rev}\nolimits_{n}Y$). However,
$X\setminus\left(  X\setminus Y\right)  =X\cap Y$ (by
(\ref{pf.thm.rev-rev.g.AB}), applied to $A=X$ and $B=Y$). Thus, we can rewrite
(\ref{pf.thm.rev-rev.g.1}) as $\operatorname*{rev}\nolimits_{n}\left(  X\cap
Y\right)  =\left(  \operatorname*{rev}\nolimits_{n}X\right)  \cap\left(
\operatorname*{rev}\nolimits_{n}Y\right)  $. This proves Theorem
\ref{thm.rev-rev} \textbf{(g)}. \medskip

\textbf{(h)} The definition of $\operatorname*{rev}\nolimits_{n}\left(
\left[  n-1\right]  \right)  $ yields%
\begin{align*}
\operatorname*{rev}\nolimits_{n}\left(  \left[  n-1\right]  \right)   &
=\left\{  n-x\ \mid\ x\in\left[  n-1\right]  \right\} \\
&  =\left\{  n-x\ \mid\ x\in\left\{  1,2,\ldots,n-1\right\}  \right\}
\ \ \ \ \ \ \ \ \ \ \left(  \text{since }\left[  n-1\right]  =\left\{
1,2,\ldots,n-1\right\}  \right) \\
&  =\left\{  n-1,\ n-2,\ \ldots,\ n-\left(  n-1\right)  \right\} \\
&  =\left\{  n-1,\ n-2,\ \ldots,\ 1\right\} \\
&  =\left\{  1,2,\ldots,n-1\right\}  =\left[  n-1\right]  .
\end{align*}
This proves Theorem \ref{thm.rev-rev} \textbf{(h)}. \medskip

\textbf{(i)} Let $\alpha\in\operatorname*{Comp}\nolimits_{n}$ be a
composition. Write this composition $\alpha$ in the form $\alpha=\left(
\alpha_{1},\alpha_{2},\ldots,\alpha_{k}\right)  $. Then, $\operatorname*{rev}%
\alpha=\left(  \alpha_{k},\alpha_{k-1},\ldots,\alpha_{1}\right)  $ (by the
definition of $\operatorname*{rev}\alpha$). Also, the definition of
$\left\vert \alpha\right\vert $ yields $\left\vert \alpha\right\vert
=\alpha_{1}+\alpha_{2}+\cdots+\alpha_{k}$.

From $\alpha\in\operatorname*{Comp}\nolimits_{n}$, we see that $\alpha$ is a
composition of $n$ (since $\operatorname*{Comp}\nolimits_{n}$ is the set of
all compositions of $n$). In other words, $\alpha$ is a composition having
size $n$. Therefore, $\left\vert \alpha\right\vert =n$ (since $\left\vert
\alpha\right\vert $ is the size of $\alpha$, but we know that $\alpha$ has
size $n$).

For each $i\in\left\{  0,1,\ldots,k\right\}  $, we define two numbers%
\begin{align*}
u_{i}  &  :=\alpha_{1}+\alpha_{2}+\cdots+\alpha_{i}%
\ \ \ \ \ \ \ \ \ \ \text{and}\\
v_{i}  &  :=\alpha_{i+1}+\alpha_{i+2}+\cdots+\alpha_{k}.
\end{align*}
Each $i\in\left\{  0,1,\ldots,k\right\}  $ satisfies%
\begin{align*}
&  \underbrace{u_{i}}_{=\alpha_{1}+\alpha_{2}+\cdots+\alpha_{i}}%
+\underbrace{v_{i}}_{=\alpha_{i+1}+\alpha_{i+2}+\cdots+\alpha_{k}}\\
&  =\left(  \alpha_{1}+\alpha_{2}+\cdots+\alpha_{i}\right)  +\left(
\alpha_{i+1}+\alpha_{i+2}+\cdots+\alpha_{k}\right) \\
&  =\alpha_{1}+\alpha_{2}+\cdots+\alpha_{k}=\left\vert \alpha\right\vert
\ \ \ \ \ \ \ \ \ \ \left(  \text{since }\left\vert \alpha\right\vert
=\alpha_{1}+\alpha_{2}+\cdots+\alpha_{k}\right) \\
&  =n
\end{align*}
and therefore%
\begin{equation}
v_{i}=n-u_{i}. \label{pf.thm.rev-rev.i.1}%
\end{equation}


From $\alpha=\left(  \alpha_{1},\alpha_{2},\ldots,\alpha_{k}\right)  $, we
obtain
\begin{align*}
D\left(  \alpha\right)   &  =D\left(  \alpha_{1},\alpha_{2},\ldots,\alpha
_{k}\right)  =\left\{  \underbrace{\alpha_{1}+\alpha_{2}+\cdots+\alpha_{i}%
}_{\substack{=u_{i}\\\text{(since }u_{i}\text{ is defined}\\\text{to be
}\alpha_{1}+\alpha_{2}+\cdots+\alpha_{i}\text{)}}}\ \mid\ i\in\left[
k-1\right]  \right\} \\
&  \ \ \ \ \ \ \ \ \ \ \ \ \ \ \ \ \ \ \ \ \left(  \text{by the definition of
}D\left(  \alpha_{1},\alpha_{2},\ldots,\alpha_{k}\right)  \right) \\
&  =\left\{  u_{i}\ \mid\ i\in\left[  k-1\right]  \right\}  =\left\{
u_{1},u_{2},\ldots,u_{k-1}\right\}  .
\end{align*}
The definition of $\operatorname*{rev}\nolimits_{n}\left(  D\left(
\alpha\right)  \right)  $ yields%
\begin{align*}
\operatorname*{rev}\nolimits_{n}\left(  D\left(  \alpha\right)  \right)   &
=\left\{  n-x\ \mid\ x\in D\left(  \alpha\right)  \right\} \\
&  =\left\{  n-x\ \mid\ x\in\left\{  u_{1},u_{2},\ldots,u_{k-1}\right\}
\right\}  \ \ \ \ \ \ \ \ \ \ \left(  \text{since }D\left(  \alpha\right)
=\left\{  u_{1},u_{2},\ldots,u_{k-1}\right\}  \right) \\
&  =\left\{  n-u_{1},\ n-u_{2},\ \ldots,\ n-u_{k-1}\right\} \\
&  =\left\{  \underbrace{n-u_{i}}_{\substack{=v_{i}\\\text{(by
(\ref{pf.thm.rev-rev.i.1}))}}}\ \mid\ i\in\left[  k-1\right]  \right\}
=\left\{  v_{i}\ \mid\ i\in\left[  k-1\right]  \right\} \\
&  =\left\{  v_{1},v_{2},\ldots,v_{k-1}\right\}  .
\end{align*}
Comparing this with%
\begin{align*}
D\left(  \operatorname*{rev}\alpha\right)   &  =D\left(  \alpha_{k}%
,\alpha_{k-1},\ldots,\alpha_{1}\right)  \ \ \ \ \ \ \ \ \ \ \left(
\text{since }\operatorname*{rev}\alpha=\left(  \alpha_{k},\alpha_{k-1}%
,\ldots,\alpha_{1}\right)  \right) \\
&  =\left\{  \underbrace{\alpha_{k}+\alpha_{k-1}+\cdots+\alpha_{k-i+1}%
}_{\substack{=\alpha_{k-i+1}+\alpha_{k-i+2}+\cdots+\alpha_{k}\\=v_{k-i}%
\\\text{(since }v_{k-i}\text{ is defined}\\\text{to be }\alpha_{k-i+1}%
+\alpha_{k-i+2}+\cdots+\alpha_{k}\text{)}}}\ \mid\ i\in\left[  k-1\right]
\right\} \\
&  \ \ \ \ \ \ \ \ \ \ \ \ \ \ \ \ \ \ \ \ \left(  \text{by the definition of
}D\left(  \alpha_{k},\alpha_{k-1},\ldots,\alpha_{1}\right)  \right) \\
&  =\left\{  v_{k-i}\ \mid\ i\in\left[  k-1\right]  \right\}  =\left\{
v_{k-1},v_{k-2},\ldots,v_{k-\left(  k-1\right)  }\right\} \\
&  =\left\{  v_{k-1},v_{k-2},\ldots,v_{1}\right\}  =\left\{  v_{1}%
,v_{2},\ldots,v_{k-1}\right\}  ,
\end{align*}
we obtain $D\left(  \operatorname*{rev}\alpha\right)  =\operatorname*{rev}%
\nolimits_{n}\left(  D\left(  \alpha\right)  \right)  $. This proves Theorem
\ref{thm.rev-rev} \textbf{(i)}.
\end{proof}

\begin{corollary}
\label{cor.revrev}Let $n\in\mathbb{N}$, and let $\alpha\in\operatorname*{Comp}%
\nolimits_{n}$. Then, $\operatorname*{rev}\nolimits_{n}\left(  D\left(
\operatorname*{rev}\alpha\right)  \right)  =D\left(  \alpha\right)  $.
\end{corollary}

\begin{proof}
[Proof of Corollary \ref{cor.revrev}.]We have $\alpha\in\operatorname*{Comp}%
\nolimits_{n}$. In other words, $\alpha$ is a composition of $n$. That is,
$\alpha$ is a composition having size $n$. In other words, $\alpha
\in\operatorname*{Comp}$ and $\left\vert \alpha\right\vert =n$. Hence,
Proposition \ref{prop.rev-comp-size} yields $\left\vert \operatorname*{rev}%
\alpha\right\vert =\left\vert \alpha\right\vert =n$. In other words, the
composition $\operatorname*{rev}\alpha$ has size $n$. In other words,
$\operatorname*{rev}\alpha$ is a composition of $n$. In other words,
$\operatorname*{rev}\alpha\in\operatorname*{Comp}\nolimits_{n}$. Hence,
$D\left(  \operatorname*{rev}\alpha\right)  \in\mathcal{P}\left(  \left[
n-1\right]  \right)  $ (since $D$ is a map $\operatorname*{Comp}%
\nolimits_{n}\rightarrow\mathcal{P}\left(  \left[  n-1\right]  \right)  $). In
other words, $D\left(  \operatorname*{rev}\alpha\right)  $ is a subset of
$\left[  n-1\right]  $. Hence, $\operatorname*{rev}\nolimits_{n}\left(
D\left(  \operatorname*{rev}\alpha\right)  \right)  $ is well-defined.

Furthermore, $D\left(  \alpha\right)  \in\mathcal{P}\left(  \left[
n-1\right]  \right)  $ (since $D$ is a map $\operatorname*{Comp}%
\nolimits_{n}\rightarrow\mathcal{P}\left(  \left[  n-1\right]  \right)  $). In
other words, $D\left(  \alpha\right)  $ is a subset of $\left[  n-1\right]  $.

Theorem \ref{thm.rev-rev} \textbf{(i)} yields $D\left(  \operatorname*{rev}%
\alpha\right)  =\operatorname*{rev}\nolimits_{n}\left(  D\left(
\alpha\right)  \right)  $. Thus,%
\[
\operatorname*{rev}\nolimits_{n}\left(  \underbrace{D\left(
\operatorname*{rev}\alpha\right)  }_{=\operatorname*{rev}\nolimits_{n}\left(
D\left(  \alpha\right)  \right)  }\right)  =\operatorname*{rev}\nolimits_{n}%
\left(  \operatorname*{rev}\nolimits_{n}\left(  D\left(  \alpha\right)
\right)  \right)  =D\left(  \alpha\right)
\]
(by Theorem \ref{thm.rev-rev} \textbf{(b)}, applied to $X=D\left(
\alpha\right)  $). This proves Corollary \ref{cor.revrev}.
\end{proof}

\begin{corollary}
\label{cor.rev-bijn}Let $n\in\mathbb{N}$. Then, the map%
\begin{align*}
\operatorname*{Comp}\nolimits_{n} &  \rightarrow\operatorname*{Comp}%
\nolimits_{n},\\
\delta &  \mapsto\operatorname*{rev}\delta
\end{align*}
is a bijection.
\end{corollary}

\begin{proof}
[Proof of Corollary \ref{cor.rev-bijn}.]Each $\delta\in\operatorname*{Comp}%
\nolimits_{n}$ satisfies $\operatorname*{rev}\delta\in\operatorname*{Comp}%
\nolimits_{n}$\ \ \ \ \footnote{\textit{Proof.} Let $\delta\in
\operatorname*{Comp}\nolimits_{n}$. Thus, $\delta$ is a composition of $n$. In
other words, $\delta$ is a composition that has size $n$. In other words,
$\delta$ is a composition and satisfies $\left\vert \delta\right\vert =n$.
Now, Proposition \ref{prop.rev-comp-size} (applied to $\alpha=\delta$) yields
$\left\vert \operatorname*{rev}\delta\right\vert =\left\vert \delta\right\vert
=n$. Hence, $\operatorname*{rev}\delta$ is a composition that has size $n$
(since it has size $\left\vert \operatorname*{rev}\delta\right\vert =n$). In
other words, $\operatorname*{rev}\delta$ is a composition of $n$. In other
words, $\operatorname*{rev}\delta\in\operatorname*{Comp}\nolimits_{n}$. Qed.}.
Hence, the map%
\begin{align*}
\operatorname*{Comp}\nolimits_{n} &  \rightarrow\operatorname*{Comp}%
\nolimits_{n},\\
\delta &  \mapsto\operatorname*{rev}\delta
\end{align*}
is well-defined. It remains to prove that this map is a bijection.

Let us denote this map by $\operatorname*{rev}$ (since the image of any
$\delta\in\operatorname*{Comp}$ under this map is already being called
$\operatorname*{rev}\delta$). Thus, we must prove that this map
$\operatorname*{rev}$ is a bijection.

But this is easy: Every $\alpha\in\operatorname*{Comp}\nolimits_{n}$
satisfies
\begin{align*}
\left(  \operatorname*{rev}\circ\operatorname*{rev}\right)  \left(
\alpha\right)   &  =\operatorname*{rev}\left(  \operatorname*{rev}%
\alpha\right)  =\alpha\ \ \ \ \ \ \ \ \ \ \left(  \text{by Proposition
\ref{prop.rev-rev}}\right)  \\
&  =\operatorname*{id}\left(  \alpha\right)  .
\end{align*}
Thus, $\operatorname*{rev}\circ\operatorname*{rev}=\operatorname*{id}$. Hence,
the map $\operatorname*{rev}$ is inverse to itself. Thus, the map
$\operatorname*{rev}$ is invertible, i.e., bijective. In other words, it is a
bijection. This proves Corollary \ref{cor.rev-bijn}.
\end{proof}

\begin{proposition}
\label{prop.rev.D-inc}Let $n\in\mathbb{N}$. Let $\alpha\in\operatorname*{Comp}%
\nolimits_{n}$ and $\beta\in\operatorname*{Comp}\nolimits_{n}$ be arbitrary.
Then, we have the logical equivalence%
\[
\left(  D\left(  \operatorname*{rev}\beta\right)  \subseteq D\left(
\operatorname*{rev}\alpha\right)  \right)  \ \Longleftrightarrow\ \left(
D\left(  \beta\right)  \subseteq D\left(  \alpha\right)  \right)  .
\]

\end{proposition}

\begin{proof}
[Proof of Proposition \ref{prop.rev.D-inc}.] We have $\alpha\in
\operatorname*{Comp}\nolimits_{n}$ and thus $D\left(  \alpha\right)
\in\mathcal{P}\left(  \left[  n-1\right]  \right)  $ (since
$D:\operatorname*{Comp}\nolimits_{n}\rightarrow\mathcal{P}\left(  \left[
n-1\right]  \right)  $ is a map). In other words, $D\left(  \alpha\right)  $
is a subset of $\left[  n-1\right]  $. Similarly, $D\left(  \beta\right)  $ is
a subset of $\left[  n-1\right]  $.

Theorem \ref{thm.rev-rev} \textbf{(i)} yields $D\left(  \operatorname*{rev}%
\alpha\right)  =\operatorname*{rev}\nolimits_{n}\left(  D\left(
\alpha\right)  \right)  $. Also, Theorem \ref{thm.rev-rev} \textbf{(i)}
(applied to $\beta$ instead of $\alpha$) yields $D\left(  \operatorname*{rev}%
\beta\right)  =\operatorname*{rev}\nolimits_{n}\left(  D\left(  \beta\right)
\right)  $.

Now, if $D\left(  \beta\right)  \subseteq D\left(  \alpha\right)  $, then
$\operatorname*{rev}\nolimits_{n}\left(  D\left(  \beta\right)  \right)
\subseteq\operatorname*{rev}\nolimits_{n}\left(  D\left(  \alpha\right)
\right)  $ (by Theorem \ref{thm.rev-rev} \textbf{(c)}, applied to $X=D\left(
\beta\right)  $ and $Y=D\left(  \alpha\right)  $) and therefore%
\[
D\left(  \operatorname*{rev}\beta\right)  =\operatorname*{rev}\nolimits_{n}%
\left(  D\left(  \beta\right)  \right)  \subseteq\operatorname*{rev}%
\nolimits_{n}\left(  D\left(  \alpha\right)  \right)  =D\left(
\operatorname*{rev}\alpha\right)
\]
(since $D\left(  \operatorname*{rev}\alpha\right)  =\operatorname*{rev}%
\nolimits_{n}\left(  D\left(  \alpha\right)  \right)  $). In other words, the
implication
\begin{equation}
\left(  D\left(  \beta\right)  \subseteq D\left(  \alpha\right)  \right)
\ \Longrightarrow\ \left(  D\left(  \operatorname*{rev}\beta\right)  \subseteq
D\left(  \operatorname*{rev}\alpha\right)  \right)
\label{pf.prop.rev.D-inc.back}%
\end{equation}
holds.

Proposition \ref{prop.rev-rev} yields $\operatorname*{rev}\left(
\operatorname*{rev}\alpha\right)  =\alpha$. Similarly, $\operatorname*{rev}%
\left(  \operatorname*{rev}\beta\right)  =\beta$.

However, Corollary \ref{cor.rev-bijn} says that the map%
\begin{align*}
\operatorname*{Comp}\nolimits_{n} &  \rightarrow\operatorname*{Comp}%
\nolimits_{n},\\
\delta &  \mapsto\operatorname*{rev}\delta
\end{align*}
is a bijection. Thus, in particular, this map is well-defined. In other words,
for any $\delta\in\operatorname*{Comp}\nolimits_{n}$, we have
$\operatorname*{rev}\delta\in\operatorname*{Comp}\nolimits_{n}$. Applying this
to $\delta=\alpha$, we obtain $\operatorname*{rev}\alpha\in
\operatorname*{Comp}\nolimits_{n}$ (since $\alpha\in\operatorname*{Comp}%
\nolimits_{n}$). Similarly, $\operatorname*{rev}\beta\in\operatorname*{Comp}%
\nolimits_{n}$. Thus, we can apply the implication
(\ref{pf.prop.rev.D-inc.back}) to $\operatorname*{rev}\alpha$ and
$\operatorname*{rev}\beta$ instead of $\alpha$ and $\beta$. Hence, we obtain
the implication%
\[
\left(  D\left(  \operatorname*{rev}\beta\right)  \subseteq D\left(
\operatorname*{rev}\alpha\right)  \right)  \ \Longrightarrow\ \left(  D\left(
\operatorname*{rev}\left(  \operatorname*{rev}\beta\right)  \right)  \subseteq
D\left(  \operatorname*{rev}\left(  \operatorname*{rev}\alpha\right)  \right)
\right)  .
\]
In view of $\operatorname*{rev}\left(  \operatorname*{rev}\alpha\right)
=\alpha$ and $\operatorname*{rev}\left(  \operatorname*{rev}\beta\right)
=\beta$, we can rewrite this as
\[
\left(  D\left(  \operatorname*{rev}\beta\right)  \subseteq D\left(
\operatorname*{rev}\alpha\right)  \right)  \ \Longrightarrow\ \left(  D\left(
\beta\right)  \subseteq D\left(  \alpha\right)  \right)  .
\]
Combining this implication with (\ref{pf.prop.rev.D-inc.back}), we obtain the
logical equivalence%
\[
\left(  D\left(  \operatorname*{rev}\beta\right)  \subseteq D\left(
\operatorname*{rev}\alpha\right)  \right)  \ \Longleftrightarrow\ \left(
D\left(  \beta\right)  \subseteq D\left(  \alpha\right)  \right)  .
\]
This proves Proposition \ref{prop.rev.D-inc}.
\end{proof}

\section{The omega operation}

\begin{proposition}
\label{prop.omega-wd}Let $n\in\mathbb{N}$. Let $\gamma\in\operatorname*{Comp}%
\nolimits_{n}$. Then, there exists a unique composition $\delta$ of $n$
satisfying
\[
D\left(  \delta\right)  =\left[  n-1\right]  \setminus D\left(
\operatorname*{rev}\gamma\right)  .
\]

\end{proposition}

\begin{proof}
[Proof of Proposition \ref{prop.omega-wd}.]The set $\left[  n-1\right]
\setminus D\left(  \operatorname*{rev}\gamma\right)  $ is clearly a subset of
$\left[  n-1\right]  $, and thus belongs to $\mathcal{P}\left(  \left[
n-1\right]  \right)  $. Hence, there exists a unique $\delta\in
\operatorname*{Comp}\nolimits_{n}$ satisfying $D\left(  \delta\right)
=\left[  n-1\right]  \setminus D\left(  \operatorname*{rev}\gamma\right)  $
(since the map $D:\operatorname*{Comp}\nolimits_{n}\rightarrow\mathcal{P}%
\left(  \left[  n-1\right]  \right)  $ is a bijection). In other words, there
exists a unique composition $\delta$ of $n$ satisfying $D\left(
\delta\right)  =\left[  n-1\right]  \setminus D\left(  \operatorname*{rev}%
\gamma\right)  $ (because a composition $\delta$ of $n$ is the same as an
element $\delta\in\operatorname*{Comp}\nolimits_{n}$). This proves Proposition
\ref{prop.omega-wd}.
\end{proof}

We now define another operation on compositions:

\begin{definition}
\label{def.omega-on-comps}Let $n\in\mathbb{N}$. For any composition $\gamma
\in\operatorname*{Comp}\nolimits_{n}$, we let $\omega\left(  \gamma\right)  $
denote the unique composition $\delta$ of $n$ satisfying%
\[
D\left(  \delta\right)  =\left[  n-1\right]  \setminus D\left(
\operatorname*{rev}\gamma\right)  .
\]
(This $\omega\left(  \gamma\right)  $ is indeed well-defined, according to
Proposition \ref{prop.omega-wd}.)
\end{definition}

We observe the following simple properties of these compositions
$\omega\left(  \gamma\right)  $:

\begin{proposition}
\label{prop.omega-invol}Let $n\in\mathbb{N}$. Let $\gamma\in
\operatorname*{Comp}\nolimits_{n}$. Then:

\begin{enumerate}
\item[\textbf{(a)}] We have $\omega\left(  \gamma\right)  \in
\operatorname*{Comp}\nolimits_{n}$.

\item[\textbf{(b)}] We have $D\left(  \omega\left(  \gamma\right)  \right)
=\left[  n-1\right]  \setminus D\left(  \operatorname*{rev}\gamma\right)  $.

\item[\textbf{(c)}] We have $D\left(  \omega\left(  \gamma\right)  \right)
=\left[  n-1\right]  \setminus\operatorname*{rev}\nolimits_{n}\left(  D\left(
\gamma\right)  \right)  $.

\item[\textbf{(d)}] We have $\omega\left(  \omega\left(  \gamma\right)
\right)  =\gamma$.
\end{enumerate}
\end{proposition}

\begin{proof}
[Proof of Proposition \ref{prop.omega-invol}.]We have defined $\omega\left(
\gamma\right)  $ to be the unique composition $\delta$ of $n$ satisfying
$D\left(  \delta\right)  =\left[  n-1\right]  \setminus D\left(
\operatorname*{rev}\gamma\right)  $. Thus, $\omega\left(  \gamma\right)  $ is
a composition of $n$ and satisfies $D\left(  \omega\left(  \gamma\right)
\right)  =\left[  n-1\right]  \setminus D\left(  \operatorname*{rev}%
\gamma\right)  $. This proves Proposition \ref{prop.omega-invol} \textbf{(b)}.
Moreover, we have $\omega\left(  \gamma\right)  \in\operatorname*{Comp}%
\nolimits_{n}$ (since $\omega\left(  \gamma\right)  $ is a composition of
$n$); this proves Proposition \ref{prop.omega-invol} \textbf{(a)}.

It remains to prove parts \textbf{(c)} and \textbf{(d)}. \medskip

\textbf{(c)} Theorem \ref{thm.rev-rev} \textbf{(i)} (applied to $\alpha
=\gamma$) yields $D\left(  \operatorname*{rev}\gamma\right)
=\operatorname*{rev}\nolimits_{n}\left(  D\left(  \gamma\right)  \right)  $.
Now,%
\[
D\left(  \omega\left(  \gamma\right)  \right)  =\left[  n-1\right]
\setminus\underbrace{D\left(  \operatorname*{rev}\gamma\right)  }%
_{=\operatorname*{rev}\nolimits_{n}\left(  D\left(  \gamma\right)  \right)
}=\left[  n-1\right]  \setminus\operatorname*{rev}\nolimits_{n}\left(
D\left(  \gamma\right)  \right)  .
\]
This proves Proposition \ref{prop.omega-invol} \textbf{(c)}. \medskip

\textbf{(d)} We observe that $\gamma$ is a composition of $n$ (since
$\gamma\in\operatorname*{Comp}\nolimits_{n}$). In other words, $\gamma$ is a
composition having size $n$. In other words, $\gamma\in\operatorname*{Comp}$
and $\left\vert \gamma\right\vert =n$. However, Proposition
\ref{prop.rev-comp-size} (applied to $\alpha=\gamma$) yields $\left\vert
\operatorname*{rev}\gamma\right\vert =\left\vert \gamma\right\vert =n$. Hence,
$\operatorname*{rev}\gamma$ is a composition having size $\left\vert
\operatorname*{rev}\gamma\right\vert =n$. In other words, $\operatorname*{rev}%
\gamma$ is a composition of $n$. Hence, $\operatorname*{rev}\gamma
\in\operatorname*{Comp}\nolimits_{n}$. Thus, $D\left(  \operatorname*{rev}%
\gamma\right)  \in\mathcal{P}\left(  \left[  n-1\right]  \right)  $ (since $D$
is a map $\operatorname*{Comp}\nolimits_{n}\rightarrow\mathcal{P}\left(
\left[  n-1\right]  \right)  $). In other words, $D\left(  \operatorname*{rev}%
\gamma\right)  $ is a subset of $\left[  n-1\right]  $. Furthermore, $D\left(
\gamma\right)  \in\mathcal{P}\left(  \left[  n-1\right]  \right)  $ (since
$\gamma\in\operatorname*{Comp}\nolimits_{n}$ and since $D$ is a map
$\operatorname*{Comp}\nolimits_{n}\rightarrow\mathcal{P}\left(  \left[
n-1\right]  \right)  $). In other words, $D\left(  \gamma\right)  $ is a
subset of $\left[  n-1\right]  $. Also, $\left[  n-1\right]  $ is a subset of
$\left[  n-1\right]  $ as well.

Theorem \ref{thm.rev-rev} \textbf{(i)} (applied to $\alpha=\omega\left(
\gamma\right)  $) yields%
\begin{align*}
&  D\left(  \operatorname*{rev}\left(  \omega\left(  \gamma\right)  \right)
\right) \\
&  =\operatorname*{rev}\nolimits_{n}\left(  \underbrace{D\left(  \omega\left(
\gamma\right)  \right)  }_{=\left[  n-1\right]  \setminus D\left(
\operatorname*{rev}\gamma\right)  }\right) \\
&  =\operatorname*{rev}\nolimits_{n}\left(  \left[  n-1\right]  \setminus
D\left(  \operatorname*{rev}\gamma\right)  \right) \\
&  =\underbrace{\operatorname*{rev}\nolimits_{n}\left(  \left[  n-1\right]
\right)  }_{\substack{=\left[  n-1\right]  \\\text{(by Theorem
\ref{thm.rev-rev} \textbf{(h)})}}}\setminus\underbrace{\operatorname*{rev}%
\nolimits_{n}\left(  D\left(  \operatorname*{rev}\gamma\right)  \right)
}_{\substack{=D\left(  \gamma\right)  \\\text{(by Corollary \ref{cor.revrev}%
,}\\\text{applied to }\alpha=\gamma\text{)}}}\\
&  \ \ \ \ \ \ \ \ \ \ \ \ \ \ \ \ \ \ \ \ \left(  \text{by Theorem
\ref{thm.rev-rev} \textbf{(f)}, applied to }X=\left[  n-1\right]  \text{ and
}Y=D\left(  \operatorname*{rev}\gamma\right)  \right) \\
&  =\left[  n-1\right]  \setminus D\left(  \gamma\right)  .
\end{align*}


We have $\omega\left(  \omega\left(  \gamma\right)  \right)  \in
\operatorname*{Comp}\nolimits_{n}$ (by Proposition \ref{prop.omega-invol}
\textbf{(a)}, applied to $\omega\left(  \gamma\right)  $ instead of $\gamma$).
Moreover, Proposition \ref{prop.omega-invol} \textbf{(b)} (applied to
$\omega\left(  \gamma\right)  $ instead of $\gamma$) yields%
\begin{align*}
D\left(  \omega\left(  \omega\left(  \gamma\right)  \right)  \right)   &
=\left[  n-1\right]  \setminus\underbrace{D\left(  \operatorname*{rev}\left(
\omega\left(  \gamma\right)  \right)  \right)  }_{=\left[  n-1\right]
\setminus D\left(  \gamma\right)  }\\
&  =\left[  n-1\right]  \setminus\left(  \left[  n-1\right]  \setminus
D\left(  \gamma\right)  \right) \\
&  =\left[  n-1\right]  \cap D\left(  \gamma\right)
\ \ \ \ \ \ \ \ \ \ \left(
\begin{array}
[c]{c}%
\text{since }X\setminus\left(  X\setminus Y\right)  =X\cap Y\text{ for}\\
\text{any two sets }X\text{ and }Y
\end{array}
\right) \\
&  =D\left(  \gamma\right)  \ \ \ \ \ \ \ \ \ \ \left(  \text{since }D\left(
\gamma\right)  \text{ is a subset of }\left[  n-1\right]  \right)  .
\end{align*}


Recall that the map $D:\operatorname*{Comp}\nolimits_{n}\rightarrow
\mathcal{P}\left(  \left[  n-1\right]  \right)  $ is a bijection. Hence, this
map is bijective, therefore injective. Thus, any $\alpha,\beta\in
\operatorname*{Comp}\nolimits_{n}$ satisfying $D\left(  \alpha\right)
=D\left(  \beta\right)  $ must satisfy $\alpha=\beta$. We can apply this to
$\alpha=\omega\left(  \omega\left(  \gamma\right)  \right)  $ and
$\beta=\gamma$ (since $\gamma\in\operatorname*{Comp}\nolimits_{n}$ and
$\omega\left(  \omega\left(  \gamma\right)  \right)  \in\operatorname*{Comp}%
\nolimits_{n}$ and $D\left(  \omega\left(  \omega\left(  \gamma\right)
\right)  \right)  =D\left(  \gamma\right)  $), and thus we obtain
$\omega\left(  \omega\left(  \gamma\right)  \right)  =\gamma$. This proves
Proposition \ref{prop.omega-invol} \textbf{(d)}.
\end{proof}

\begin{proposition}
\label{prop.omega-obs1}Let $n$ be a positive integer. Let $\alpha
\in\operatorname*{Comp}\nolimits_{n}$ and $\gamma\in\operatorname*{Comp}%
\nolimits_{n}$. Then:

\begin{enumerate}
\item[\textbf{(a)}] We have
\[
\left\vert D\left(  \omega\left(  \gamma\right)  \right)  \cap D\left(
\alpha\right)  \right\vert =\ell\left(  \alpha\right)  -1-\left\vert D\left(
\gamma\right)  \cap D\left(  \operatorname*{rev}\alpha\right)  \right\vert .
\]


\item[\textbf{(b)}] We have
\[
\left\vert D\left(  \omega\left(  \gamma\right)  \right)  \setminus D\left(
\alpha\right)  \right\vert =n-\ell\left(  \alpha\right)  -\left\vert D\left(
\gamma\right)  \setminus D\left(  \operatorname*{rev}\alpha\right)
\right\vert .
\]

\end{enumerate}
\end{proposition}

\begin{proof}
[Proof of Proposition \ref{prop.omega-obs1}.]We have $\alpha\in
\operatorname*{Comp}\nolimits_{n}$. In other words, $\alpha$ is a composition
of $n$. That is, $\alpha$ is a composition having size $n$. In other words,
$\alpha\in\operatorname*{Comp}$ and $\left\vert \alpha\right\vert =n$. The
same argument (applied to $\gamma$ instead of $\alpha$) yields $\gamma
\in\operatorname*{Comp}$ and $\left\vert \gamma\right\vert =n$.

We have $n\geq1$ (since $n$ is a positive integer) and thus $n-1\in\mathbb{N}%
$. Hence, $\left\vert \left[  n-1\right]  \right\vert =n-1$.

Also, we have $\left\vert \alpha\right\vert =n\geq1>0$. Hence, Proposition
\ref{prop.D-size.pos} yields%
\begin{equation}
\left\vert D\left(  \alpha\right)  \right\vert =\ell\left(  \alpha\right)  -1.
\label{pf.prop.eta.S.Da}%
\end{equation}
Moreover, $D\left(  \alpha\right)  \in\mathcal{P}\left(  \left[  n-1\right]
\right)  $ (since $D$ is a map $\operatorname*{Comp}\nolimits_{n}%
\rightarrow\mathcal{P}\left(  \left[  n-1\right]  \right)  $); in other words,
$D\left(  \alpha\right)  $ is a subset of $\left[  n-1\right]  $. The same
argument (applied to $\gamma$ instead of $\alpha$) shows that $D\left(
\gamma\right)  $ is a subset of $\left[  n-1\right]  $. That is, we have
$D\left(  \gamma\right)  \subseteq\left[  n-1\right]  $. Hence,
$\operatorname*{rev}\nolimits_{n}\left(  D\left(  \gamma\right)  \right)
\subseteq\left[  n-1\right]  $ as well (by Theorem \ref{thm.rev-rev}
\textbf{(a)}, applied to $X=D\left(  \gamma\right)  $).

Also, $D\left(  \operatorname*{rev}\alpha\right)  $ is a subset of $\left[
n-1\right]  $ (this can be easily proved in the same way as in the proof of
Corollary \ref{cor.revrev} above).

Proposition \ref{prop.omega-invol} \textbf{(c)} yields $D\left(  \omega\left(
\gamma\right)  \right)  =\left[  n-1\right]  \setminus\operatorname*{rev}%
\nolimits_{n}\left(  D\left(  \gamma\right)  \right)  $.

However, for any three sets $X$, $Y$ and $Z$, we have $\left(  X\setminus
Y\right)  \cap Z=\left(  X\cap Z\right)  \setminus Y$. Applying this to
$X=\left[  n-1\right]  $ and $Y=\operatorname*{rev}\nolimits_{n}\left(
D\left(  \gamma\right)  \right)  $ and $Z=D\left(  \alpha\right)  $, we obtain%
\begin{align*}
\left(  \left[  n-1\right]  \setminus\operatorname*{rev}\nolimits_{n}\left(
D\left(  \gamma\right)  \right)  \right)  \cap D\left(  \alpha\right)   &
=\underbrace{\left(  \left[  n-1\right]  \cap D\left(  \alpha\right)  \right)
}_{\substack{=D\left(  \alpha\right)  \\\text{(since }D\left(  \alpha\right)
\text{ is a}\\\text{subset of }\left[  n-1\right]  \text{)}}}\setminus
\operatorname*{rev}\nolimits_{n}\left(  D\left(  \gamma\right)  \right) \\
&  =D\left(  \alpha\right)  \setminus\operatorname*{rev}\nolimits_{n}\left(
D\left(  \gamma\right)  \right)  .
\end{align*}
Thus,%
\begin{align*}
\underbrace{D\left(  \omega\left(  \gamma\right)  \right)  }_{=\left[
n-1\right]  \setminus\operatorname*{rev}\nolimits_{n}\left(  D\left(
\gamma\right)  \right)  }\cap\,D\left(  \alpha\right)   &  =\left(  \left[
n-1\right]  \setminus\operatorname*{rev}\nolimits_{n}\left(  D\left(
\gamma\right)  \right)  \right)  \cap D\left(  \alpha\right) \\
&  =D\left(  \alpha\right)  \setminus\operatorname*{rev}\nolimits_{n}\left(
D\left(  \gamma\right)  \right)  .
\end{align*}
Therefore,%
\begin{align}
\left\vert D\left(  \omega\left(  \gamma\right)  \right)  \cap D\left(
\alpha\right)  \right\vert  &  =\left\vert D\left(  \alpha\right)
\setminus\operatorname*{rev}\nolimits_{n}\left(  D\left(  \gamma\right)
\right)  \right\vert \nonumber\\
&  =\left\vert D\left(  \alpha\right)  \right\vert -\left\vert D\left(
\alpha\right)  \cap\operatorname*{rev}\nolimits_{n}\left(  D\left(
\gamma\right)  \right)  \right\vert \label{pf.prop.eta.S.O1.3}%
\end{align}
(since any finite sets $X$ and $Y$ satisfy $\left\vert X\setminus Y\right\vert
=\left\vert X\right\vert -\left\vert X\cap Y\right\vert $).

However, Theorem \ref{thm.rev-rev} \textbf{(g)} (applied to $X=D\left(
\gamma\right)  $ and $Y=D\left(  \operatorname*{rev}\alpha\right)  $) yields%
\begin{align}
\operatorname*{rev}\nolimits_{n}\left(  D\left(  \gamma\right)  \cap D\left(
\operatorname*{rev}\alpha\right)  \right)   &  =\operatorname*{rev}%
\nolimits_{n}\left(  D\left(  \gamma\right)  \right)  \cap
\underbrace{\operatorname*{rev}\nolimits_{n}\left(  D\left(
\operatorname*{rev}\alpha\right)  \right)  }_{\substack{=D\left(
\alpha\right)  \\\text{(by Corollary \ref{cor.revrev})}}}=\operatorname*{rev}%
\nolimits_{n}\left(  D\left(  \gamma\right)  \right)  \cap D\left(
\alpha\right) \nonumber\\
&  =D\left(  \alpha\right)  \cap\operatorname*{rev}\nolimits_{n}\left(
D\left(  \gamma\right)  \right)  . \label{pf.prop.omega-obs1.4}%
\end{align}
However, $D\left(  \gamma\right)  \cap D\left(  \operatorname*{rev}%
\alpha\right)  $ is a subset of $\left[  n-1\right]  $ (since $D\left(
\gamma\right)  \cap D\left(  \operatorname*{rev}\alpha\right)  \subseteq
D\left(  \gamma\right)  \subseteq\left[  n-1\right]  $). Hence, Theorem
\ref{thm.rev-rev} \textbf{(d)} (applied to $X=D\left(  \gamma\right)  \cap
D\left(  \operatorname*{rev}\alpha\right)  $) yields%
\[
\left\vert \operatorname*{rev}\nolimits_{n}\left(  D\left(  \gamma\right)
\cap D\left(  \operatorname*{rev}\alpha\right)  \right)  \right\vert
=\left\vert D\left(  \gamma\right)  \cap D\left(  \operatorname*{rev}%
\alpha\right)  \right\vert .
\]
In view of (\ref{pf.prop.omega-obs1.4}), we can rewrite this as%
\[
\left\vert D\left(  \alpha\right)  \cap\operatorname*{rev}\nolimits_{n}\left(
D\left(  \gamma\right)  \right)  \right\vert =\left\vert D\left(
\gamma\right)  \cap D\left(  \operatorname*{rev}\alpha\right)  \right\vert .
\]
Therefore, (\ref{pf.prop.eta.S.O1.3}) becomes%
\begin{align*}
\left\vert D\left(  \omega\left(  \gamma\right)  \right)  \cap D\left(
\alpha\right)  \right\vert  &  =\underbrace{\left\vert D\left(  \alpha\right)
\right\vert }_{\substack{=\ell\left(  \alpha\right)  -1\\\text{(by
(\ref{pf.prop.eta.S.Da}))}}}-\underbrace{\left\vert D\left(  \alpha\right)
\cap\operatorname*{rev}\nolimits_{n}\left(  D\left(  \gamma\right)  \right)
\right\vert }_{=\left\vert D\left(  \gamma\right)  \cap D\left(
\operatorname*{rev}\alpha\right)  \right\vert }\\
&  =\ell\left(  \alpha\right)  -1-\left\vert D\left(  \gamma\right)  \cap
D\left(  \operatorname*{rev}\alpha\right)  \right\vert .
\end{align*}
This proves Proposition \ref{prop.omega-obs1} \textbf{(a)}. \medskip

\textbf{(b)} From Proposition \ref{prop.omega-invol} \textbf{(c)}, we obtain
$D\left(  \omega\left(  \gamma\right)  \right)  =\left[  n-1\right]
\setminus\operatorname*{rev}\nolimits_{n}\left(  D\left(  \gamma\right)
\right)  $.

However, if two finite sets $X$ and $Y$ satisfy $Y\subseteq X$, then
$\left\vert X\setminus Y\right\vert =\left\vert X\right\vert -\left\vert
Y\right\vert $. Applying this to $X=\left[  n-1\right]  $ and
$Y=\operatorname*{rev}\nolimits_{n}\left(  D\left(  \gamma\right)  \right)  $,
we obtain
\begin{align*}
\left\vert \left[  n-1\right]  \setminus\operatorname*{rev}\nolimits_{n}%
\left(  D\left(  \gamma\right)  \right)  \right\vert  &
=\underbrace{\left\vert \left[  n-1\right]  \right\vert }_{\substack{=n-1}%
}-\underbrace{\left\vert \operatorname*{rev}\nolimits_{n}\left(  D\left(
\gamma\right)  \right)  \right\vert }_{\substack{=\left\vert D\left(
\gamma\right)  \right\vert \\\text{(by Theorem \ref{thm.rev-rev}
\textbf{(d)},}\\\text{applied to }X=D\left(  \gamma\right)  \text{)}%
}}\ \ \ \ \ \ \ \ \ \ \left(  \text{since }\operatorname*{rev}\nolimits_{n}%
\left(  D\left(  \gamma\right)  \right)  \subseteq\left[  n-1\right]  \right)
\\
&  =n-1-\left\vert D\left(  \gamma\right)  \right\vert .
\end{align*}
In view of $D\left(  \omega\left(  \gamma\right)  \right)  =\left[
n-1\right]  \setminus\operatorname*{rev}\nolimits_{n}\left(  D\left(
\gamma\right)  \right)  $, we can rewrite this as%
\begin{equation}
\left\vert D\left(  \omega\left(  \gamma\right)  \right)  \right\vert
=n-1-\left\vert D\left(  \gamma\right)  \right\vert .
\label{pf.prop.eta.S.O1.Dwgsize}%
\end{equation}


Next, recall that $\left\vert X\setminus Y\right\vert =\left\vert X\right\vert
-\left\vert X\cap Y\right\vert $ for any two finite sets $X$ and $Y$. From
this equality, we obtain%
\begin{equation}
\left\vert D\left(  \omega\left(  \gamma\right)  \right)  \setminus D\left(
\alpha\right)  \right\vert =\left\vert D\left(  \omega\left(  \gamma\right)
\right)  \right\vert -\left\vert D\left(  \omega\left(  \gamma\right)
\right)  \cap D\left(  \alpha\right)  \right\vert \label{pf.prop.eta.S.O1.b1}%
\end{equation}
and%
\begin{equation}
\left\vert D\left(  \gamma\right)  \setminus D\left(  \operatorname*{rev}%
\alpha\right)  \right\vert =\left\vert D\left(  \gamma\right)  \right\vert
-\left\vert D\left(  \gamma\right)  \cap D\left(  \operatorname*{rev}%
\alpha\right)  \right\vert . \label{pf.prop.eta.S.O1.b2}%
\end{equation}
Adding these two equalities together, we find%
\begin{align*}
&  \left\vert D\left(  \omega\left(  \gamma\right)  \right)  \setminus
D\left(  \alpha\right)  \right\vert +\left\vert D\left(  \gamma\right)
\setminus D\left(  \operatorname*{rev}\alpha\right)  \right\vert \\
&  =\underbrace{\left\vert D\left(  \omega\left(  \gamma\right)  \right)
\right\vert }_{\substack{=n-1-\left\vert D\left(  \gamma\right)  \right\vert
\\\text{(by (\ref{pf.prop.eta.S.O1.Dwgsize}))}}}-\underbrace{\left\vert
D\left(  \omega\left(  \gamma\right)  \right)  \cap D\left(  \alpha\right)
\right\vert }_{\substack{=\ell\left(  \alpha\right)  -1-\left\vert D\left(
\gamma\right)  \cap D\left(  \operatorname*{rev}\alpha\right)  \right\vert
\\\text{(by Proposition \ref{prop.omega-obs1} \textbf{(a)})}}}+\left\vert
D\left(  \gamma\right)  \right\vert -\left\vert D\left(  \gamma\right)  \cap
D\left(  \operatorname*{rev}\alpha\right)  \right\vert \\
&  =n-1-\left\vert D\left(  \gamma\right)  \right\vert -\left(  \ell\left(
\alpha\right)  -1-\left\vert D\left(  \gamma\right)  \cap D\left(
\operatorname*{rev}\alpha\right)  \right\vert \right)  +\left\vert D\left(
\gamma\right)  \right\vert -\left\vert D\left(  \gamma\right)  \cap D\left(
\operatorname*{rev}\alpha\right)  \right\vert \\
&  =n-\ell\left(  \alpha\right)  .
\end{align*}
In other words,%
\[
\left\vert D\left(  \omega\left(  \gamma\right)  \right)  \setminus D\left(
\alpha\right)  \right\vert =n-\ell\left(  \alpha\right)  -\left\vert D\left(
\gamma\right)  \setminus D\left(  \operatorname*{rev}\alpha\right)
\right\vert .
\]
This proves Proposition \ref{prop.omega-obs1} \textbf{(b)}.
\end{proof}

\section{Concatenation}

\subsection{Definition and basic properties}

The simplest binary operation on compositions is concatenation:

\begin{definition}
\label{def.concat}The \emph{concatenation} of two compositions $\beta=\left(
\beta_{1},\beta_{2},\ldots,\beta_{p}\right)  $ and $\gamma=\left(  \gamma
_{1},\gamma_{2},\ldots,\gamma_{q}\right)  $ is defined to be the composition
\[
\left(  \beta_{1},\beta_{2},\ldots,\beta_{p},\gamma_{1},\gamma_{2}%
,\ldots,\gamma_{q}\right)  .
\]
It is denoted by $\beta\gamma$.
\end{definition}

It is clear that any composition $\alpha$ satisfies $\alpha\varnothing
=\varnothing\alpha=\alpha$ (where $\varnothing$ denotes the empty composition,
as before). The next fact is also evident:

\begin{proposition}
\label{prop.concat.lsize}Let $\beta$ and $\gamma$ be two compositions. Then:

\begin{enumerate}
\item[\textbf{(a)}] We have $\ell\left(  \beta\gamma\right)  =\ell\left(
\beta\right)  +\ell\left(  \gamma\right)  $.

\item[\textbf{(b)}] We have $\left\vert \beta\gamma\right\vert =\left\vert
\beta\right\vert +\left\vert \gamma\right\vert $.
\end{enumerate}
\end{proposition}

\begin{proof}
[Proof of Proposition \ref{prop.concat.lsize}.]Write the compositions $\beta$
and $\gamma$ in the forms $\beta=\left(  \beta_{1},\beta_{2},\ldots,\beta
_{p}\right)  $ and $\gamma=\left(  \gamma_{1},\gamma_{2},\ldots,\gamma
_{q}\right)  $. Thus, the definition of $\beta\gamma$ yields%
\[
\beta\gamma=\left(  \beta_{1},\beta_{2},\ldots,\beta_{p},\gamma_{1},\gamma
_{2},\ldots,\gamma_{q}\right)  .
\]
Hence, the definition of $\ell\left(  \beta\gamma\right)  $ yields
$\ell\left(  \beta\gamma\right)  =p+q$, whereas the definition of $\left\vert
\beta\gamma\right\vert $ yields
\begin{equation}
\left\vert \beta\gamma\right\vert =\beta_{1}+\beta_{2}+\cdots+\beta_{p}%
+\gamma_{1}+\gamma_{2}+\cdots+\gamma_{q}. \label{pf.prop.concat.lsize.1}%
\end{equation}


However, from $\beta=\left(  \beta_{1},\beta_{2},\ldots,\beta_{p}\right)  $,
we obtain $\ell\left(  \beta\right)  =p$ and $\left\vert \beta\right\vert
=\beta_{1}+\beta_{2}+\cdots+\beta_{p}$. Moreover, from $\gamma=\left(
\gamma_{1},\gamma_{2},\ldots,\gamma_{q}\right)  $, we obtain $\ell\left(
\gamma\right)  =q$ and $\left\vert \gamma\right\vert =\gamma_{1}+\gamma
_{2}+\cdots+\gamma_{q}$. Thus,%
\[
\underbrace{\ell\left(  \beta\right)  }_{=p}+\underbrace{\ell\left(
\gamma\right)  }_{=q}=p+q=\ell\left(  \beta\gamma\right)
\ \ \ \ \ \ \ \ \ \ \left(  \text{since }\ell\left(  \beta\gamma\right)
=p+q\right)  .
\]
This proves Proposition \ref{prop.concat.lsize} \textbf{(a)}. \medskip

\textbf{(b)} Adding the equalities $\left\vert \beta\right\vert =\beta
_{1}+\beta_{2}+\cdots+\beta_{p}$ and $\left\vert \gamma\right\vert =\gamma
_{1}+\gamma_{2}+\cdots+\gamma_{q}$ together, we obtain
\[
\left\vert \beta\right\vert +\left\vert \gamma\right\vert =\beta_{1}+\beta
_{2}+\cdots+\beta_{p}+\gamma_{1}+\gamma_{2}+\cdots+\gamma_{q}=\left\vert
\beta\gamma\right\vert
\]
(by (\ref{pf.prop.concat.lsize.1})). This proves Proposition
\ref{prop.concat.lsize} \textbf{(b)}.
\end{proof}

\subsection{Concatenation and reversal}

Concatenation and reversal interact in a nice way:

\begin{proposition}
\label{prop.concat.rev}Let $\beta$ and $\gamma$ be two compositions. Then,
$\operatorname*{rev}\left(  \beta\gamma\right)  =\left(  \operatorname*{rev}%
\gamma\right)  \left(  \operatorname*{rev}\beta\right)  $.
\end{proposition}

\begin{proof}
[Proof of Proposition \ref{prop.concat.rev}.]Write the compositions $\beta$
and $\gamma$ in the forms $\beta=\left(  \beta_{1},\beta_{2},\ldots,\beta
_{p}\right)  $ and $\gamma=\left(  \gamma_{1},\gamma_{2},\ldots,\gamma
_{q}\right)  $. Thus, the definition of $\beta\gamma$ yields%
\[
\beta\gamma=\left(  \beta_{1},\beta_{2},\ldots,\beta_{p},\gamma_{1},\gamma
_{2},\ldots,\gamma_{q}\right)  .
\]
Hence, the definition of $\operatorname*{rev}\left(  \beta\gamma\right)  $
yields%
\begin{equation}
\operatorname*{rev}\left(  \beta\gamma\right)  =\left(  \gamma_{q}%
,\gamma_{q-1},\ldots,\gamma_{1},\beta_{p},\beta_{p-1},\ldots,\beta_{1}\right)
. \label{pf.prop.concat.rev.1}%
\end{equation}
However, the definition of $\operatorname*{rev}\beta$ yields
$\operatorname*{rev}\beta=\left(  \beta_{p},\beta_{p-1},\ldots,\beta
_{1}\right)  $ (since $\beta=\left(  \beta_{1},\beta_{2},\ldots,\beta
_{p}\right)  $). Furthermore, the definition of $\operatorname*{rev}\gamma$
yields $\operatorname*{rev}\gamma=\left(  \gamma_{q},\gamma_{q-1}%
,\ldots,\gamma_{1}\right)  $ (since $\gamma=\left(  \gamma_{1},\gamma
_{2},\ldots,\gamma_{q}\right)  $). Thus,%
\begin{align*}
&  \underbrace{\left(  \operatorname*{rev}\gamma\right)  }_{=\left(
\gamma_{q},\gamma_{q-1},\ldots,\gamma_{1}\right)  }\underbrace{\left(
\operatorname*{rev}\beta\right)  }_{=\left(  \beta_{p},\beta_{p-1}%
,\ldots,\beta_{1}\right)  }\\
&  =\left(  \gamma_{q},\gamma_{q-1},\ldots,\gamma_{1}\right)  \left(
\beta_{p},\beta_{p-1},\ldots,\beta_{1}\right) \\
&  =\left(  \gamma_{q},\gamma_{q-1},\ldots,\gamma_{1},\beta_{p},\beta
_{p-1},\ldots,\beta_{1}\right)
\end{align*}
(by the definition of concatenation). Comparing this with
(\ref{pf.prop.concat.rev.1}), we obtain $\operatorname*{rev}\left(
\beta\gamma\right)  =\left(  \operatorname*{rev}\gamma\right)  \left(
\operatorname*{rev}\beta\right)  $. This proves Proposition
\ref{prop.concat.rev}.
\end{proof}

\subsection{Concatenation and partial sums}

We shall next show some less trivial properties of concatenations of
compositions. We will need the following notation:

\begin{definition}
\label{def.K+m}If $K$ is a set of integers, and if $m$ is an integer, then we
define two sets $K+m$ and $K-m$ by%
\begin{align*}
K+m  &  :=\left\{  k+m\ \mid\ k\in K\right\}  ,\\
K-m  &  :=\left\{  k-m\ \mid\ k\in K\right\}  .
\end{align*}
Clearly, both of these sets $K+m$ and $K-m$ are again sets of integers.
\end{definition}

For example, $\left\{  2,3,5\right\}  +10=\left\{  12,13,15\right\}  $ and
$\left\{  2,3,5\right\}  -1=\left\{  1,2,4\right\}  $. Visually, you can think
of $K+m$ as being the set $K$, moved to the right by $m$ units on the number
line. Similarly, $K-m$ is the set $K$, moved to the left by $m$ units on the
number line.

Clearly, if $K$ is any set of integers, and if $m$ is an integer, then
$\left(  K+m\right)  -m=K$ and $\left(  K-m\right)  +m=K$.

Now, if we know the sizes and the partial sum sets of two compositions $\beta$
and $\gamma$, then we can compute the partial sum set of their concatenation
$\beta\gamma$ as follows:

\begin{proposition}
\label{prop.concat.D}Let $\beta$ and $\gamma$ be two compositions such that
$\beta\neq\varnothing$ and $\gamma\neq\varnothing$. Let $m=\left\vert
\beta\right\vert $. Then,
\[
D\left(  \beta\gamma\right)  =\left\{  m\right\}  \cup D\left(  \beta\right)
\cup\left(  D\left(  \gamma\right)  +m\right)  .
\]

\end{proposition}

\begin{proof}
[Proof of Proposition \ref{prop.concat.D}.]Write the compositions $\beta$ and
$\gamma$ in the forms $\beta=\left(  \beta_{1},\beta_{2},\ldots,\beta
_{p}\right)  $ and $\gamma=\left(  \gamma_{1},\gamma_{2},\ldots,\gamma
_{q}\right)  $. From $\beta\neq\varnothing$, we easily obtain $p\neq
0$\ \ \ \ \footnote{\textit{Proof.} If we had $p=0$, then we would have%
\begin{align*}
\beta &  =\left(  \beta_{1},\beta_{2},\ldots,\beta_{p}\right)  =\left(
\beta_{1},\beta_{2},\ldots,\beta_{0}\right)  \ \ \ \ \ \ \ \ \ \ \left(
\text{since }p=0\right) \\
&  =\left(  {}\right)  =\varnothing,
\end{align*}
which would contradict $\beta\neq\varnothing$. Hence, we cannot have $p=0$.
Thus, we have $p\neq0$.}. Similarly, from $\gamma\neq\varnothing$, we obtain
$q\neq0$. Also, $m=\left\vert \beta\right\vert =\beta_{1}+\beta_{2}%
+\cdots+\beta_{p}$ (by the definition of $\left\vert \beta\right\vert $, since
$\beta=\left(  \beta_{1},\beta_{2},\ldots,\beta_{p}\right)  $). Thus,
$\beta_{1}+\beta_{2}+\cdots+\beta_{p}=m$.

From $\beta=\left(  \beta_{1},\beta_{2},\ldots,\beta_{p}\right)  $, we obtain
\begin{equation}
D\left(  \beta\right)  =\left\{  \beta_{1},\ \ \beta_{1}+\beta_{2}%
,\ \ \beta_{1}+\beta_{2}+\beta_{3},\ \ \ldots,\ \ \beta_{1}+\beta_{2}%
+\cdots+\beta_{p-1}\right\}  \label{pf.prop.concat.D.Db=}%
\end{equation}
(by the definition of $D\left(  \beta\right)  $).

From $\gamma=\left(  \gamma_{1},\gamma_{2},\ldots,\gamma_{q}\right)  $, we
obtain%
\[
D\left(  \gamma\right)  =\left\{  \gamma_{1},\ \ \gamma_{1}+\gamma
_{2},\ \ \gamma_{1}+\gamma_{2}+\gamma_{3},\ \ \ldots,\ \ \gamma_{1}+\gamma
_{2}+\cdots+\gamma_{q-1}\right\}  .
\]
However, the definition of $D\left(  \gamma\right)  +m$ yields%
\begin{align}
&  D\left(  \gamma\right)  +m\nonumber\\
&  =\left\{  \underbrace{k+m}_{=m+k}\ \mid\ k\in D\left(  \gamma\right)
\right\}  =\left\{  m+k\ \mid\ k\in D\left(  \gamma\right)  \right\}
\nonumber\\
&  =\left\{  m+k\ \mid\ k\in\left\{  \gamma_{1},\ \ \gamma_{1}+\gamma
_{2},\ \ \gamma_{1}+\gamma_{2}+\gamma_{3},\ \ \ldots,\ \ \gamma_{1}+\gamma
_{2}+\cdots+\gamma_{q-1}\right\}  \right\} \nonumber\\
&  \ \ \ \ \ \ \ \ \ \ \left(  \text{since }D\left(  \gamma\right)  =\left\{
\gamma_{1},\ \ \gamma_{1}+\gamma_{2},\ \ \gamma_{1}+\gamma_{2}+\gamma
_{3},\ \ \ldots,\ \ \gamma_{1}+\gamma_{2}+\cdots+\gamma_{q-1}\right\}  \right)
\nonumber\\
&  =\left\{  m+\gamma_{1},\ \ m+\gamma_{1}+\gamma_{2},\ \ m+\gamma_{1}%
+\gamma_{2}+\gamma_{3},\ \ \ldots,\ \ m+\gamma_{1}+\gamma_{2}+\cdots
+\gamma_{q-1}\right\}  . \label{pf.prop.concat.D.Dgm=}%
\end{align}


Now, recall that $\beta=\left(  \beta_{1},\beta_{2},\ldots,\beta_{p}\right)  $
and $\gamma=\left(  \gamma_{1},\gamma_{2},\ldots,\gamma_{q}\right)  $. Hence,
the definition of $\beta\gamma$ yields%
\[
\beta\gamma=\left(  \beta_{1},\beta_{2},\ldots,\beta_{p},\gamma_{1},\gamma
_{2},\ldots,\gamma_{q}\right)  .
\]
Hence, the definition of $D\left(  \beta\gamma\right)  $ yields\footnote{Note
the tacit use of $p\neq0$ and $q\neq0$ in this computation.}%
\begin{align*}
&  D\left(  \beta\gamma\right) \\
=  &  \ \{\beta_{1},\ \ \beta_{1}+\beta_{2},\ \ \beta_{1}+\beta_{2}+\beta
_{3},\ \ \ldots,\ \ \beta_{1}+\beta_{2}+\cdots+\beta_{p-1},\\
&  \ \ \ \ \beta_{1}+\beta_{2}+\cdots+\beta_{p},\\
&  \ \ \ \ \beta_{1}+\beta_{2}+\cdots+\beta_{p}+\gamma_{1},\\
&  \ \ \ \ \beta_{1}+\beta_{2}+\cdots+\beta_{p}+\gamma_{1}+\gamma_{2},\\
&  \ \ \ \ \beta_{1}+\beta_{2}+\cdots+\beta_{p}+\gamma_{1}+\gamma_{2}%
+\gamma_{3},\\
&  \ \ \ \ \ldots,\\
&  \ \ \ \ \beta_{1}+\beta_{2}+\cdots+\beta_{p}+\gamma_{1}+\gamma_{2}%
+\cdots+\gamma_{q-1}\}\\
=  &  \ \{\beta_{1},\ \ \beta_{1}+\beta_{2},\ \ \beta_{1}+\beta_{2}+\beta
_{3},\ \ \ldots,\ \ \beta_{1}+\beta_{2}+\cdots+\beta_{p-1},\\
&  \ \ \ \ m,\ \ m+\gamma_{1},\ \ m+\gamma_{1}+\gamma_{2},\ \ m+\gamma
_{1}+\gamma_{2}+\gamma_{3},\ \ \ldots,\ \ m+\gamma_{1}+\gamma_{2}%
+\cdots+\gamma_{q-1}\}\\
&  \ \ \ \ \ \ \ \ \ \ \ \ \ \ \ \ \ \ \ \ \left(  \text{since }\beta
_{1}+\beta_{2}+\cdots+\beta_{p}=m\right) \\
=  &  \ \underbrace{\left\{  \beta_{1},\ \ \beta_{1}+\beta_{2},\ \ \beta
_{1}+\beta_{2}+\beta_{3},\ \ \ldots,\ \ \beta_{1}+\beta_{2}+\cdots+\beta
_{p-1}\right\}  }_{\substack{=D\left(  \beta\right)  \\\text{(by
(\ref{pf.prop.concat.D.Db=}))}}}\\
&  \ \ \ \ \ \ \ \ \ \ \cup\left\{  m\right\} \\
&  \ \ \ \ \ \ \ \ \ \ \cup\underbrace{\left\{  m+\gamma_{1},\ \ m+\gamma
_{1}+\gamma_{2},\ \ m+\gamma_{1}+\gamma_{2}+\gamma_{3},\ \ \ldots
,\ \ m+\gamma_{1}+\gamma_{2}+\cdots+\gamma_{q-1}\right\}  }%
_{\substack{=D\left(  \gamma\right)  +m\\\text{(by
(\ref{pf.prop.concat.D.Dgm=}))}}}\\
=  &  \ D\left(  \beta\right)  \cup\left\{  m\right\}  \cup\left(  D\left(
\gamma\right)  +\left\{  m\right\}  \right)  =\left\{  m\right\}  \cup
D\left(  \beta\right)  \cup\left(  D\left(  \gamma\right)  +m\right)  .
\end{align*}
This proves Proposition \ref{prop.concat.D}.
\end{proof}

The following is a variant of Proposition \ref{prop.concat.D} that avoids the
requirements that $\beta\neq\varnothing$ and $\gamma\neq\varnothing$:

\begin{proposition}
\label{prop.concat.D0}Let $\beta$ and $\gamma$ be two compositions. Let
$m=\left\vert \beta\right\vert $ and $n=\left\vert \gamma\right\vert $. Then,
\[
D\left(  \beta\gamma\right)  =\left(  \left\{  m\right\}  \cup D\left(
\beta\right)  \cup\left(  D\left(  \gamma\right)  +m\right)  \right)
\cap\left[  m+n-1\right]  .
\]

\end{proposition}

\begin{proof}
[Proof of Proposition \ref{prop.concat.D0}.]We know that $\gamma$ is a
composition having size $n$ (since the size of $\gamma$ is $\left\vert
\gamma\right\vert =n$). In other words, $\gamma$ is a composition of $n$. In
other words, $\gamma\in\operatorname*{Comp}\nolimits_{n}$ (since
$\operatorname*{Comp}\nolimits_{n}$ is the set of all compositions of $n$).

We know that $\beta$ is a composition having size $m$ (since the size of
$\beta$ is $\left\vert \beta\right\vert =m$). In other words, $\beta$ is a
composition of $m$. In other words, $\beta\in\operatorname*{Comp}%
\nolimits_{m}$ (since $\operatorname*{Comp}\nolimits_{m}$ is the set of all
compositions of $m$).

We have $0\notin\left[  n-1\right]  $ (since the set $\left[  n-1\right]
=\left\{  1,2,\ldots,n-1\right\}  $ does not contain $0$) and $m\notin\left[
m-1\right]  $ (since the set $\left[  m-1\right]  =\left\{  1,2,\ldots
,m-1\right\}  $ does not contain $m$).

We are in one of the following three cases:

\textit{Case 1:} We have $\beta=\varnothing$.

\textit{Case 2:} We have $\gamma=\varnothing$.

\textit{Case 3:} We have neither $\beta=\varnothing$ nor $\gamma=\varnothing$.

Let us first consider Case 1. In this case, we have $\beta=\varnothing$. Thus,
$D\left(  \beta\right)  =D\left(  \varnothing\right)  =\varnothing$ (by the
definition of the map $D:\operatorname*{Comp}\nolimits_{0}\rightarrow
\mathcal{P}\left(  \left[  0-1\right]  \right)  $).

Moreover, $m=\left\vert \beta\right\vert $. In view of $\beta=\varnothing$,
this rewrites as $m=\left\vert \varnothing\right\vert =0$. Thus, $D\left(
\gamma\right)  +\underbrace{m}_{=0}=D\left(  \gamma\right)  +0=D\left(
\gamma\right)  $ (because any set $K$ of integers satisfies $K+0=K$).

Recall that $D$ is a map $\operatorname*{Comp}\nolimits_{n}\rightarrow
\mathcal{P}\left(  \left[  n-1\right]  \right)  $. Hence, $D\left(
\gamma\right)  \in\mathcal{P}\left(  \left[  n-1\right]  \right)  $ (since
$\gamma\in\operatorname*{Comp}\nolimits_{n}$). In other words, $D\left(
\gamma\right)  \subseteq\left[  n-1\right]  $.

Now,%
\begin{align}
&  \left(  \left\{  \underbrace{m}_{=0}\right\}  \cup\underbrace{D\left(
\beta\right)  }_{=\varnothing}\cup\underbrace{\left(  D\left(  \gamma\right)
+m\right)  }_{=D\left(  \gamma\right)  }\right)  \cap\left[  \underbrace{m}%
_{=0}+n-1\right] \nonumber\\
&  =\left(  \underbrace{\left\{  0\right\}  \cup\varnothing}_{=\left\{
0\right\}  }\cup D\left(  \gamma\right)  \right)  \cap\left[
\underbrace{0+n-1}_{=n-1}\right] \nonumber\\
&  =\left(  \left\{  0\right\}  \cup D\left(  \gamma\right)  \right)
\cap\left[  n-1\right]  . \label{pf.prop.concat.D0.c1.1}%
\end{align}


However, recall that any three sets $X_{1},X_{2},Y$ satisfy%
\[
\left(  X_{1}\cup X_{2}\right)  \cap Y=\left(  X_{1}\cap Y\right)  \cup\left(
X_{2}\cap Y\right)  .
\]
Applying this to $X_{1}=\left\{  0\right\}  $, $X_{2}=D\left(  \gamma\right)
$ and $Y=\left[  n-1\right]  $, we obtain%
\begin{align*}
\left(  \left\{  0\right\}  \cup D\left(  \gamma\right)  \right)  \cap\left[
n-1\right]   &  =\underbrace{\left(  \left\{  0\right\}  \cap\left[
n-1\right]  \right)  }_{\substack{=\varnothing\\\text{(since }0\notin\left[
n-1\right]  \text{)}}}\cup\underbrace{\left(  D\left(  \gamma\right)
\cap\left[  n-1\right]  \right)  }_{\substack{=D\left(  \gamma\right)
\\\text{(since }D\left(  \gamma\right)  \subseteq\left[  n-1\right]  \text{)}%
}}\\
&  =\varnothing\cup D\left(  \gamma\right)  =D\left(  \gamma\right)  .
\end{align*}
Thus, (\ref{pf.prop.concat.D0.c1.1}) rewrites as%
\begin{align*}
&  \left(  \left\{  m\right\}  \cup D\left(  \beta\right)  \cup\left(
D\left(  \gamma\right)  +m\right)  \right)  \cap\left[  m+n-1\right] \\
&  =D\left(  \gamma\right)  =D\left(  \beta\gamma\right)
\ \ \ \ \ \ \ \ \ \ \left(  \text{since }\gamma=\beta\gamma\text{ (because
}\underbrace{\beta}_{=\varnothing}\gamma=\varnothing\gamma=\gamma
\text{)}\right)  .
\end{align*}
Hence, Proposition \ref{prop.concat.D0} is proved in Case 1.

Let us now consider Case 2. In this case, we have $\gamma=\varnothing$. Hence,
$D\left(  \gamma\right)  =D\left(  \varnothing\right)  =\varnothing$. Hence,
$\underbrace{D\left(  \gamma\right)  }_{=\varnothing}+m=\varnothing
+m=\varnothing$ (since $\varnothing+k=\varnothing$ for any integer $k$).

Moreover, $n=\left\vert \gamma\right\vert $. In view of $\gamma=\varnothing$,
this rewrites as $n=\left\vert \varnothing\right\vert =0$.

Recall that $D$ is a map $\operatorname*{Comp}\nolimits_{m}\rightarrow
\mathcal{P}\left(  \left[  m-1\right]  \right)  $. Hence, $D\left(
\beta\right)  \in\mathcal{P}\left(  \left[  m-1\right]  \right)  $ (since
$\beta\in\operatorname*{Comp}\nolimits_{m}$). In other words, $D\left(
\beta\right)  \subseteq\left[  m-1\right]  $.

Now,%
\begin{align}
&  \left(  \left\{  m\right\}  \cup D\left(  \beta\right)  \cup
\underbrace{\left(  D\left(  \gamma\right)  +m\right)  }_{=\varnothing
}\right)  \cap\left[  m+\underbrace{n}_{=0}-1\right] \nonumber\\
&  =\underbrace{\left(  \left\{  m\right\}  \cup D\left(  \beta\right)
\cup\varnothing\right)  }_{=\left\{  m\right\}  \cup D\left(  \beta\right)
}\cap\left[  \underbrace{m+0-1}_{=m-1}\right] \nonumber\\
&  =\left(  \left\{  m\right\}  \cup D\left(  \beta\right)  \right)
\cap\left[  m-1\right]  . \label{pf.prop.concat.D0.c2.1}%
\end{align}


However, recall that any three sets $X_{1},X_{2},Y$ satisfy%
\[
\left(  X_{1}\cup X_{2}\right)  \cap Y=\left(  X_{1}\cap Y\right)  \cup\left(
X_{2}\cap Y\right)  .
\]
Applying this to $X_{1}=\left\{  m\right\}  $, $X_{2}=D\left(  \beta\right)  $
and $Y=\left[  m-1\right]  $, we obtain%
\begin{align*}
\left(  \left\{  m\right\}  \cup D\left(  \beta\right)  \right)  \cap\left[
m-1\right]   &  =\underbrace{\left(  \left\{  m\right\}  \cap\left[
m-1\right]  \right)  }_{\substack{=\varnothing\\\text{(since }m\notin\left[
m-1\right]  \text{)}}}\cup\underbrace{\left(  D\left(  \beta\right)
\cap\left[  m-1\right]  \right)  }_{\substack{=D\left(  \beta\right)
\\\text{(since }D\left(  \beta\right)  \subseteq\left[  m-1\right]  \text{)}%
}}\\
&  =\varnothing\cup D\left(  \beta\right)  =D\left(  \beta\right)  .
\end{align*}
Thus, (\ref{pf.prop.concat.D0.c2.1}) rewrites as%
\begin{align*}
&  \left(  \left\{  m\right\}  \cup D\left(  \beta\right)  \cup\left(
D\left(  \gamma\right)  +m\right)  \right)  \cap\left[  m+n-1\right] \\
&  =D\left(  \beta\right)  =D\left(  \beta\gamma\right)
\ \ \ \ \ \ \ \ \ \ \left(  \text{since }\beta=\beta\gamma\text{ (because
}\beta\underbrace{\gamma}_{=\varnothing}=\beta\varnothing=\beta\text{)}%
\right)  .
\end{align*}
Hence, Proposition \ref{prop.concat.D0} is proved in Case 2.

Now, let us consider Case 3. In this case, we have neither $\beta=\varnothing$
nor $\gamma=\varnothing$. In other words, we have $\beta\neq\varnothing$ and
$\gamma\neq\varnothing$. Thus, Proposition \ref{prop.concat.D} yields%
\[
D\left(  \beta\gamma\right)  =\left\{  m\right\}  \cup D\left(  \beta\right)
\cup\left(  D\left(  \gamma\right)  +m\right)  .
\]
However, Proposition \ref{prop.concat.lsize} \textbf{(b)} yields $\left\vert
\beta\gamma\right\vert =\underbrace{\left\vert \beta\right\vert }%
_{=m}+\underbrace{\left\vert \gamma\right\vert }_{=n}=m+n$. Thus, the
composition $\beta\gamma$ has size $\left\vert \beta\gamma\right\vert =m+n$.
In other words, $\beta\gamma$ is a composition of $m+n$. In other words,
$\beta\gamma\in\operatorname*{Comp}\nolimits_{m+n}$. Hence, $D\left(
\beta\gamma\right)  \in\mathcal{P}\left(  \left[  m+n-1\right]  \right)  $
(since $D$ is a map $\operatorname*{Comp}\nolimits_{m+n}\rightarrow
\mathcal{P}\left(  \left[  m+n-1\right]  \right)  $). In other words,
$D\left(  \beta\gamma\right)  \subseteq\left[  m+n-1\right]  $. Hence,
$D\left(  \beta\gamma\right)  \cap\left[  m+n-1\right]  =D\left(  \beta
\gamma\right)  $, so that%
\begin{align*}
D\left(  \beta\gamma\right)   &  =\underbrace{D\left(  \beta\gamma\right)
}_{=\left\{  m\right\}  \cup D\left(  \beta\right)  \cup\left(  D\left(
\gamma\right)  +m\right)  }\cap\left[  m+n-1\right] \\
&  =\left(  \left\{  m\right\}  \cup D\left(  \beta\right)  \cup\left(
D\left(  \gamma\right)  +m\right)  \right)  \cap\left[  m+n-1\right]  .
\end{align*}
Therefore, Proposition \ref{prop.concat.D0} is proved in Case 3.

We have now proved Proposition \ref{prop.concat.D0} in each of the three Cases
1, 2 and 3. This completes the proof of Proposition \ref{prop.concat.D0}.
\end{proof}

Conversely, given two compositions $\beta$ and $\gamma$, we can reconstruct
the partial sum sets $D\left(  \beta\right)  $ and $D\left(  \gamma\right)  $
if we know the size $\left\vert \beta\right\vert $ and the partial sum set
$D\left(  \beta\gamma\right)  $ as follows:

\begin{proposition}
\label{prop.concat.Dcut}Let $\beta$ and $\gamma$ be two compositions. Let
$m=\left\vert \beta\right\vert $. Then:

\begin{enumerate}
\item[\textbf{(a)}] We have $D\left(  \beta\right)  =D\left(  \beta
\gamma\right)  \cap\left[  m-1\right]  $.

\item[\textbf{(b)}] We have $D\left(  \gamma\right)  =\left(  D\left(
\beta\gamma\right)  \setminus\left[  m\right]  \right)  -m$.
\end{enumerate}
\end{proposition}

\begin{proof}
[Proof of Proposition \ref{prop.concat.Dcut}.]Let $n=\left\vert \gamma
\right\vert $. Then, as in the above proof of Proposition \ref{prop.concat.D0}%
, we can show that $\beta\in\operatorname*{Comp}\nolimits_{m}$ and $\gamma
\in\operatorname*{Comp}\nolimits_{n}$.

Recall that $D$ is a map $\operatorname*{Comp}\nolimits_{n}\rightarrow
\mathcal{P}\left(  \left[  n-1\right]  \right)  $. Hence, $D\left(
\gamma\right)  \in\mathcal{P}\left(  \left[  n-1\right]  \right)  $ (since
$\gamma\in\operatorname*{Comp}\nolimits_{n}$). In other words, $D\left(
\gamma\right)  \subseteq\left[  n-1\right]  $. The same argument (applied to
$\beta$ and $m$ instead of $\gamma$ and $n$) yields $D\left(  \beta\right)
\subseteq\left[  m-1\right]  $. Also, note that $\left[  m-1\right]
\subseteq\left[  m+n-1\right]  $ (since $m-1\leq m+n-1$ (because $n\geq0$)).
\medskip

\textbf{(a)} Let $x\in D\left(  \beta\right)  $. We shall show that $x\in
D\left(  \beta\gamma\right)  \cap\left[  m-1\right]  $.

Indeed, we observe that%
\[
x\in D\left(  \beta\right)  \subseteq\left\{  m\right\}  \cup D\left(
\beta\right)  \cup\left(  D\left(  \gamma\right)  +m\right)  .
\]
Combining this with $x\in D\left(  \beta\right)  \subseteq\left[  m-1\right]
\subseteq\left[  m+n-1\right]  $, we obtain%
\begin{align*}
x  &  \in\left(  \left\{  m\right\}  \cup D\left(  \beta\right)  \cup\left(
D\left(  \gamma\right)  +m\right)  \right)  \cap\left[  m+n-1\right] \\
&  =D\left(  \beta\gamma\right)  \ \ \ \ \ \ \ \ \ \ \left(  \text{by
Proposition \ref{prop.concat.D0}}\right)  .
\end{align*}
Combining this with $x\in\left[  m-1\right]  $, we obtain $x\in D\left(
\beta\gamma\right)  \cap\left[  m-1\right]  $.

Forget that we fixed $x$. We thus have shown that $x\in D\left(  \beta
\gamma\right)  \cap\left[  m-1\right]  $ for each $x\in D\left(  \beta\right)
$. In other words,%
\begin{equation}
D\left(  \beta\right)  \subseteq D\left(  \beta\gamma\right)  \cap\left[
m-1\right]  . \label{pf.prop.concat.Dcut.a.1}%
\end{equation}


On the other hand, let $y\in D\left(  \beta\gamma\right)  \cap\left[
m-1\right]  $. Thus, $y\in D\left(  \beta\gamma\right)  $ and $y\in\left[
m-1\right]  $. From $y\in\left[  m-1\right]  =\left\{  1,2,\ldots,m-1\right\}
$, we obtain $y\leq m-1<m$. Thus, we cannot have $y\in\left\{  m\right\}  $
(because $y\in\left\{  m\right\}  $ would entail $y=m$, which would contradict
$y<m$). Furthermore, we cannot have $y\in D\left(  \gamma\right)  +m$ (because
$y\in D\left(  \gamma\right)  +m$ would entail that $y\geq m$%
\ \ \ \ \footnote{\textit{Proof.} Assume that $y\in D\left(  \gamma\right)
+m$. We must show that $y\geq m$.
\par
We have $y\in D\left(  \gamma\right)  +m=\left\{  k+m\ \mid\ k\in D\left(
\gamma\right)  \right\}  $ (by the definition of $D\left(  \gamma\right)
+m$). In other words, $y=k+m$ for some $k\in D\left(  \gamma\right)  $.
Consider this $k$. From $k\in D\left(  \gamma\right)  \subseteq\left[
n-1\right]  =\left\{  1,2,\ldots,n-1\right\}  $, we obtain $k\geq1>0$. Hence,
$y=\underbrace{k}_{>0}+m>m$, thus $y\geq m$.}, which would contradict $y<m$).

However,
\begin{align*}
y  &  \in D\left(  \beta\gamma\right) \\
&  =\left(  \left\{  m\right\}  \cup D\left(  \beta\right)  \cup\left(
D\left(  \gamma\right)  +m\right)  \right)  \cap\left[  m+n-1\right]
\ \ \ \ \ \ \ \ \ \ \left(  \text{by Proposition \ref{prop.concat.D0}}\right)
\\
&  \subseteq\left\{  m\right\}  \cup D\left(  \beta\right)  \cup\left(
D\left(  \gamma\right)  +m\right)  .
\end{align*}
In other words, we have $y\in\left\{  m\right\}  $ or $y\in D\left(
\beta\right)  $ or $y\in D\left(  \gamma\right)  +m$. Hence, we must have
$y\in D\left(  \beta\right)  $ (since we cannot have $y\in\left\{  m\right\}
$, and we cannot have $y\in D\left(  \gamma\right)  +m$).

Forget that we fixed $y$. We thus have shown that $y\in D\left(  \beta\right)
$ for each $y\in D\left(  \beta\gamma\right)  \cap\left[  m-1\right]  $. In
other words,%
\[
D\left(  \beta\gamma\right)  \cap\left[  m-1\right]  \subseteq D\left(
\beta\right)  .
\]
Combining this with (\ref{pf.prop.concat.Dcut.a.1}), we obtain $D\left(
\beta\right)  =D\left(  \beta\gamma\right)  \cap\left[  m-1\right]  $. This
proves Proposition \ref{prop.concat.Dcut} \textbf{(a)}. \medskip

\textbf{(b)} The definition of $D\left(  \gamma\right)  +m$ yields
\begin{equation}
D\left(  \gamma\right)  +m=\left\{  k+m\ \mid\ k\in D\left(  \gamma\right)
\right\}  . \label{pf.prop.concat.Dcut.b.1}%
\end{equation}
The definition of $\left(  D\left(  \beta\gamma\right)  \setminus\left[
m\right]  \right)  -m$ yields%
\begin{equation}
\left(  D\left(  \beta\gamma\right)  \setminus\left[  m\right]  \right)
-m=\left\{  k-m\ \mid\ k\in D\left(  \beta\gamma\right)  \setminus\left[
m\right]  \right\}  . \label{pf.prop.concat.Dcut.b.2}%
\end{equation}


Let $x\in D\left(  \gamma\right)  $. We shall show that $x\in\left(  D\left(
\beta\gamma\right)  \setminus\left[  m\right]  \right)  -m$.

Indeed, we have $x\in D\left(  \gamma\right)  \subseteq\left[  n-1\right]
=\left\{  1,2,\ldots,n-1\right\}  $, so that
\[
x+m\in\left\{  m+1,m+2,\ldots,m+n-1\right\}  \subseteq\left\{  1,2,\ldots
,m+n-1\right\}  =\left[  m+n-1\right]  .
\]
Also, from $x\in\left\{  1,2,\ldots,n-1\right\}  $, we obtain $x\geq1>0$, and
therefore $\underbrace{x}_{>0}+m>m$, so that $x+m\notin\left[  m\right]
$\ \ \ \ \footnote{\textit{Proof.} If we had $x+m\in\left[  m\right]  $, then
we would have $x+m\leq m$ (since $x+m\in\left[  m\right]  =\left\{
1,2,\ldots,m\right\}  $), but this would contradict $x+m>m$. Hence, we cannot
have $x+m\in\left[  m\right]  $. In other words, we have $x+m\notin\left[
m\right]  $.}.

Next, we recall that $x\in D\left(  \gamma\right)  $. Thus, the number $x+m$
can be written in the form $k+m$ for some $k\in D\left(  \gamma\right)  $
(namely, for $k=x$). In other words, $x+m\in\left\{  k+m\ \mid\ k\in D\left(
\gamma\right)  \right\}  $. In view of (\ref{pf.prop.concat.Dcut.b.1}), we can
rewrite this as $x+m\in D\left(  \gamma\right)  +m$. Hence,%
\[
x+m\in D\left(  \gamma\right)  +m\subseteq\left\{  m\right\}  \cup D\left(
\beta\right)  \cup\left(  D\left(  \gamma\right)  +m\right)  .
\]
Combining this with $x+m\in\left[  m+n-1\right]  $, we obtain%
\begin{align*}
x+m  &  \in\left(  \left\{  m\right\}  \cup D\left(  \beta\right)  \cup\left(
D\left(  \gamma\right)  +m\right)  \right)  \cap\left[  m+n-1\right] \\
&  =D\left(  \beta\gamma\right)  \ \ \ \ \ \ \ \ \ \ \left(  \text{by
Proposition \ref{prop.concat.D0}}\right)  .
\end{align*}
Combining this with $x+m\notin\left[  m\right]  $, we obtain $x+m\in D\left(
\beta\gamma\right)  \setminus\left[  m\right]  $. We also have $x=\left(
x+m\right)  -m$. Therefore, $x$ has the form $k-m$ for some $k\in D\left(
\beta\gamma\right)  \setminus\left[  m\right]  $ (namely, for $k=x+m$),
because $x+m\in D\left(  \beta\gamma\right)  \setminus\left[  m\right]  $. In
other words, $x\in\left\{  k-m\ \mid\ k\in D\left(  \beta\gamma\right)
\setminus\left[  m\right]  \right\}  $. In view of
(\ref{pf.prop.concat.Dcut.b.2}), this rewrites as $x\in\left(  D\left(
\beta\gamma\right)  \setminus\left[  m\right]  \right)  -m$.

Forget that we fixed $x$. We thus have shown that $x\in\left(  D\left(
\beta\gamma\right)  \setminus\left[  m\right]  \right)  -m$ for each $x\in
D\left(  \gamma\right)  $. In other words,%
\begin{equation}
D\left(  \gamma\right)  \subseteq\left(  D\left(  \beta\gamma\right)
\setminus\left[  m\right]  \right)  -m. \label{pf.prop.concat.Dcut.b.3}%
\end{equation}


On the other hand, let $y\in\left(  D\left(  \beta\gamma\right)
\setminus\left[  m\right]  \right)  -m$. Thus,
\[
y\in\left(  D\left(  \beta\gamma\right)  \setminus\left[  m\right]  \right)
-m=\left\{  k-m\ \mid\ k\in D\left(  \beta\gamma\right)  \setminus\left[
m\right]  \right\}
\]
(by (\ref{pf.prop.concat.Dcut.b.2})). In other words, $y=k-m$ for some $k\in
D\left(  \beta\gamma\right)  \setminus\left[  m\right]  $. Consider this $k$,
and denote it by $z$. Thus, $y=z-m$ and $z\in D\left(  \beta\gamma\right)
\setminus\left[  m\right]  $.

From $z\in D\left(  \beta\gamma\right)  \setminus\left[  m\right]  $, we
obtain $z\in D\left(  \beta\gamma\right)  $ and $z\notin\left[  m\right]  $.
In particular,%
\begin{align*}
z  &  \in D\left(  \beta\gamma\right) \\
&  =\left(  \left\{  m\right\}  \cup D\left(  \beta\right)  \cup\left(
D\left(  \gamma\right)  +m\right)  \right)  \cap\left[  m+n-1\right]
\ \ \ \ \ \ \ \ \ \ \left(  \text{by Proposition \ref{prop.concat.D0}}\right)
\\
&  \subseteq\left[  m+n-1\right]  =\left\{  1,2,\ldots,m+n-1\right\}  .
\end{align*}
Combining this with $z\notin\left[  m\right]  =\left\{  1,2,\ldots,m\right\}
$, we obtain%
\[
z\in\left\{  1,2,\ldots,m+n-1\right\}  \setminus\left\{  1,2,\ldots,m\right\}
=\left\{  m+1,m+2,\ldots,m+n-1\right\}  .
\]
Hence, $z\geq m+1>m$.

Furthermore,%
\begin{align*}
z  &  \in\left(  \left\{  m\right\}  \cup D\left(  \beta\right)  \cup\left(
D\left(  \gamma\right)  +m\right)  \right)  \cap\left[  m+n-1\right] \\
&  \subseteq\left\{  m\right\}  \cup D\left(  \beta\right)  \cup\left(
D\left(  \gamma\right)  +m\right)  .
\end{align*}
In other words, we have $z\in\left\{  m\right\}  $ or $z\in D\left(
\beta\right)  $ or $z\in D\left(  \gamma\right)  +m$. However, we cannot have
$z\in\left\{  m\right\}  $\ \ \ \ \footnote{\textit{Proof.} Assume the
contrary. Thus, $z\in\left\{  m\right\}  $. Hence, $z=m$, which contradicts
$z>m$. This contradiction shows that our assumption was wrong. Thus, we cannot
have $z\in\left\{  m\right\}  $.}, and we also cannot have $z\in D\left(
\beta\right)  $\ \ \ \ \footnote{\textit{Proof.} Assume the contrary. Thus,
$z\in D\left(  \beta\right)  \subseteq\left[  m-1\right]  =\left\{
1,2,\ldots,m-1\right\}  $. Hence, $z\leq m-1<m$, which contradicts $z>m$. This
contradiction shows that our assumption was wrong. Thus, we cannot have $z\in
D\left(  \beta\right)  $.}. Hence, we must have $z\in D\left(  \gamma\right)
+m$ (since we have $z\in\left\{  m\right\}  $ or $z\in D\left(  \beta\right)
$ or $z\in D\left(  \gamma\right)  +m$). In view of
(\ref{pf.prop.concat.Dcut.b.1}), this rewrites as
\[
z\in\left\{  k+m\ \mid\ k\in D\left(  \gamma\right)  \right\}  .
\]
In other words, $z=k+m$ for some $k\in D\left(  \gamma\right)  $. Consider
this $k$. We have $y=\underbrace{z}_{=k+m}-m=k+m-m=k\in D\left(
\gamma\right)  $.

Forget that we fixed $y$. We thus have shown that $y\in D\left(
\gamma\right)  $ for each $y\in\left(  D\left(  \beta\gamma\right)
\setminus\left[  m\right]  \right)  -m$. In other words,%
\[
\left(  D\left(  \beta\gamma\right)  \setminus\left[  m\right]  \right)
-m\subseteq D\left(  \gamma\right)  .
\]
Combining this with (\ref{pf.prop.concat.Dcut.b.3}), we obtain $\left(
D\left(  \beta\gamma\right)  \setminus\left[  m\right]  \right)  -m=D\left(
\gamma\right)  $. This proves Proposition \ref{prop.concat.Dcut} \textbf{(b)}.
\end{proof}

\subsection{Further lemmas}

The next few propositions and lemmas will be used in a later proof.

\begin{proposition}
\label{prop.concat.2}Let $\beta$, $\gamma$, $\beta^{\prime}$ and
$\gamma^{\prime}$ be four compositions such that $\left\vert \beta^{\prime
}\right\vert =\left\vert \beta\right\vert $ and $D\left(  \beta^{\prime
}\right)  \subseteq D\left(  \beta\right)  $ and $\left\vert \gamma^{\prime
}\right\vert =\left\vert \gamma\right\vert $ and $D\left(  \gamma^{\prime
}\right)  \subseteq D\left(  \gamma\right)  $. Then, $D\left(  \beta^{\prime
}\gamma^{\prime}\right)  \subseteq D\left(  \beta\gamma\right)  $.
\end{proposition}

\begin{proof}
[Proof of Proposition \ref{prop.concat.2}.]Let $m=\left\vert \beta\right\vert
$ and $n=\left\vert \gamma\right\vert $. Thus, $\left\vert \beta^{\prime
}\right\vert =\left\vert \beta\right\vert =m$ and $\left\vert \gamma^{\prime
}\right\vert =\left\vert \gamma\right\vert =n$.

It is easy to see that if $K$ and $L$ are two sets of integers satisfying
$K\subseteq L$, and if $k$ is any integer, then $K+k\subseteq L+k$. Applying
this to $K=D\left(  \gamma^{\prime}\right)  $ and $L=D\left(  \gamma\right)  $
and $k=m$, we obtain $D\left(  \gamma^{\prime}\right)  +m\subseteq D\left(
\gamma\right)  +m$ (since $D\left(  \gamma^{\prime}\right)  \subseteq D\left(
\gamma\right)  $).

Now, Proposition \ref{prop.concat.D0} yields%
\begin{equation}
D\left(  \beta\gamma\right)  =\left(  \left\{  m\right\}  \cup D\left(
\beta\right)  \cup\left(  D\left(  \gamma\right)  +m\right)  \right)
\cap\left[  m+n-1\right]  . \label{pf.prop.concat.2.1}%
\end{equation}


Also, we have $m=\left\vert \beta^{\prime}\right\vert $ (since $\left\vert
\beta^{\prime}\right\vert =m$) and $n=\left\vert \gamma^{\prime}\right\vert $
(since $\left\vert \gamma^{\prime}\right\vert =n$). Hence, Proposition
\ref{prop.concat.D0} (applied to $\beta^{\prime}$ and $\gamma^{\prime}$
instead of $\beta$ and $\gamma$) yields%
\begin{align*}
D\left(  \beta^{\prime}\gamma^{\prime}\right)   &  =\left(  \left\{
m\right\}  \cup\underbrace{D\left(  \beta^{\prime}\right)  }_{\subseteq
D\left(  \beta\right)  }\cup\underbrace{\left(  D\left(  \gamma^{\prime
}\right)  +m\right)  }_{\subseteq D\left(  \gamma\right)  +m}\right)
\cap\left[  m+n-1\right] \\
&  \subseteq\left(  \left\{  m\right\}  \cup D\left(  \beta\right)
\cup\left(  D\left(  \gamma\right)  +m\right)  \right)  \cap\left[
m+n-1\right]  =D\left(  \beta\gamma\right)
\end{align*}
(by (\ref{pf.prop.concat.2.1})). This proves Proposition \ref{prop.concat.2}.
\end{proof}

\begin{proposition}
\label{prop.concat.4u}Let $\alpha\in\operatorname*{Comp}$ be any composition,
and let $m\in\mathbb{N}$. Then, there exists at most one pair $\left(
\beta,\gamma\right)  $ of compositions such that $\left\vert \beta\right\vert
=m$ and $\beta\gamma=\alpha$.
\end{proposition}

\begin{proof}
[Proof of Proposition \ref{prop.concat.4u}.]Let $\left(  \beta^{\prime}%
,\gamma^{\prime}\right)  $ and $\left(  \beta^{\prime\prime},\gamma
^{\prime\prime}\right)  $ be two pairs $\left(  \beta,\gamma\right)  $ of
compositions such that $\left\vert \beta\right\vert =m$ and $\beta
\gamma=\alpha$. Thus, $\left(  \beta^{\prime},\gamma^{\prime}\right)  $ and
$\left(  \beta^{\prime\prime},\gamma^{\prime\prime}\right)  $ are two pairs of
compositions and have the property that $\left\vert \beta^{\prime}\right\vert
=m$ and $\beta^{\prime}\gamma^{\prime}=\alpha$ and $\left\vert \beta
^{\prime\prime}\right\vert =m$ and $\beta^{\prime\prime}\gamma^{\prime\prime
}=\alpha$. Thus, $\alpha=\beta^{\prime\prime}\gamma^{\prime\prime}$.

We have $m=\left\vert \beta^{\prime}\right\vert $ (since $\left\vert
\beta^{\prime}\right\vert =m$). Thus, Proposition \ref{prop.concat.Dcut}
\textbf{(a)} (applied to $\beta^{\prime}$ and $\gamma^{\prime}$ instead of
$\beta$ and $\gamma$) yields
\[
D\left(  \beta^{\prime}\right)  =D\left(  \underbrace{\beta^{\prime}%
\gamma^{\prime}}_{=\alpha}\right)  \cap\left[  m-1\right]  =D\left(
\alpha\right)  \cap\left[  m-1\right]  .
\]
The same argument (applied to $\beta^{\prime\prime}$ and $\gamma^{\prime
\prime}$ instead of $\beta^{\prime}$ and $\gamma^{\prime}$) yields
\[
D\left(  \beta^{\prime\prime}\right)  =D\left(  \alpha\right)  \cap\left[
m-1\right]  .
\]
Comparing these two equalities, we find $D\left(  \beta^{\prime}\right)
=D\left(  \beta^{\prime\prime}\right)  $.

Now, $\beta^{\prime}$ is a composition having size $\left\vert \beta^{\prime
}\right\vert =m$. In other words, $\beta^{\prime}$ is a composition of $m$. In
other words, $\beta^{\prime}\in\operatorname*{Comp}\nolimits_{m}$. Similarly,
$\beta^{\prime\prime}\in\operatorname*{Comp}\nolimits_{m}$.

Recall that the map $D:\operatorname*{Comp}\nolimits_{n}\rightarrow
\mathcal{P}\left(  \left[  n-1\right]  \right)  $ is a bijection. Similarly,
the map $D:\operatorname*{Comp}\nolimits_{m}\rightarrow\mathcal{P}\left(
\left[  m-1\right]  \right)  $ is a bijection. Hence, this map $D$ is
bijective, thus injective. In other words, if $\varphi,\psi\in
\operatorname*{Comp}\nolimits_{m}$ satisfy $D\left(  \varphi\right)  =D\left(
\psi\right)  $, then $\varphi=\psi$. Applying this to $\varphi=\beta^{\prime}$
and $\psi=\beta^{\prime\prime}$, we obtain $\beta^{\prime}=\beta^{\prime
\prime}$ (since $\beta^{\prime}\in\operatorname*{Comp}\nolimits_{m}$ and
$\beta^{\prime\prime}\in\operatorname*{Comp}\nolimits_{m}$ and $D\left(
\beta^{\prime}\right)  =D\left(  \beta^{\prime\prime}\right)  $).

Furthermore, Proposition \ref{prop.concat.Dcut} \textbf{(b)} (applied to
$\beta^{\prime}$ and $\gamma^{\prime}$ instead of $\beta$ and $\gamma$)
yields
\[
D\left(  \gamma^{\prime}\right)  =\left(  D\left(  \underbrace{\beta^{\prime
}\gamma^{\prime}}_{=\alpha}\right)  \setminus\left[  m\right]  \right)
-m=\left(  D\left(  \alpha\right)  \setminus\left[  m\right]  \right)  -m.
\]
The same argument (applied to $\beta^{\prime\prime}$ and $\gamma^{\prime
\prime}$ instead of $\beta^{\prime}$ and $\gamma^{\prime}$) yields
\[
D\left(  \gamma^{\prime\prime}\right)  =\left(  D\left(  \alpha\right)
\setminus\left[  m\right]  \right)  -m.
\]
Comparing these two equalities, we find $D\left(  \gamma^{\prime}\right)
=D\left(  \gamma^{\prime\prime}\right)  $.

Set $n=\left\vert \gamma^{\prime}\right\vert $. Then, from $\beta^{\prime
}\gamma^{\prime}=\alpha$, we obtain $\alpha=\beta^{\prime}\gamma^{\prime}$.
Thus,
\begin{align*}
\left\vert \alpha\right\vert  &  =\left\vert \beta^{\prime}\gamma^{\prime
}\right\vert =\underbrace{\left\vert \beta^{\prime}\right\vert }%
_{=m}+\underbrace{\left\vert \gamma^{\prime}\right\vert }_{=n}%
\ \ \ \ \ \ \ \ \ \ \left(
\begin{array}
[c]{c}%
\text{by Proposition \ref{prop.concat.lsize} \textbf{(b)},}\\
\text{applied to }\beta=\beta^{\prime}\text{ and }\gamma=\gamma^{\prime}%
\end{array}
\right) \\
&  =m+n.
\end{align*}
Hence,%
\begin{align*}
m+n  &  =\left\vert \alpha\right\vert =\left\vert \beta^{\prime\prime}%
\gamma^{\prime\prime}\right\vert \ \ \ \ \ \ \ \ \ \ \left(  \text{since
}\alpha=\beta^{\prime\prime}\gamma^{\prime\prime}\right) \\
&  =\underbrace{\left\vert \beta^{\prime\prime}\right\vert }_{=m}+\left\vert
\gamma^{\prime\prime}\right\vert \ \ \ \ \ \ \ \ \ \ \left(
\begin{array}
[c]{c}%
\text{by Proposition \ref{prop.concat.lsize} \textbf{(b)},}\\
\text{applied to }\beta=\beta^{\prime\prime}\text{ and }\gamma=\gamma
^{\prime\prime}%
\end{array}
\right) \\
&  =m+\left\vert \gamma^{\prime\prime}\right\vert .
\end{align*}
Subtracting $m$ from this equality, we obtain $n=\left\vert \gamma
^{\prime\prime}\right\vert $.

Now, $\gamma^{\prime}$ is a composition having size $\left\vert \gamma
^{\prime}\right\vert =n$ (since $n=\left\vert \gamma^{\prime}\right\vert $).
In other words, $\gamma^{\prime}$ is a composition of $n$. In other words,
$\gamma^{\prime}\in\operatorname*{Comp}\nolimits_{n}$. Similarly,
$\gamma^{\prime\prime}\in\operatorname*{Comp}\nolimits_{n}$ (since
$n=\left\vert \gamma^{\prime\prime}\right\vert $).

Recall that the map $D:\operatorname*{Comp}\nolimits_{n}\rightarrow
\mathcal{P}\left(  \left[  n-1\right]  \right)  $ is a bijection. Hence, this
map $D$ is bijective, thus injective. In other words, if $\varphi,\psi
\in\operatorname*{Comp}\nolimits_{n}$ satisfy $D\left(  \varphi\right)
=D\left(  \psi\right)  $, then $\varphi=\psi$. Applying this to $\varphi
=\gamma^{\prime}$ and $\psi=\gamma^{\prime\prime}$, we obtain $\gamma^{\prime
}=\gamma^{\prime\prime}$ (since $\gamma^{\prime}\in\operatorname*{Comp}%
\nolimits_{n}$ and $\gamma^{\prime\prime}\in\operatorname*{Comp}\nolimits_{n}$
and $D\left(  \gamma^{\prime}\right)  =D\left(  \gamma^{\prime\prime}\right)
$).

Now, $\left(  \underbrace{\beta^{\prime}}_{=\beta^{\prime\prime}%
},\underbrace{\gamma^{\prime}}_{=\gamma^{\prime\prime}}\right)  =\left(
\beta^{\prime\prime},\gamma^{\prime\prime}\right)  $.

Forget that we fixed $\left(  \beta^{\prime},\gamma^{\prime}\right)  $ and
$\left(  \beta^{\prime\prime},\gamma^{\prime\prime}\right)  $. We thus have
shown that if $\left(  \beta^{\prime},\gamma^{\prime}\right)  $ and $\left(
\beta^{\prime\prime},\gamma^{\prime\prime}\right)  $ are two pairs $\left(
\beta,\gamma\right)  $ of compositions such that $\left\vert \beta\right\vert
=m$ and $\beta\gamma=\alpha$, then $\left(  \beta^{\prime},\gamma^{\prime
}\right)  =\left(  \beta^{\prime\prime},\gamma^{\prime\prime}\right)  $. In
other words, there exists at most one pair $\left(  \beta,\gamma\right)  $ of
compositions such that $\left\vert \beta\right\vert =m$ and $\beta
\gamma=\alpha$. This proves Proposition \ref{prop.concat.4u}.
\end{proof}

Next, we shall show a nearly trivial lemma:

\begin{lemma}
\label{lem.concat.4a.1}Let $m\in\mathbb{N}$. Let $K$ be a subset of $\left\{
1,2,3,\ldots\right\}  $. Then,
\[
\left(  K\cap\left[  m-1\right]  \right)  \cup\left(  K\setminus\left[
m\right]  \right)  =K\setminus\left\{  m\right\}  .
\]

\end{lemma}

\begin{proof}
[Proof of Lemma \ref{lem.concat.4a.1}.]Any element $k\in K$ is an element of
$\left\{  1,2,3,\ldots\right\}  $ (since $K$ is a subset of $\left\{
1,2,3,\ldots\right\}  $) and therefore is a positive integer. Hence, for any
element $k\in K$, we have the following chain of logical equivalences:%
\begin{align*}
\left(  k\in\left[  m-1\right]  \right)  \  &  \Longleftrightarrow\ \left(
k\leq m-1\right)  \ \ \ \ \ \ \ \ \ \ \left(  \text{since }k\text{ is a
positive integer}\right) \\
&  \Longleftrightarrow\ \left(  k<m\right)  \ \ \ \ \ \ \ \ \ \ \left(
\text{since }k\text{ and }m\text{ are integers}\right)  .
\end{align*}
Thus,%
\[
\left\{  k\in K\ \mid\ k\in\left[  m-1\right]  \right\}  =\left\{  k\in
K\ \mid\ k<m\right\}  .
\]


Recall again that any element $k\in K$ is a positive integer. Hence, for any
element $k\in K$, we have the following chain of logical equivalences:%
\begin{align*}
\left(  k\notin\left[  m\right]  \right)  \  &  \Longleftrightarrow\ \left(
\text{we don't have }k\in\left[  m\right]  \right) \\
&  \Longleftrightarrow\ \left(  \text{we don't have }k\leq m\right)
\ \ \ \ \ \ \ \ \ \ \left(
\begin{array}
[c]{c}%
\text{since }k\text{ is a positive integer,}\\
\text{and thus the statement \textquotedblleft}k\in\left[  m\right]
\text{\textquotedblright}\\
\text{is equivalent to \textquotedblleft}k\leq m\text{\textquotedblright}%
\end{array}
\right) \\
&  \Longleftrightarrow\ \left(  k>m\right)  .
\end{align*}
Hence,%
\[
\left\{  k\in K\ \mid\ k\notin\left[  m\right]  \right\}  =\left\{  k\in
K\ \mid\ k>m\right\}  .
\]


Now,
\begin{align*}
&  \underbrace{\left(  K\cap\left[  m-1\right]  \right)  }%
_{\substack{=\left\{  k\in K\ \mid\ k\in\left[  m-1\right]  \right\}
\\=\left\{  k\in K\ \mid\ k<m\right\}  }}\cup\underbrace{\left(
K\setminus\left[  m\right]  \right)  }_{\substack{\left\{  k\in K\ \mid
\ k\notin\left[  m\right]  \right\}  \\=\left\{  k\in K\ \mid\ k>m\right\}
}}\\
&  =\left\{  k\in K\ \mid\ k<m\right\}  \cup\left\{  k\in K\ \mid\ k>m\right\}
\\
&  =\left\{  k\in K\ \mid\ k<m\text{ or }k>m\right\} \\
&  =\left\{  k\in K\ \mid\ k\neq m\right\}  \ \ \ \ \ \ \ \ \ \ \left(
\begin{array}
[c]{c}%
\text{since the statement \textquotedblleft}k<m\text{ or }%
k>m\text{\textquotedblright}\\
\text{is equivalent to \textquotedblleft}k\neq m\text{\textquotedblright}%
\end{array}
\right) \\
&  =K\setminus\left\{  m\right\}  .
\end{align*}
This proves Lemma \ref{lem.concat.4a.1}.
\end{proof}

Our next proposition characterizes the concatenation $\varphi\psi$ of two
compositions $\varphi$ and $\psi$ in terms of how its partial sum set
$D\left(  \alpha\right)  $ relates to $D\left(  \varphi\right)  $ and
$D\left(  \psi\right)  $:

\begin{proposition}
\label{prop.concat.4a}Let $n\in\mathbb{N}$ and $m\in\mathbb{N}$. Let
$\alpha\in\operatorname*{Comp}\nolimits_{m+n}$ be any composition of $m+n$
such that $m\in D\left(  \alpha\right)  \cup\left\{  0,m+n\right\}  $.

Let $\varphi\in\operatorname*{Comp}\nolimits_{m}$ be a composition that
satisfies $D\left(  \varphi\right)  =D\left(  \alpha\right)  \cap\left[
m-1\right]  $.

Let $\psi\in\operatorname*{Comp}\nolimits_{n}$ be a composition that satisfies
$D\left(  \psi\right)  =\left(  D\left(  \alpha\right)  \setminus\left[
m\right]  \right)  -m$.

Then, $\varphi\psi=\alpha$.
\end{proposition}

\begin{proof}
[Proof of Proposition \ref{prop.concat.4a}.]We have $\varphi\in
\operatorname*{Comp}\nolimits_{m}$. In other words, $\varphi$ is a composition
of $m$. In other words, $\varphi$ is a composition having size $m$. In other
words, $\varphi\in\operatorname*{Comp}$ and $\left\vert \varphi\right\vert
=m$. Similarly, from $\psi\in\operatorname*{Comp}\nolimits_{n}$, we obtain
$\psi\in\operatorname*{Comp}$ and $\left\vert \psi\right\vert =n$.

Now, Proposition \ref{prop.concat.lsize} \textbf{(b)} (applied to
$\beta=\varphi$ and $\gamma=\psi$) yields $\left\vert \varphi\psi\right\vert
=\underbrace{\left\vert \varphi\right\vert }_{=m}+\underbrace{\left\vert
\psi\right\vert }_{=n}=m+n$. Thus, $\varphi\psi$ is a composition having size
$\left\vert \varphi\psi\right\vert =m+n$. In other words, $\varphi\psi$ is a
composition of $m+n$. In other words, $\varphi\psi\in\operatorname*{Comp}%
\nolimits_{m+n}$.

Recall that the map $D:\operatorname*{Comp}\nolimits_{m+n}\rightarrow
\mathcal{P}\left(  \left[  m+n-1\right]  \right)  $ is a bijection. Hence,
this map $D$ is bijective, thus injective. Furthermore, from $\alpha
\in\operatorname*{Comp}\nolimits_{m+n}$, we obtain $D\left(  \alpha\right)
\in\mathcal{P}\left(  \left[  m+n-1\right]  \right)  $ (since $D$ is a map
from $\operatorname*{Comp}\nolimits_{m+n}$ to $\mathcal{P}\left(  \left[
m+n-1\right]  \right)  $). In other words, $D\left(  \alpha\right)
\subseteq\left[  m+n-1\right]  $. Hence,%
\[
D\left(  \alpha\right)  \subseteq\left[  m+n-1\right]  \subseteq\left\{
1,2,3,\ldots\right\}  .
\]
In other words, $D\left(  \alpha\right)  $ is a subset of $\left\{
1,2,3,\ldots\right\}  $. Hence, Lemma \ref{lem.concat.4a.1} (applied to
$K=D\left(  \alpha\right)  $) yields%
\begin{equation}
\left(  D\left(  \alpha\right)  \cap\left[  m-1\right]  \right)  \cup\left(
D\left(  \alpha\right)  \setminus\left[  m\right]  \right)  =D\left(
\alpha\right)  \setminus\left\{  m\right\}  . \label{pf.prop.concat.4a.1}%
\end{equation}


If $K$ is any set of integers, then $\left(  K-m\right)  +m=K$ (indeed, this
follows easily from Definition \ref{def.K+m}). Applying this to $K=D\left(
\alpha\right)  \setminus\left[  m\right]  $, we obtain%
\[
\left(  \left(  D\left(  \alpha\right)  \setminus\left[  m\right]  \right)
-m\right)  +m=D\left(  \alpha\right)  \setminus\left[  m\right]  .
\]
In view of $D\left(  \psi\right)  =\left(  D\left(  \alpha\right)
\setminus\left[  m\right]  \right)  -m$, we can rewrite this as%
\begin{equation}
D\left(  \psi\right)  +m=D\left(  \alpha\right)  \setminus\left[  m\right]  .
\label{pf.prop.concat.4a.Dpsi+m}%
\end{equation}


Proposition \ref{prop.concat.D0} (applied to $\beta=\varphi$ and $\gamma=\psi
$) yields
\begin{align}
D\left(  \varphi\psi\right)   &  =\left(  \left\{  m\right\}  \cup
\underbrace{D\left(  \varphi\right)  }_{=D\left(  \alpha\right)  \cap\left[
m-1\right]  }\cup\underbrace{\left(  D\left(  \psi\right)  +m\right)
}_{\substack{=D\left(  \alpha\right)  \setminus\left[  m\right]  \\\text{(by
(\ref{pf.prop.concat.4a.Dpsi+m}))}}}\right)  \cap\left[  m+n-1\right]
\nonumber\\
&  =\left(  \left\{  m\right\}  \cup\underbrace{\left(  D\left(
\alpha\right)  \cap\left[  m-1\right]  \right)  \cup\left(  D\left(
\alpha\right)  \setminus\left[  m\right]  \right)  }_{\substack{=D\left(
\alpha\right)  \setminus\left\{  m\right\}  \\\text{(by
(\ref{pf.prop.concat.4a.1}))}}}\right)  \cap\left[  m+n-1\right] \nonumber\\
&  =\underbrace{\left(  \left\{  m\right\}  \cup\left(  D\left(
\alpha\right)  \setminus\left\{  m\right\}  \right)  \right)  }%
_{\substack{=\left\{  m\right\}  \cup D\left(  \alpha\right)  \\\text{(since
}\left(  X\cup\left(  Y\setminus X\right)  \right)  =X\cup Y\\\text{for any
two sets }X\text{ and }Y\text{)}}}\cap\left[  m+n-1\right] \nonumber\\
&  =\left(  \left\{  m\right\}  \cup D\left(  \alpha\right)  \right)
\cap\left[  m+n-1\right]  . \label{pf.prop.concat.4a.3}%
\end{align}


Now, we recall that $m\in D\left(  \alpha\right)  \cup\left\{  0,m+n\right\}
$ (by assumption). Hence, $\left\{  m\right\}  \subseteq D\left(
\alpha\right)  \cup\left\{  0,m+n\right\}  $. Thus,%
\begin{align*}
\underbrace{\left\{  m\right\}  }_{\subseteq D\left(  \alpha\right)
\cup\left\{  0,m+n\right\}  }\cup D\left(  \alpha\right)   &  \subseteq\left(
D\left(  \alpha\right)  \cup\left\{  0,m+n\right\}  \right)  \cup D\left(
\alpha\right) \\
&  =\underbrace{D\left(  \alpha\right)  \cup D\left(  \alpha\right)
}_{=D\left(  \alpha\right)  }\cup\left\{  0,m+n\right\} \\
&  =D\left(  \alpha\right)  \cup\left\{  0,m+n\right\}  .
\end{align*}
Hence,%
\begin{align*}
&  \underbrace{\left(  \left\{  m\right\}  \cup D\left(  \alpha\right)
\right)  }_{\subseteq D\left(  \alpha\right)  \cup\left\{  0,m+n\right\}
}\cap\left[  m+n-1\right] \\
&  \subseteq\left(  D\left(  \alpha\right)  \cup\left\{  0,m+n\right\}
\right)  \cap\left[  m+n-1\right] \\
&  =\left(  D\left(  \alpha\right)  \cap\left[  m+n-1\right]  \right)
\cup\underbrace{\left(  \left\{  0,m+n\right\}  \cap\left[  m+n-1\right]
\right)  }_{\substack{=\varnothing\\\text{(since neither }0\text{ nor
}m+n\text{ belongs}\\\text{to the set }\left[  m+n-1\right]  \text{)}}}\\
&  \ \ \ \ \ \ \ \ \ \ \ \ \ \ \ \ \ \ \ \ \left(
\begin{array}
[c]{c}%
\text{since }\left(  X\cup Y\right)  \cap Z=\left(  X\cap Z\right)
\cup\left(  Y\cap Z\right) \\
\text{for any three sets }X\text{, }Y\text{ and }Z
\end{array}
\right) \\
&  =\left(  D\left(  \alpha\right)  \cap\left[  m+n-1\right]  \right)
\cup\varnothing=D\left(  \alpha\right)  \cap\left[  m+n-1\right] \\
&  =D\left(  \alpha\right)  \ \ \ \ \ \ \ \ \ \ \left(  \text{since }D\left(
\alpha\right)  \subseteq\left[  m+n-1\right]  \right)  .
\end{align*}
Combining this inclusion with%
\begin{align*}
D\left(  \alpha\right)   &  =\underbrace{D\left(  \alpha\right)  }%
_{\subseteq\left(  \left\{  m\right\}  \cup D\left(  \alpha\right)  \right)
}\cap\left[  m+n-1\right]  \ \ \ \ \ \ \ \ \ \ \left(  \text{since }D\left(
\alpha\right)  \subseteq\left[  m+n-1\right]  \right) \\
&  \subseteq\left(  \left\{  m\right\}  \cup D\left(  \alpha\right)  \right)
\cap\left[  m+n-1\right]  ,
\end{align*}
we obtain%
\[
\left(  \left\{  m\right\}  \cup D\left(  \alpha\right)  \right)  \cap\left[
m+n-1\right]  =D\left(  \alpha\right)  .
\]
Hence, we can rewrite (\ref{pf.prop.concat.4a.3}) as $D\left(  \varphi
\psi\right)  =D\left(  \alpha\right)  $.

Now, recall that the map $D:\operatorname*{Comp}\nolimits_{m+n}\rightarrow
\mathcal{P}\left(  \left[  m+n-1\right]  \right)  $ is injective. Hence, if
$\zeta$ and $\eta$ are two elements of $\operatorname*{Comp}\nolimits_{m+n}$
satisfying $D\left(  \zeta\right)  =D\left(  \eta\right)  $, then $\zeta=\eta
$. Applying this to $\zeta=\varphi\psi$ and $\eta=\alpha$, we obtain
$\varphi\psi=\alpha$ (since $\varphi\psi\in\operatorname*{Comp}\nolimits_{m+n}%
$ and $\alpha\in\operatorname*{Comp}\nolimits_{m+n}$ and $D\left(  \varphi
\psi\right)  =D\left(  \alpha\right)  $). This proves Proposition
\ref{prop.concat.4a}.
\end{proof}

\subsection{Concatenation and coarsenings}

We shall next study the interaction between concatenation and coarsenings.
First, we define coarsenings:

\begin{definition}
If $\gamma$ is a composition, then $C\left(  \gamma\right)  $ shall denote the
set of all compositions $\beta\in\operatorname*{Comp}\nolimits_{\left\vert
\gamma\right\vert }$ satisfying $D\left(  \beta\right)  \subseteq D\left(
\gamma\right)  $.

The compositions belonging to $C\left(  \gamma\right)  $ are often called the
\emph{coarsenings} of $\gamma$.
\end{definition}

\begin{example}
Let $\gamma$ be the composition $\left(  4,1,2\right)  $. Then, the set
$C\left(  \gamma\right)  $ consists of the compositions $\beta\in
\operatorname*{Comp}\nolimits_{7}$ satisfying $D\left(  \beta\right)
\subseteq D\left(  \gamma\right)  =\left\{  4,5\right\}  $. Thus,%
\[
C\left(  \gamma\right)  =\left\{  \left(  7\right)  ,\ \left(  5,2\right)
,\ \left(  4,3\right)  ,\ \left(  4,1,2\right)  \right\}  .
\]
So the coarsenings of $\gamma$ are the four compositions $\left(  7\right)  $,
$\left(  5,2\right)  $, $\left(  4,3\right)  $ and $\left(  4,1,2\right)  $.
\end{example}

An equivalent definition of the coarsenings of a composition $\gamma$ can be
informally given as follows: If $\gamma$ is a composition, then a
\emph{coarsening} of $\gamma$ means a composition obtained by
\textquotedblleft combining\textquotedblright\ some groups of consecutive
entries of $\gamma$ (that is, replacing them by their sums). For instance, one
of the many coarsenings of a composition $\left(  \alpha_{1},\alpha_{2}%
,\alpha_{3},\alpha_{4},\alpha_{5},\alpha_{6},\alpha_{7}\right)  $ is $\left(
\alpha_{1}+\alpha_{2},\ \alpha_{3},\ \alpha_{4}+\alpha_{5}+\alpha_{6}%
,\ \alpha_{7}\right)  $. We shall not formalize this equivalent definition, as
we will not use it.

The following lemma is a trivial consequence of the definition of a
coarsening, restated for convenience:

\begin{lemma}
\label{lem.coarser.1}Let $\gamma$ be a composition.

\begin{enumerate}
\item[\textbf{(a)}] If $\nu\in C\left(  \gamma\right)  $, then $\nu
\in\operatorname*{Comp}$ and $\left\vert \nu\right\vert =\left\vert
\gamma\right\vert $ and $D\left(  \nu\right)  \subseteq D\left(
\gamma\right)  $.

\item[\textbf{(b)}] If $\nu\in\operatorname*{Comp}$ is a composition that
satisfies $\left\vert \nu\right\vert =\left\vert \gamma\right\vert $ and
$D\left(  \nu\right)  \subseteq D\left(  \gamma\right)  $, then $\nu\in
C\left(  \gamma\right)  $.
\end{enumerate}
\end{lemma}

\begin{proof}
\textbf{(a)} Assume that $\nu\in C\left(  \gamma\right)  $. According to the
definition of $C\left(  \gamma\right)  $, this means that $\nu$ is a
composition $\beta\in\operatorname*{Comp}\nolimits_{\left\vert \gamma
\right\vert }$ satisfying $D\left(  \beta\right)  \subseteq D\left(
\gamma\right)  $. In other words, $\nu\in\operatorname*{Comp}%
\nolimits_{\left\vert \gamma\right\vert }$ and $D\left(  \nu\right)  \subseteq
D\left(  \gamma\right)  $. From $\nu\in\operatorname*{Comp}%
\nolimits_{\left\vert \gamma\right\vert }$, we obtain $\nu\in
\operatorname*{Comp}$ and $\left\vert \nu\right\vert =\left\vert
\gamma\right\vert $. Thus, we have $\nu\in\operatorname*{Comp}$ and
$\left\vert \nu\right\vert =\left\vert \gamma\right\vert $ and $D\left(
\nu\right)  \subseteq D\left(  \gamma\right)  $. This proves Lemma
\ref{lem.coarser.1} \textbf{(a)}.

\textbf{(b)} Assume that $\nu\in\operatorname*{Comp}$ is a composition that
satisfies $\left\vert \nu\right\vert =\left\vert \gamma\right\vert $ and
$D\left(  \nu\right)  \subseteq D\left(  \gamma\right)  $. From $\nu
\in\operatorname*{Comp}$ and $\left\vert \nu\right\vert =\left\vert
\gamma\right\vert $, we obtain $\nu\in\operatorname*{Comp}%
\nolimits_{\left\vert \gamma\right\vert }$. Thus, $\nu$ is a composition
$\beta\in\operatorname*{Comp}\nolimits_{\left\vert \gamma\right\vert }$
satisfying $D\left(  \beta\right)  \subseteq D\left(  \gamma\right)  $ (since
$D\left(  \nu\right)  \subseteq D\left(  \gamma\right)  $). In other words,
$\nu\in C\left(  \gamma\right)  $ (by the definition of $C\left(
\gamma\right)  $). This proves Lemma \ref{lem.coarser.1} \textbf{(b)}.
\end{proof}

We can now restate Proposition \ref{prop.concat.2} in terms of coarsenings:

\begin{proposition}
\label{prop.concat.2C}Let $\beta$ and $\gamma$ be two compositions. Let
$\mu\in C\left(  \beta\right)  $ and $\nu\in C\left(  \gamma\right)  $. Then,
$\mu\nu\in C\left(  \beta\gamma\right)  $.
\end{proposition}

\begin{proof}
[Proof of Proposition \ref{prop.concat.2C}.]We have $\nu\in C\left(
\gamma\right)  $. Thus, Lemma \ref{lem.coarser.1} \textbf{(a)} yields that
$\nu\in\operatorname*{Comp}$ and $\left\vert \nu\right\vert =\left\vert
\gamma\right\vert $ and $D\left(  \nu\right)  \subseteq D\left(
\gamma\right)  $. The same argument (applied to $\mu$ and $\beta$ instead of
$\nu$ and $\gamma$) shows that $\mu\in\operatorname*{Comp}$ and $\left\vert
\mu\right\vert =\left\vert \beta\right\vert $ and $D\left(  \mu\right)
\subseteq D\left(  \beta\right)  $. Hence, Proposition \ref{prop.concat.2}
(applied to $\beta^{\prime}=\mu$ and $\gamma^{\prime}=\nu$) yields $D\left(
\mu\nu\right)  \subseteq D\left(  \beta\gamma\right)  $.

However, Proposition \ref{prop.concat.lsize} \textbf{(b)} yields $\left\vert
\beta\gamma\right\vert =\left\vert \beta\right\vert +\left\vert \gamma
\right\vert $.

Furthermore, Proposition \ref{prop.concat.lsize} \textbf{(b)} (applied to
$\mu$ and $\nu$ instead of $\beta$ and $\gamma$) yields
\[
\left\vert \mu\nu\right\vert =\underbrace{\left\vert \mu\right\vert
}_{=\left\vert \beta\right\vert }+\underbrace{\left\vert \nu\right\vert
}_{=\left\vert \gamma\right\vert }=\left\vert \beta\right\vert +\left\vert
\gamma\right\vert =\left\vert \beta\gamma\right\vert
\ \ \ \ \ \ \ \ \ \ \left(  \text{since }\left\vert \beta\gamma\right\vert
=\left\vert \beta\right\vert +\left\vert \gamma\right\vert \right)  .
\]
Thus, we now know that $\mu\nu\in\operatorname*{Comp}$ and $\left\vert \mu
\nu\right\vert =\left\vert \beta\gamma\right\vert $ and $D\left(  \mu
\nu\right)  \subseteq D\left(  \beta\gamma\right)  $. Hence, Lemma
\ref{lem.coarser.1} \textbf{(b)} (applied to $\beta\gamma$ and $\mu\nu$
instead of $\gamma$ and $\nu$) yields that $\mu\nu\in C\left(  \beta
\gamma\right)  $. This proves Proposition \ref{prop.concat.2C}.
\end{proof}

The following proposition is a sort of converse to Proposition
\ref{prop.concat.2C}:

\begin{proposition}
\label{prop.concat.4}Let $\alpha$ be a composition. Let $\mu$ and $\nu$ be two
compositions satisfying $\mu\nu\in C\left(  \alpha\right)  $. Then, there
exists a unique pair $\left(  \beta,\gamma\right)  \in\operatorname*{Comp}%
\times\operatorname*{Comp}$ of compositions satisfying $\beta\gamma=\alpha$
and $\mu\in C\left(  \beta\right)  $ and $\nu\in C\left(  \gamma\right)  $.
\end{proposition}

\begin{proof}
[Proof of Proposition \ref{prop.concat.4}.]Let $m=\left\vert \mu\right\vert $
and $n=\left\vert \nu\right\vert $. Then, $\mu\in\operatorname*{Comp}%
\nolimits_{m}$ (since $\mu$ is a composition that satisfies $\left\vert
\mu\right\vert =m$) and $\nu\in\operatorname*{Comp}\nolimits_{n}$ (since $\nu$
is a composition that satisfies $\left\vert \nu\right\vert =n$). Also, from
$\mu\nu\in C\left(  \alpha\right)  $, we conclude (by the definition of
$C\left(  \alpha\right)  $) that $\mu\nu\in\operatorname*{Comp}%
\nolimits_{\left\vert \alpha\right\vert }$ and $D\left(  \mu\nu\right)
\subseteq D\left(  \alpha\right)  $. Now, from $\mu\nu\in\operatorname*{Comp}%
\nolimits_{\left\vert \alpha\right\vert }$, we obtain $\left\vert \mu
\nu\right\vert =\left\vert \alpha\right\vert $. Thus, $\left\vert
\alpha\right\vert =\left\vert \mu\nu\right\vert $.

On the other hand, Proposition \ref{prop.concat.lsize} \textbf{(b)} (applied
to $\beta=\mu$ and $\gamma=\nu$) yields
\[
\left\vert \mu\nu\right\vert =\underbrace{\left\vert \mu\right\vert }%
_{=m}+\underbrace{\left\vert \nu\right\vert }_{=n}=m+n.
\]
Hence, $\left\vert \alpha\right\vert =\left\vert \mu\nu\right\vert =m+n$, so
that $\alpha\in\operatorname*{Comp}\nolimits_{m+n}$. Thus, $D\left(
\alpha\right)  \in\mathcal{P}\left(  \left[  m+n-1\right]  \right)  $ (since
$D:\operatorname*{Comp}\nolimits_{m+n}\rightarrow\mathcal{P}\left(  \left[
m+n-1\right]  \right)  $ is a bijection). In other words, $D\left(
\alpha\right)  \subseteq\left[  m+n-1\right]  $.

It is furthermore easy to see that%
\[
m\in D\left(  \alpha\right)  \cup\left\{  0,m+n\right\}
\]
\footnote{\textit{Proof.} We are in one of the following three cases:
\par
\textit{Case 1:} We have $m=0$.
\par
\textit{Case 2:} We have $n=0$.
\par
\textit{Case 3:} We have neither $m=0$ nor $n=0$.
\par
Let us first consider Case 1. In this case, we have $m=0$. Hence,
$m=0\in\left\{  0,m+n\right\}  \subseteq D\left(  \alpha\right)  \cup\left\{
0,m+n\right\}  $. Thus, $m\in D\left(  \alpha\right)  \cup\left\{
0,m+n\right\}  $ is proved in Case 1.
\par
Let us next consider Case 2. In this case, we have $n=0$. Hence,
$m+\underbrace{n}_{=0}=m$, so that $m=m+n\in\left\{  0,m+n\right\}  \subseteq
D\left(  \alpha\right)  \cup\left\{  0,m+n\right\}  $. Thus, $m\in D\left(
\alpha\right)  \cup\left\{  0,m+n\right\}  $ is proved in Case 2.
\par
Now, let us consider Case 3. In this case, we have neither $m=0$ nor $n=0$.
Hence, $m\neq0$ and $n\neq0$. Therefore, $\mu\neq\varnothing$ (since
$\left\vert \mu\right\vert =m\neq0=\left\vert \varnothing\right\vert $) and
$\nu\neq\varnothing$ (since $\left\vert \nu\right\vert =n\neq0=\left\vert
\varnothing\right\vert $). Hence, Proposition \ref{prop.concat.D} (applied to
$\mu$ and $\nu$ instead of $\beta$ and $\gamma$) yields%
\begin{equation}
D\left(  \mu\nu\right)  =\left\{  m\right\}  \cup D\left(  \mu\right)
\cup\left(  D\left(  \nu\right)  +m\right)  . \label{pf.prop.concat.4.fn1.1}%
\end{equation}
Now,
\begin{align*}
m  &  \in\left\{  m\right\}  \subseteq\left\{  m\right\}  \cup D\left(
\mu\right)  \cup\left(  D\left(  \nu\right)  +m\right) \\
&  =D\left(  \mu\nu\right)  \ \ \ \ \ \ \ \ \ \ \left(  \text{by
(\ref{pf.prop.concat.4.fn1.1})}\right) \\
&  \subseteq D\left(  \alpha\right)  \subseteq D\left(  \alpha\right)
\cup\left\{  0,m+n\right\}  .
\end{align*}
Thus, $m\in D\left(  \alpha\right)  \cup\left\{  0,m+n\right\}  $ is proved in
Case 3.
\par
Hence, we have proved $m\in D\left(  \alpha\right)  \cup\left\{
0,m+n\right\}  $ in all three Cases 1, 2 and 3. Thus, $m\in D\left(
\alpha\right)  \cup\left\{  0,m+n\right\}  $ always holds.}.

We have $D\left(  \alpha\right)  \cap\left[  m-1\right]  \subseteq\left[
m-1\right]  $, so that $D\left(  \alpha\right)  \cap\left[  m-1\right]
\in\mathcal{P}\left(  \left[  m-1\right]  \right)  $.

Furthermore, it is easy to see that $\left(  D\left(  \alpha\right)
\setminus\left[  m\right]  \right)  -m\in\mathcal{P}\left(  \left[
n-1\right]  \right)  $\ \ \ \ \footnote{\textit{Proof.} Let $g\in\left(
D\left(  \alpha\right)  \setminus\left[  m\right]  \right)  -m$. We shall show
that $g\in\left[  n-1\right]  $.
\par
Indeed,%
\[
g\in\left(  D\left(  \alpha\right)  \setminus\left[  m\right]  \right)
-m=\left\{  k-m\ \mid\ k\in D\left(  \alpha\right)  \setminus\left[  m\right]
\right\}
\]
(by the definition of $\left(  D\left(  \alpha\right)  \setminus\left[
m\right]  \right)  -m$). In other words, $g=k-m$ for some $k\in D\left(
\alpha\right)  \setminus\left[  m\right]  $. Consider this $k$.
\par
We have $k\in D\left(  \alpha\right)  \setminus\left[  m\right]  $, so that
$k\in D\left(  \alpha\right)  $ and $k\notin\left[  m\right]  $. From $k\in
D\left(  \alpha\right)  \subseteq\left[  m+n-1\right]  $, we obtain $1\leq
k\leq m+n-1$. If we had $k\leq m$, then we would have $k\in\left[  m\right]  $
(since $1\leq k\leq m$), which would contradict $k\notin\left[  m\right]  $.
Thus, we cannot have $k\leq m$. Hence, we must have $k>m$. Thus, $k\geq m+1$
(since $k$ and $m$ are integers), so that $k-m\geq1$. Furthermore, from $k\leq
m+n-1$, we obtain $k-m\leq n-1$. Combining $k-m\geq1$ with $k-m\leq n-1$, we
find $k-m\in\left\{  1,2,\ldots,n-1\right\}  =\left[  n-1\right]  $. Thus,
$g=k-m\in\left[  n-1\right]  $.
\par
Forget now that we fixed $g$. We thus have shown that $g\in\left[  n-1\right]
$ for each $g\in\left(  D\left(  \alpha\right)  \setminus\left[  m\right]
\right)  -m$. In other words, $\left(  D\left(  \alpha\right)  \setminus
\left[  m\right]  \right)  -m\subseteq\left[  n-1\right]  $. In other words,
$\left(  D\left(  \alpha\right)  \setminus\left[  m\right]  \right)
-m\in\mathcal{P}\left(  \left[  n-1\right]  \right)  $.}.

Recall that the map $D:\operatorname*{Comp}\nolimits_{m}\rightarrow
\mathcal{P}\left(  \left[  m-1\right]  \right)  $ is a bijection. Hence, it is
bijective, thus surjective. Therefore, there exists some composition
$\varphi\in\operatorname*{Comp}\nolimits_{m}$ that satisfies%
\begin{equation}
D\left(  \varphi\right)  =D\left(  \alpha\right)  \cap\left[  m-1\right]
\label{pf.prop.concat.4.Dphi}%
\end{equation}
(since $D\left(  \alpha\right)  \cap\left[  m-1\right]  \in\mathcal{P}\left(
\left[  m-1\right]  \right)  $). Consider this $\varphi$.

Recall that the map $D:\operatorname*{Comp}\nolimits_{n}\rightarrow
\mathcal{P}\left(  \left[  n-1\right]  \right)  $ is a bijection. Hence, it is
bijective, thus surjective. Therefore, there exists some composition $\psi
\in\operatorname*{Comp}\nolimits_{n}$ that satisfies%
\begin{equation}
D\left(  \psi\right)  =\left(  D\left(  \alpha\right)  \setminus\left[
m\right]  \right)  -m \label{pf.prop.concat.4.Dpsi}%
\end{equation}
(since $\left(  D\left(  \alpha\right)  \setminus\left[  m\right]  \right)
-m\in\mathcal{P}\left(  \left[  n-1\right]  \right)  $). Consider this $\psi$.

Proposition \ref{prop.concat.4a} yields that
\[
\varphi\psi=\alpha.
\]
Also, $\varphi\in\operatorname*{Comp}\nolimits_{m}\subseteq
\operatorname*{Comp}$ and $\left\vert \varphi\right\vert =m$ (since
$\varphi\in\operatorname*{Comp}\nolimits_{m}$). Furthermore, $\psi
\in\operatorname*{Comp}\nolimits_{n}\subseteq\operatorname*{Comp}$ and
$\left\vert \psi\right\vert =n$ (since $\psi\in\operatorname*{Comp}%
\nolimits_{n}$). From $\varphi\in\operatorname*{Comp}$ and $\psi
\in\operatorname*{Comp}$, we obtain $\left(  \varphi,\psi\right)
\in\operatorname*{Comp}\times\operatorname*{Comp}$.

Proposition \ref{prop.concat.Dcut} \textbf{(a)} (applied to $\beta=\mu$ and
$\gamma=\nu$) yields%
\[
D\left(  \mu\right)  =\underbrace{D\left(  \mu\nu\right)  }_{\subseteq
D\left(  \alpha\right)  }\cap\left[  m-1\right]  \subseteq D\left(
\alpha\right)  \cap\left[  m-1\right]  =D\left(  \varphi\right)
\]
(by (\ref{pf.prop.concat.4.Dphi})). Also, we have $\left\vert \mu\right\vert
=m=\left\vert \varphi\right\vert $ (since $\left\vert \varphi\right\vert =m$),
so that $\mu\in\operatorname*{Comp}\nolimits_{\left\vert \varphi\right\vert }%
$. Thus, $\mu$ is a composition $\beta\in\operatorname*{Comp}%
\nolimits_{\left\vert \varphi\right\vert }$ satisfying $D\left(  \beta\right)
\subseteq D\left(  \varphi\right)  $ (since we have shown that $D\left(
\mu\right)  \subseteq D\left(  \varphi\right)  $). In other words, $\mu\in
C\left(  \varphi\right)  $ (by the definition of $C\left(  \varphi\right)  $).

Proposition \ref{prop.concat.Dcut} \textbf{(b)} (applied to $\beta=\mu$ and
$\gamma=\nu$) yields%
\begin{equation}
D\left(  \nu\right)  =\left(  D\left(  \mu\nu\right)  \setminus\left[
m\right]  \right)  -m. \label{pf.prop.concat.4.Dnu.1}%
\end{equation}
However, $\underbrace{D\left(  \mu\nu\right)  }_{\subseteq D\left(
\alpha\right)  }\setminus\left[  m\right]  \subseteq D\left(  \alpha\right)
\setminus\left[  m\right]  $. But it is easy to see that if $k$ is any
integer, and if $K$ and $K^{\prime}$ are two sets of integers satisfying
$K\subseteq K^{\prime}$, then $K-k\subseteq K^{\prime}-k$. Applying this to
$k=m$ and $K=D\left(  \mu\nu\right)  \setminus\left[  m\right]  $ and
$K^{\prime}=D\left(  \alpha\right)  \setminus\left[  m\right]  $, we conclude
that $\left(  D\left(  \mu\nu\right)  \setminus\left[  m\right]  \right)
-m\subseteq\left(  D\left(  \alpha\right)  \setminus\left[  m\right]  \right)
-m$ (since $D\left(  \mu\nu\right)  \setminus\left[  m\right]  \subseteq
D\left(  \alpha\right)  \setminus\left[  m\right]  $). In view of
(\ref{pf.prop.concat.4.Dnu.1}), we can rewrite this as%
\[
D\left(  \nu\right)  \subseteq \left(  D\left(  \alpha\right)  \setminus\left[
m\right]  \right)  -m=D\left(  \psi\right)
\]
(by (\ref{pf.prop.concat.4.Dpsi})). Also, we have $\left\vert \nu\right\vert
=n=\left\vert \psi\right\vert $ (since $\left\vert \psi\right\vert =n$), so
that $\nu\in\operatorname*{Comp}\nolimits_{\left\vert \psi\right\vert }$.
Thus, $\nu$ is a composition $\beta\in\operatorname*{Comp}%
\nolimits_{\left\vert \psi\right\vert }$ satisfying $D\left(  \beta\right)
\subseteq D\left(  \psi\right)  $ (since we have shown that $D\left(
\nu\right)  \subseteq D\left(  \psi\right)  $). In other words, $\nu\in
C\left(  \psi\right)  $ (by the definition of $C\left(  \psi\right)  $).

We have now shown that $\left(  \varphi,\psi\right)  \in\operatorname*{Comp}%
\times\operatorname*{Comp}$ is a pair of compositions satisfying $\varphi
\psi=\alpha$ and $\mu\in C\left(  \varphi\right)  $ and $\nu\in C\left(
\psi\right)  $. Hence, there exists \textbf{at least} one pair $\left(
\beta,\gamma\right)  \in\operatorname*{Comp}\times\operatorname*{Comp}$ of
compositions satisfying $\beta\gamma=\alpha$ and $\mu\in C\left(
\beta\right)  $ and $\nu\in C\left(  \gamma\right)  $ (namely, the pair
$\left(  \varphi,\psi\right)  $).

It remains to show that there exists \textbf{only} one such pair. So let us
show this now.

Indeed, let $\left(  \beta^{\prime},\gamma^{\prime}\right)  $ and $\left(
\beta^{\prime\prime},\gamma^{\prime\prime}\right)  $ be two pairs $\left(
\beta,\gamma\right)  \in\operatorname*{Comp}\times\operatorname*{Comp}$ of
compositions satisfying $\beta\gamma=\alpha$ and $\mu\in C\left(
\beta\right)  $ and $\nu\in C\left(  \gamma\right)  $. We must prove that
$\left(  \beta^{\prime},\gamma^{\prime}\right)  =\left(  \beta^{\prime\prime
},\gamma^{\prime\prime}\right)  $.

We know that $\left(  \beta^{\prime},\gamma^{\prime}\right)  $ is a pair
$\left(  \beta,\gamma\right)  \in\operatorname*{Comp}\times
\operatorname*{Comp}$ of compositions satisfying $\beta\gamma=\alpha$ and
$\mu\in C\left(  \beta\right)  $ and $\nu\in C\left(  \gamma\right)  $. In
other words, $\left(  \beta^{\prime},\gamma^{\prime}\right)  \in
\operatorname*{Comp}\times\operatorname*{Comp}$ is a pair of compositions
satisfying $\beta^{\prime}\gamma^{\prime}=\alpha$ and $\mu\in C\left(
\beta^{\prime}\right)  $ and $\nu\in C\left(  \gamma^{\prime}\right)  $. From
$\mu\in C\left(  \beta^{\prime}\right)  $, we easily obtain $\left\vert
\beta^{\prime}\right\vert =m$\ \ \ \ \footnote{\textit{Proof.} We have $\mu\in
C\left(  \beta^{\prime}\right)  $. By the definition of $C\left(
\beta^{\prime}\right)  $, this means that $\mu$ is a composition $\beta
\in\operatorname*{Comp}\nolimits_{\left\vert \beta^{\prime}\right\vert }$
satisfying $D\left(  \beta\right)  \subseteq D\left(  \beta^{\prime}\right)
$. In other words, $\mu\in\operatorname*{Comp}\nolimits_{\left\vert
\beta^{\prime}\right\vert }$ and $D\left(  \mu\right)  \subseteq D\left(
\beta^{\prime}\right)  $. Hence, $\mu\in\operatorname*{Comp}%
\nolimits_{\left\vert \beta^{\prime}\right\vert }$, so that $\left\vert
\mu\right\vert =\left\vert \beta^{\prime}\right\vert $. Thus, $\left\vert
\beta^{\prime}\right\vert =\left\vert \mu\right\vert =m$.}.

We have now shown that $\left\vert \beta^{\prime}\right\vert =m$ and
$\beta^{\prime}\gamma^{\prime}=\alpha$. In other words, $\left(  \beta
^{\prime},\gamma^{\prime}\right)  $ is a pair $\left(  \beta,\gamma\right)  $
of compositions such that $\left\vert \beta\right\vert =m$ and $\beta
\gamma=\alpha$. The same argument (applied to $\left(  \beta^{\prime\prime
},\gamma^{\prime\prime}\right)  $ instead of $\left(  \beta^{\prime}%
,\gamma^{\prime}\right)  $) shows that $\left(  \beta^{\prime\prime}%
,\gamma^{\prime\prime}\right)  $ is such a pair as well.

However, Proposition \ref{prop.concat.4u} shows that there exists at most one
pair $\left(  \beta,\gamma\right)  $ of compositions such that $\left\vert
\beta\right\vert =m$ and $\beta\gamma=\alpha$. Hence, any two such pairs
$\left(  \beta,\gamma\right)  $ must be equal. Since $\left(  \beta^{\prime
},\gamma^{\prime}\right)  $ and $\left(  \beta^{\prime\prime},\gamma
^{\prime\prime}\right)  $ are two such pairs (as we have shown in the previous
paragraph), we thus can conclude that $\left(  \beta^{\prime},\gamma^{\prime
}\right)  $ and $\left(  \beta^{\prime\prime},\gamma^{\prime\prime}\right)  $
must be equal. In other words, $\left(  \beta^{\prime},\gamma^{\prime}\right)
=\left(  \beta^{\prime\prime},\gamma^{\prime\prime}\right)  $.

Now, forget that we fixed $\left(  \beta^{\prime},\gamma^{\prime}\right)  $
and $\left(  \beta^{\prime\prime},\gamma^{\prime\prime}\right)  $. We thus
have shown that if $\left(  \beta^{\prime},\gamma^{\prime}\right)  $ and
$\left(  \beta^{\prime\prime},\gamma^{\prime\prime}\right)  $ are two pairs
$\left(  \beta,\gamma\right)  \in\operatorname*{Comp}\times
\operatorname*{Comp}$ of compositions satisfying $\beta\gamma=\alpha$ and
$\mu\in C\left(  \beta\right)  $ and $\nu\in C\left(  \gamma\right)  $, then
$\left(  \beta^{\prime},\gamma^{\prime}\right)  =\left(  \beta^{\prime\prime
},\gamma^{\prime\prime}\right)  $. In other words, any two pairs $\left(
\beta,\gamma\right)  \in\operatorname*{Comp}\times\operatorname*{Comp}$ of
compositions satisfying $\beta\gamma=\alpha$ and $\mu\in C\left(
\beta\right)  $ and $\nu\in C\left(  \gamma\right)  $ must be equal. In other
words, there exists \textbf{at most one} such pair $\left(  \beta
,\gamma\right)  $. Since we also know that there exists \textbf{at least one}
such pair $\left(  \beta,\gamma\right)  $ (because we have proved this further
above), we thus conclude that there exists a \textbf{unique} such pair
$\left(  \beta,\gamma\right)  $. This proves Proposition \ref{prop.concat.4}.
\end{proof}

We can combine Propositions \ref{prop.concat.2C} and \ref{prop.concat.4} into
a convenient package:

\begin{proposition}
\label{prop.concat.sum=sum}Let $\left(  A,+,0\right)  $ be an abelian group.
Let $u_{\mu,\nu}$ be an element of $A$ for each pair $\left(  \mu,\nu\right)
\in\operatorname*{Comp}\times\operatorname*{Comp}$ of two compositions. Let
$\alpha\in\operatorname*{Comp}$. Then,%
\[
\sum_{\substack{\left(  \mu,\nu\right)  \in\operatorname*{Comp}\times
\operatorname*{Comp};\\\mu\nu\in C\left(  \alpha\right)  }}u_{\mu,\nu}%
=\sum_{\substack{\left(  \beta,\gamma\right)  \in\operatorname*{Comp}%
\times\operatorname*{Comp};\\\beta\gamma=\alpha}}\ \ \sum_{\mu\in C\left(
\beta\right)  }\ \ \sum_{\nu\in C\left(  \gamma\right)  }u_{\mu,\nu}.
\]

\end{proposition}

\begin{proof}
[Proof of Proposition \ref{prop.concat.sum=sum}.]We have\footnote{As
usual, we allow infinite sums, as long as all but finitely many of their
addends are $0$.}
\begin{align}
&  \sum_{\substack{\left(  \beta,\gamma\right)  \in\operatorname*{Comp}%
\times\operatorname*{Comp};\\\beta\gamma=\alpha}}\ \ \underbrace{\sum_{\mu\in
C\left(  \beta\right)  }}_{\substack{=\sum_{\substack{\mu\in
\operatorname*{Comp};\\\mu\in C\left(  \beta\right)  }}\\\text{(since
}C\left(  \beta\right)  \subseteq\operatorname*{Comp}\text{)}}%
}\ \ \underbrace{\sum_{\nu\in C\left(  \gamma\right)  }}_{\substack{=\sum
_{\substack{\nu\in\operatorname*{Comp};\\\nu\in C\left(  \gamma\right)
}}\\\text{(since }C\left(  \gamma\right)  \subseteq\operatorname*{Comp}%
\text{)}}}u_{\mu,\nu}\nonumber\\
&  =\underbrace{\sum_{\substack{\left(  \beta,\gamma\right)  \in
\operatorname*{Comp}\times\operatorname*{Comp};\\\beta\gamma=\alpha}%
}\ \ \sum_{\substack{\mu\in\operatorname*{Comp};\\\mu\in C\left(
\beta\right)  }}\ \ \sum_{\substack{\nu\in\operatorname*{Comp};\\\nu\in
C\left(  \gamma\right)  }}}_{=\sum_{\mu\in\operatorname*{Comp}}\ \ \sum
_{\nu\in\operatorname*{Comp}}\ \ \sum_{\substack{\left(  \beta,\gamma\right)
\in\operatorname*{Comp}\times\operatorname*{Comp};\\\beta\gamma=\alpha
;\\\mu\in C\left(  \beta\right)  ;\\\nu\in C\left(  \gamma\right)  }}}%
u_{\mu,\nu}\nonumber\\
&  =\underbrace{\sum_{\mu\in\operatorname*{Comp}}\ \ \sum_{\nu\in
\operatorname*{Comp}}}_{=\sum_{\left(  \mu,\nu\right)  \in\operatorname*{Comp}%
\times\operatorname*{Comp}}}\ \ \sum_{\substack{\left(  \beta,\gamma\right)
\in\operatorname*{Comp}\times\operatorname*{Comp};\\\beta\gamma=\alpha
;\\\mu\in C\left(  \beta\right)  ;\\\nu\in C\left(  \gamma\right)  }%
}u_{\mu,\nu}\nonumber\\
&  =\sum_{\left(  \mu,\nu\right)  \in\operatorname*{Comp}\times
\operatorname*{Comp}}\ \ \sum_{\substack{\left(  \beta,\gamma\right)
\in\operatorname*{Comp}\times\operatorname*{Comp};\\\beta\gamma=\alpha
;\\\mu\in C\left(  \beta\right)  ;\\\nu\in C\left(  \gamma\right)  }%
}u_{\mu,\nu}. \label{pf.prop.concat.sum=sum.1}%
\end{align}


Now, we claim the following:

\begin{statement}
\textit{Claim 1:} Let $\left(  \mu,\nu\right)  \in\operatorname*{Comp}%
\times\operatorname*{Comp}$ be such that $\mu\nu\in C\left(  \alpha\right)  $.
Then,
\[
\sum_{\substack{\left(  \beta,\gamma\right)  \in\operatorname*{Comp}%
\times\operatorname*{Comp};\\\beta\gamma=\alpha;\\\mu\in C\left(
\beta\right)  ;\\\nu\in C\left(  \gamma\right)  }}u_{\mu,\nu}=u_{\mu,\nu}.
\]

\end{statement}

[\textit{Proof of Claim 1:} Proposition \ref{prop.concat.4} shows that there
exists a unique pair $\left(  \beta,\gamma\right)  \in\operatorname*{Comp}%
\times\operatorname*{Comp}$ of compositions satisfying $\beta\gamma=\alpha$
and $\mu\in C\left(  \beta\right)  $ and $\nu\in C\left(  \gamma\right)  $. In
other words, the sum $\sum_{\substack{\left(  \beta,\gamma\right)
\in\operatorname*{Comp}\times\operatorname*{Comp};\\\beta\gamma=\alpha
;\\\mu\in C\left(  \beta\right)  ;\\\nu\in C\left(  \gamma\right)  }%
}u_{\mu,\nu}$ has exactly one addend. Hence, this sum equals $u_{\mu,\nu}$.
This proves Claim 1.]

\begin{statement}
\textit{Claim 2:} Let $\left(  \mu,\nu\right)  \in\operatorname*{Comp}%
\times\operatorname*{Comp}$ be such that $\mu\nu\notin C\left(  \alpha\right)
$. Then,
\[
\sum_{\substack{\left(  \beta,\gamma\right)  \in\operatorname*{Comp}%
\times\operatorname*{Comp};\\\beta\gamma=\alpha;\\\mu\in C\left(
\beta\right)  ;\\\nu\in C\left(  \gamma\right)  }}u_{\mu,\nu}=0.
\]

\end{statement}

[\textit{Proof of Claim 2:} If $\left(  \beta,\gamma\right)  \in
\operatorname*{Comp}\times\operatorname*{Comp}$ is a pair of compositions
satisfying $\beta\gamma=\alpha$ and $\mu\in C\left(  \beta\right)  $ and
$\nu\in C\left(  \gamma\right)  $, then Proposition \ref{prop.concat.2C} shows
that $\mu\nu\in C\left(  \underbrace{\beta\gamma}_{=\alpha}\right)  =C\left(
\alpha\right)  $, which contradicts $\mu\nu\notin C\left(  \alpha\right)  $.
Hence, there exists no pair $\left(  \beta,\gamma\right)  \in
\operatorname*{Comp}\times\operatorname*{Comp}$ of compositions satisfying
$\beta\gamma=\alpha$ and $\mu\in C\left(  \beta\right)  $ and $\nu\in C\left(
\gamma\right)  $. In other words, the sum $\sum_{\substack{\left(
\beta,\gamma\right)  \in\operatorname*{Comp}\times\operatorname*{Comp}%
;\\\beta\gamma=\alpha;\\\mu\in C\left(  \beta\right)  ;\\\nu\in C\left(
\gamma\right)  }}u_{\mu,\nu}$ is empty. Therefore, this sum equals $0$. This
proves Claim 2.]

Now, each pair $\left(  \mu,\nu\right)  \in\operatorname*{Comp}\times
\operatorname*{Comp}$ satisfies either $\mu\nu\in C\left(  \alpha\right)  $ or
$\mu\nu\notin C\left(  \alpha\right)  $ (but not both). Hence, we can split
the outer sum on the right hand side of (\ref{pf.prop.concat.sum=sum.1}) as
follows:%
\begin{align*}
&  \sum_{\left(  \mu,\nu\right)  \in\operatorname*{Comp}\times
\operatorname*{Comp}}\ \ \sum_{\substack{\left(  \beta,\gamma\right)
\in\operatorname*{Comp}\times\operatorname*{Comp};\\\beta\gamma=\alpha
;\\\mu\in C\left(  \beta\right)  ;\\\nu\in C\left(  \gamma\right)  }%
}u_{\mu,\nu}\\
&  =\sum_{\substack{\left(  \mu,\nu\right)  \in\operatorname*{Comp}%
\times\operatorname*{Comp};\\\mu\nu\in C\left(  \alpha\right)  }%
}\ \ \underbrace{\sum_{\substack{\left(  \beta,\gamma\right)  \in
\operatorname*{Comp}\times\operatorname*{Comp};\\\beta\gamma=\alpha;\\\mu\in
C\left(  \beta\right)  ;\\\nu\in C\left(  \gamma\right)  }}u_{\mu,\nu}%
}_{\substack{=u_{\mu,\nu}\\\text{(by Claim 1)}}}+\sum_{\substack{\left(
\mu,\nu\right)  \in\operatorname*{Comp}\times\operatorname*{Comp};\\\mu
\nu\notin C\left(  \alpha\right)  }}\ \ \underbrace{\sum_{\substack{\left(
\beta,\gamma\right)  \in\operatorname*{Comp}\times\operatorname*{Comp}%
;\\\beta\gamma=\alpha;\\\mu\in C\left(  \beta\right)  ;\\\nu\in C\left(
\gamma\right)  }}u_{\mu,\nu}}_{\substack{=0\\\text{(by Claim 2)}}}\\
&  =\sum_{\substack{\left(  \mu,\nu\right)  \in\operatorname*{Comp}%
\times\operatorname*{Comp};\\\mu\nu\in C\left(  \alpha\right)  }}u_{\mu,\nu
}+\underbrace{\sum_{\substack{\left(  \mu,\nu\right)  \in\operatorname*{Comp}%
\times\operatorname*{Comp};\\\mu\nu\notin C\left(  \alpha\right)  }}0}%
_{=0}=\sum_{\substack{\left(  \mu,\nu\right)  \in\operatorname*{Comp}%
\times\operatorname*{Comp};\\\mu\nu\in C\left(  \alpha\right)  }}u_{\mu,\nu}.
\end{align*}
Hence, we can rewrite (\ref{pf.prop.concat.sum=sum.1}) as%
\[
\sum_{\substack{\left(  \beta,\gamma\right)  \in\operatorname*{Comp}%
\times\operatorname*{Comp};\\\beta\gamma=\alpha}}\ \ \sum_{\mu\in C\left(
\beta\right)  }\ \ \sum_{\nu\in C\left(  \gamma\right)  }u_{\mu,\nu}%
=\sum_{\substack{\left(  \mu,\nu\right)  \in\operatorname*{Comp}%
\times\operatorname*{Comp};\\\mu\nu\in C\left(  \alpha\right)  }}u_{\mu,\nu}.
\]
This proves Proposition \ref{prop.concat.sum=sum}.
\end{proof}

\begin{thebibliography}{99999999}                                                                                         %


\bibitem[Grinbe15]{Grinbe15}\href{https://arxiv.org/abs/1509.08355v4}{Darij
Grinberg, \textit{Double posets and the antipode of QSym}, arXiv:1509.08355v4}.

\bibitem[GriVas22]{main}Darij Grinberg, Ekaterina A. Vassilieva, \textit{The
enriched $q$-monomial basis of the quasisymmetric functions}, arXiv:2309.01118.
\end{thebibliography}


\end{document}